% Liberté et Responsabilité — Philosophie Tle Tech — Cours signé OnBuch+
% Programme officiel MINESEC. Compiler avec : tectonic N15-liberte-responsabilite.tex
\newcommand{\DOCMATIERE}{Philosophie}
\newcommand{\DOCNIVEAU}{Tle Tech}
\newcommand{\DOCMODULE}{Philosophie – classe de Terminale technique}
\newcommand{\DOCLECON}{N15}
\newcommand{\DOCTITRE}{Liberté et Responsabilité}
% OnBuch+ — préambule « cours signés » ENRICHI (compilé avec tectonic / XeLaTeX)
% Système visuel « Pop » d'OnBuch+ : fond crème (jamais blanc), bordures encre
% épaisses, ombres « dures » décalées, titres Archivo Black en capitales.
% Version MATHÉMATIQUES : graphes (pgfplots/tikz), boîtes pédagogiques
% (définition, propriété, méthode, attention, le savais-tu, exercice résolu,
% exercice, TP, bilan), page de garde et signature OnBuch+ ancrée en bas.
%
% Macros attendues dans main.tex AVANT \input{preamble} :
%   \DOCMATIERE  \DOCNIVEAU  \DOCMODULE  \DOCLECON (numéro)  \DOCTITRE
\documentclass[11pt]{article}

\usepackage{fontspec}
\usepackage[french]{babel}
\usepackage[a4paper,margin=2cm,top=3.4cm,bottom=2.8cm,headheight=1.4cm,footskip=1.3cm]{geometry}
\usepackage{amsmath,amssymb,amsfonts}
\usepackage{mathtools} % \xmapsto, \coloneqq... souvent utilisés par les modèles
\usepackage{cancel} % simplifications barrées (bilans de forces)
% \abs{z} (module, valeur absolue) et \norm{u} (norme), utilisés par les modèles
\providecommand{\abs}{}\let\abs\relax\DeclarePairedDelimiter{\abs}{\lvert}{\rvert}
\providecommand{\norm}{}\let\norm\relax\DeclarePairedDelimiter{\norm}{\lVert}{\rVert}
\usepackage[table]{xcolor} % [table] : fournit aussi \rowcolors (alternance de lignes)
\usepackage{pagecolor}
\usepackage{tikz}
\usetikzlibrary{arrows.meta,positioning,calc,shapes.geometric,shapes.misc,shapes.arrows,decorations.pathmorphing,decorations.pathreplacing,decorations.markings,patterns,shadows,mindmap,trees,fit,backgrounds,matrix,intersections,angles,quotes}
\usepackage{pgfplots}
\usepackage[version=4]{mhchem} % équations nucléaires \ce{^{235}_{92}U}
\usepackage[european,siunitx]{circuitikz} % schémas électriques
\pgfplotsset{compat=1.18}
\usepgfplotslibrary{fillbetween,groupplots} % aires entre courbes (calcul intégral)
\usepackage{enumitem}
\usepackage{fancyhdr}
% [most] charge toutes les bibliothèques tcolorbox (listings, minted, raster,
% documentation...) et gonfle inutilement la mémoire TeX sur un document long
% (~300 boîtes) : on ne charge que ce qui est réellement utilisé.
\usepackage[skins,breakable]{tcolorbox}
\usepackage{graphicx}
\usepackage{colortbl}
\usepackage{booktabs}
\usepackage{tabularx}
\usepackage{multirow}
\usepackage{diagbox} % case d'en-tête diagonale des tableaux à double entrée
% Colonnes extensibles centrée / gauche / droite (C, L, R), souvent utilisées par les modèles
\newcolumntype{C}{>{\centering\arraybackslash}X}
\newcolumntype{L}{>{\raggedright\arraybackslash}X}
\newcolumntype{R}{>{\raggedleft\arraybackslash}X}
\usepackage{array}
\usepackage{multicol}
\usepackage{siunitx}
% parse-numbers=false : accepte aussi \SI{2,5 \times 10^{-3}}{...}
\sisetup{locale=FR,output-decimal-marker={,},group-minimum-digits=4,parse-numbers=false}
\DeclareSIUnit\ohm{\ensuremath{\symup{\Omega}}} % U+2126 absent de la police
\DeclareSIUnit\division{div} % réglages d'oscilloscope (ms/div, V/div)
\DeclareSIUnit\tr{tr} % tours (vitesse de rotation, ex. \SI{1500}{\tr\per\minute})
\DeclareSIUnit\molar{M} % concentration molaire (ex. \SI{3}{\molar}, \SI{1}{\micro\molar})
\DeclareSIUnit\an{an} % année (le modèle confond parfois avec \year, un compteur TeX) — ex. \SI{5}{\kilo\meter\per\an}
\DeclareSIUnit\FCFA{FCFA} % franc CFA (ex. \SI{500}{\FCFA\per\kilo\gram})
\DeclareSIUnit\degreeBrix{\textdegree Brix} % teneur en sucre d'un sirop/jus (réfractométrie)
\DeclareSIUnit\month{mois} % le modèle utilise parfois \month, un compteur TeX, comme unité de durée
\DeclareSIUnit\microliter{\micro\liter} % le modèle écrit parfois \microliter en un seul token
\DeclareSIUnit\million{millions} % dénombrement (ex. \SI{15}{\million\per\mL} spermatozoïdes)
\usepackage{qrcode}
% \degree : commande fréquemment utilisée par les modèles pour un angle ou une
% température en dehors de \SI{}{} (non fournie par les paquets chargés).
\providecommand{\degree}{^{\circ}}
% \qed (fin de démonstration), utilisé par les modèles sans amsthm
\providecommand{\qed}{\hfill$\square$}
% \barre : double barre des valeurs interdites dans un tableau de variations (tkz-tab)
\providecommand{\barre}{\ensuremath{\|}}
% environnement proof (amsthm non chargé) utilisé par les modèles
\ifdefined\proof\else\newenvironment{proof}[1][Démonstration]{\par\noindent\textit{#1.}\ }{\qed\par}\fi
\usepackage{lastpage}
\usepackage{hyperref}
\hypersetup{hidelinks}

% ============================================================
% Palette « Pop » OnBuch+ — mêmes hex que la classe P (plus_ui.dart)
% ============================================================
\definecolor{poporange}{HTML}{F59321}
\definecolor{popdark}{HTML}{E07A0C}
\definecolor{popcream}{HTML}{FFF6E8}
\definecolor{popink}{HTML}{1C1714}
\definecolor{poppurple}{HTML}{7A5AE0}
\definecolor{popgreen}{HTML}{2E9E6A}
\definecolor{poppink}{HTML}{E0457B}
\definecolor{popblue}{HTML}{2F7DE1}
\definecolor{popgold}{HTML}{C98A12}
\definecolor{popmuted}{HTML}{8A7F76}
\definecolor{popline}{HTML}{EADFD2}
\definecolor{popgray}{HTML}{8A7F76}
\colorlet{popgrey}{popgray}
% Alias tolérants (le modèle nomme parfois les couleurs différemment)
\colorlet{popwhite}{white}
\colorlet{popblack}{popink}
\colorlet{popred}{poppink}
\colorlet{popyellow}{popgold}
% Teintes claires (fonds de tableaux/encadrés) — le modèle nomme parfois
% les couleurs avec un suffixe L (ex. popblueL) sans les avoir définies.
\colorlet{poporangeL}{poporange!20}
\colorlet{popdarkL}{popdark!20}
\colorlet{poppurpleL}{poppurple!20}
\colorlet{popgreenL}{popgreen!20}
\colorlet{poppinkL}{poppink!20}
\colorlet{popblueL}{popblue!20}
\colorlet{popgoldL}{popgold!20}
\colorlet{popredL}{popred!20}
\colorlet{popgrayL}{popgray!20}
% Teintes claires pour fonds de tableaux / zones
\colorlet{popblueL}{popblue!10!white}
\colorlet{poporangeL}{poporange!14!white}
\colorlet{popgreenL}{popgreen!12!white}
\colorlet{poppurpleL}{poppurple!12!white}
\colorlet{poppinkL}{poppink!12!white}
\colorlet{popgoldL}{popgold!14!white}

\pagecolor{popcream}

% ============================================================
% Typographie
% ============================================================
\defaultfontfeatures{Ligatures=TeX}
\setmainfont{PlusJakartaSans-Medium.ttf}[Path=../../../fonts/,BoldFont=PlusJakartaSans-Bold.ttf]
% Maths en sans-serif (Fira Math) avec chiffres et lettres droites de Plus
% Jakarta, pour que les formules se fondent dans le texte.
\usepackage[math-style=ISO,bold-style=ISO]{unicode-math}
\setmathfont{FiraMath-Regular.otf}[Path=../../../fonts/]
\setmathfont{PlusJakartaSans-Medium.ttf}[Path=../../../fonts/,range={up/num,up/{Latin,latin},\mathup}]
\usepackage{icomma} % virgule décimale sans espace en mode math
% Caractères mathématiques Unicode absents des polices de texte (Plus Jakarta,
% Archivo Black) : rendus par leur équivalent mathématique, en texte comme en
% formule — sinon ils s'affichent en carré vide (ex. « Arithmétique dans ℤ »).
\usepackage{newunicodechar}
\newunicodechar{ℝ}{\ensuremath{\mathbb{R}}}
\newunicodechar{ℤ}{\ensuremath{\mathbb{Z}}}
\newunicodechar{ℂ}{\ensuremath{\mathbb{C}}}
\newunicodechar{ℕ}{\ensuremath{\mathbb{N}}}
\newunicodechar{ℚ}{\ensuremath{\mathbb{Q}}}
\newunicodechar{𝔻}{\ensuremath{\mathbb{D}}}
\newunicodechar{α}{\ensuremath{\alpha}}
\newunicodechar{β}{\ensuremath{\beta}}
\newunicodechar{γ}{\ensuremath{\gamma}}
\newunicodechar{δ}{\ensuremath{\delta}}
\newunicodechar{ε}{\ensuremath{\varepsilon}}
\newunicodechar{θ}{\ensuremath{\theta}}
\newunicodechar{λ}{\ensuremath{\lambda}}
\newunicodechar{μ}{\ensuremath{\mu}}
\newunicodechar{σ}{\ensuremath{\sigma}}
\newunicodechar{τ}{\ensuremath{\tau}}
\newunicodechar{φ}{\ensuremath{\varphi}}
\newunicodechar{ψ}{\ensuremath{\psi}}
\newunicodechar{ω}{\ensuremath{\omega}}
\newunicodechar{Σ}{\ensuremath{\Sigma}}
% majuscules produites par \MakeUppercase dans les titres (α-aminés -> Α)
\newunicodechar{Α}{\ensuremath{\alpha}}
\newunicodechar{Β}{\ensuremath{\beta}}
\newunicodechar{Γ}{\ensuremath{\Gamma}}
\newunicodechar{Δ}{\ensuremath{\Delta}}
\newunicodechar{Ε}{\ensuremath{\varepsilon}}
\newunicodechar{Θ}{\ensuremath{\Theta}}
\newunicodechar{Λ}{\ensuremath{\Lambda}}
\newunicodechar{Μ}{\ensuremath{\mu}}
\newunicodechar{Τ}{\ensuremath{\tau}}
\newunicodechar{Φ}{\ensuremath{\Phi}}
\newunicodechar{Ψ}{\ensuremath{\Psi}}
\newunicodechar{Ω}{\ensuremath{\Omega}}
\newunicodechar{↦}{\ensuremath{\mapsto}}
\newunicodechar{∈}{\ensuremath{\in}}
\newunicodechar{∉}{\ensuremath{\notin}}
\newunicodechar{∧}{\ensuremath{\wedge}}
\newunicodechar{⇔}{\ensuremath{\Leftrightarrow}}
\newunicodechar{⟺}{\ensuremath{\Longleftrightarrow}}
\newunicodechar{⇒}{\ensuremath{\Rightarrow}}
\newunicodechar{⟹}{\ensuremath{\Longrightarrow}}
\newunicodechar{∘}{\ensuremath{\circ}}
\newunicodechar{∪}{\ensuremath{\cup}}
\newunicodechar{∩}{\ensuremath{\cap}}
\newunicodechar{≡}{\ensuremath{\equiv}}
\newunicodechar{⊂}{\ensuremath{\subset}}
\newunicodechar{⊥}{\ensuremath{\perp}}
\newunicodechar{∃}{\ensuremath{\exists}}
\newunicodechar{∀}{\ensuremath{\forall}}
\newunicodechar{∛}{\ensuremath{\sqrt[3]{\,}}}
\newunicodechar{∜}{\ensuremath{\sqrt[4]{\,}}}
\newunicodechar{⃗}{\ensuremath{{}^{\to}}}
\newunicodechar{ⁿ}{\ifmmode^{n}\else\textsuperscript{\ensuremath{n}}\fi}
\newunicodechar{ˣ}{\ifmmode^{x}\else\textsuperscript{\ensuremath{x}}\fi}
\newunicodechar{ᵇ}{\ifmmode^{b}\else\textsuperscript{\ensuremath{b}}\fi}
\newunicodechar{ʳ}{\ifmmode^{r}\else\textsuperscript{\ensuremath{r}}\fi}
\newunicodechar{ᵘ}{\ifmmode^{u}\else\textsuperscript{\ensuremath{u}}\fi}
\newunicodechar{ʸ}{\ifmmode^{y}\else\textsuperscript{\ensuremath{y}}\fi}
\newunicodechar{ᵃ}{\ifmmode^{a}\else\textsuperscript{\ensuremath{a}}\fi}
\newunicodechar{ᵉ}{\ifmmode^{e}\else\textsuperscript{\ensuremath{e}}\fi}
\newunicodechar{ᵏ}{\ifmmode^{k}\else\textsuperscript{\ensuremath{k}}\fi}
\newunicodechar{ᵖ}{\ifmmode^{p}\else\textsuperscript{\ensuremath{p}}\fi}
\newunicodechar{ᴺ}{\ifmmode^{N}\else\textsuperscript{\ensuremath{N}}\fi}
\newunicodechar{ᶜ}{\ifmmode^{c}\else\textsuperscript{\ensuremath{c}}\fi}
\newunicodechar{ʰ}{\ifmmode^{h}\else\textsuperscript{\ensuremath{h}}\fi}
\newunicodechar{ᵠ}{\ifmmode^{q}\else\textsuperscript{\ensuremath{q}}\fi}
\newunicodechar{ₙ}{\ifmmode_{n}\else\textsubscript{\ensuremath{n}}\fi}
\newunicodechar{ₐ}{\ifmmode_{a}\else\textsubscript{\ensuremath{a}}\fi}
\newunicodechar{ᵢ}{\ifmmode_{i}\else\textsubscript{\ensuremath{i}}\fi}
\newunicodechar{ᵣ}{\ifmmode_{r}\else\textsubscript{\ensuremath{r}}\fi}
\newunicodechar{ₖ}{\ifmmode_{k}\else\textsubscript{\ensuremath{k}}\fi}
\newunicodechar{ᵦ}{\ifmmode_{\beta}\else\textsubscript{\ensuremath{\beta}}\fi}
\newunicodechar{ₓ}{\ifmmode_{x}\else\textsubscript{\ensuremath{x}}\fi}
\newfontfamily\popdisplay{ArchivoBlack-Regular.ttf}[Path=../../../fonts/]
\newfontfamily\poplogo{SpaceGrotesk-Bold.ttf}[Path=../../../fonts/]
\newfontfamily\popsemi{PlusJakartaSans-SemiBold.ttf}[Path=../../../fonts/]
\newfontfamily\popxbold{PlusJakartaSans-ExtraBold.ttf}[Path=../../../fonts/]
\newfontfamily\popxboldls{PlusJakartaSans-ExtraBold.ttf}[Path=../../../fonts/,LetterSpace=3.0]

\setlist[enumerate]{leftmargin=1.7em,itemsep=5pt,topsep=4pt}
\setlist[itemize]{leftmargin=1.7em,itemsep=4pt,topsep=4pt,label={\color{poporange}\textbullet}}
\setlength{\parindent}{0pt}
\setlength{\parskip}{5pt}
\renewcommand{\arraystretch}{1.25}

% pgfplots : style Pop par défaut
\pgfplotsset{
  every axis/.append style={axis line style={popink,line width=1pt},tick style={popink},
    grid style={popline,line width=0.5pt},label style={font=\small},tick label style={font=\footnotesize}},
  popcurve/.style={poporange,line width=1.8pt,no markers,smooth},
}
% Même style disponible dans un \draw TikZ simple (hors pgfplots)
\tikzset{popcurve/.style={poporange,line width=1.8pt}}
\tikzset{popgrid/.style={popline,very thin}}
\colorlet{popcurve}{poporange} % le nom du style est parfois utilisé comme couleur

% ============================================================
% Pill / squiggle
% ============================================================
\newcommand{\poppill}[3]{%
  \tikz[baseline=-0.55ex]\node[draw=popink,line width=1.2pt,rounded corners=2.6pt,%
    fill=#2,inner xsep=6.5pt,inner ysep=3pt]{%
    \popxboldls\fontsize{7}{7}\selectfont\color{#3}\MakeUppercase{#1}};%
}
\newcommand{\popsquiggle}[1]{%
  \tikz[baseline=-0.4ex]{
    \draw[#1,line width=2.6pt,line cap=round]
      plot [smooth] coordinates {(0,0.09) (0.34,0.01) (0.68,0.21) (1.02,0.01) (1.36,0.10)};
  }%
}
% Mot-clé surligné (marqueur orange sous le texte)
\newcommand{\cle}[1]{{\popxbold\color{popdark}#1}}

% ============================================================
% En-tête / pied
% ============================================================
\pagestyle{fancy}
\fancyhf{}
\renewcommand{\headrulewidth}{0pt}
\renewcommand{\footrulewidth}{0pt}
\lhead{\raisebox{-0.15ex}{\poplogo\fontsize{13}{13}\selectfont\color{popink}OnBuch\color{poporange}+}}
\rhead{\poppill{\DOCMATIERE\ \textbullet\ \DOCNIVEAU\ \textbullet\ Leçon \DOCLECON}{white}{popink}}
\lfoot{\popxbold\fontsize{7.6}{7.6}\selectfont\color{popmuted}\MakeUppercase{Cours signé OnBuch+}\ \textbullet\ {\popsemi\DOCTITRE}}
\rfoot{\tikz[baseline=-0.6ex]\node[rounded corners=8.5pt,fill=popink,minimum height=17pt,minimum width=17pt,inner xsep=5pt,inner sep=0]{\popxbold\fontsize{8}{8}\selectfont\color{popcream}\thepage\,/\,\pageref*{LastPage}};}

% ============================================================
% Titres
% ============================================================
\newcommand{\chaptitre}[2]{%
  \begin{tcolorbox}[enhanced,colback=poporange,colframe=popink,boxrule=2.6pt,arc=11pt,
      drop shadow=popink,left=15pt,right=15pt,top=12pt,bottom=11pt,before skip=3pt,after skip=18pt]
    \poppill{#2}{popink}{popcream}\par\vspace{9pt}
    {\raggedright\hyphenpenalty=10000\exhyphenpenalty=10000\popdisplay\fontsize{22}{24}\selectfont\color{popcream}\MakeUppercase{#1}\par}
  \end{tcolorbox}
}
% Section numérotée : pastille orange avec le numéro + titre Archivo
\newcounter{popsec}
\newcommand{\coursec}[1]{%
  \stepcounter{popsec}\setcounter{popsub}{0}%
  \par\vspace{16pt}\noindent
  \begin{minipage}{\linewidth}
  \tikz[baseline=-0.7ex]\node[draw=popink,line width=1.6pt,fill=poporange,rounded corners=4pt,minimum size=22pt,inner sep=0]{\popdisplay\fontsize{12}{12}\selectfont\color{popcream}\thepopsec};%
  \hspace{8pt}{\popdisplay\fontsize{15}{17}\selectfont\color{popink}\MakeUppercase{#1}}\par
  \vspace{4pt}\noindent{\color{poporange}\rule{36pt}{3.6pt}}
  \end{minipage}\par\nopagebreak\vspace{8pt}
}
\newcounter{popsub}
\newcommand{\courssub}[1]{%
  \stepcounter{popsub}%
  \par\vspace{9pt}\noindent{\popxbold\fontsize{11.5}{13}\selectfont\color{popink}{\color{poporange}\thepopsec.\thepopsub}\hspace{6pt}#1}\par\nopagebreak\vspace{4pt}
}

% ============================================================
% Boîtes pédagogiques — même recette : fond blanc, bordure épaisse colorée,
% ombre dure encre, badge plein. Titre optionnel [#1].
% ============================================================
\tcbset{popbase/.style={breakable,enhanced,colback=white,boxrule=2.3pt,arc=9pt,
  shadow={3pt}{-3pt}{0pt}{popink},left=14pt,right=14pt,top=11pt,bottom=11pt,before skip=11pt,after skip=15pt,
  fontupper=\popsemi\fontsize{10.6}{14.5}\selectfont\color{popink},
  fontlower=\popsemi\fontsize{10.6}{14.5}\selectfont\color{popink}}}
% Coche dessinée (le glyphe ✓ n'existe pas dans les polices du document)
\newcommand{\popcheck}[1]{\tikz[baseline=-0.5ex]\draw[#1,line width=1.6pt,line cap=round,line join=round](-0.13,0.0)--(-0.04,-0.09)--(0.13,0.1);}
\providecommand{\coche}{\popcheck{popgreen}}
\providecommand{\resultat}[1]{\par\smallskip\noindent\fcolorbox{popgreen}{popgreenL}{\parbox{\dimexpr\linewidth-2\fboxsep-2\fboxrule\relax}{\textbf{Résultat.} #1}}\par\smallskip}
\newcommand{\popboxhead}[3]{\poppill{#1}{#2}{white}\ifx\relax#3\relax\else\hspace{6pt}{\popxbold\fontsize{10.6}{12}\selectfont\color{popink}#3}\fi\par\vspace{7pt}}

\newtcolorbox{aretenir}[1][]{popbase,colframe=popgreen,before upper={\popboxhead{À retenir}{popgreen}{#1}}}
\newtcolorbox{exemplebox}[1][]{popbase,colframe=poporange,before upper={\popboxhead{Exemple}{poporange}{#1}}}
\newtcolorbox{definition}[1][]{popbase,colframe=popblue,before upper={\popboxhead{Définition}{popblue}{#1}}}
\newtcolorbox{propriete}[1][]{popbase,colframe=poppurple,before upper={\popboxhead{Propriété}{poppurple}{#1}}}
\newtcolorbox{methode}[1][]{popbase,colframe=popgold,before upper={\popboxhead{Méthode}{popgold}{#1}}}
\newtcolorbox{attention}[1][]{popbase,colframe=poppink,before upper={\popboxhead{Attention}{poppink}{#1}}}
\newtcolorbox{experience}[1][]{popbase,colframe=popblue,colback=popblueL,before upper={\popboxhead{Expérience}{popblue}{#1}}}
% « Le savais-tu ? » : carte violette pleine (contexte, culture, Cameroun)
\newtcolorbox{savaistu}[1][]{popbase,colback=poppurple,colframe=popink,
  fontupper=\popsemi\fontsize{10.4}{14.2}\selectfont\color{white},
  before upper={\poppill{Le savais-tu\,?}{white}{poppurple}\ifx\relax#1\relax\else\hspace{6pt}{\popxbold\color{white}#1}\fi\par\vspace{7pt}}}
% Exercice résolu : énoncé en haut, \tcblower puis la solution sur fond crème
\newcounter{popexo}\newcounter{popexr}
\newtcolorbox{exoresolu}[1][]{popbase,colframe=poporange,colbacklower=popcream,
  segmentation style={popink,line width=1.2pt,dashed},
  before upper={\stepcounter{popexr}\popboxhead{Exercice résolu \thepopexr}{poporange}{#1}},
  before lower={\poppill{Solution}{popink}{popcream}\par\vspace{6pt}}}
% Exercice (non résolu en place) — la correction est dans la partie Corrigés
\newtcolorbox{exercice}[1][]{popbase,colframe=popink,
  before upper={\stepcounter{popexo}\popboxhead{Exercice \thepopexo}{popink}{#1}}}
% Corrigé
\newtcolorbox{corrige}[1][]{popbase,colframe=popgreen,colback=popgreenL,
  before upper={\popboxhead{Corrigé}{popgreen}{#1}}}
% Objectifs / prérequis / bilan
\newtcolorbox{objectifs}{popbase,colframe=popink,colback=poporangeL,
  before upper={\popboxhead{Objectifs de la leçon}{popink}{}}}
\newtcolorbox{prerequis}{popbase,colframe=popink,colback=popblueL,
  before upper={\popboxhead{Prérequis}{popblue}{}}}
\newtcolorbox{bilan}{popbase,colframe=popink,colback=white,boxrule=2.8pt,
  before upper={\popboxhead{Fiche bilan}{poporange}{L'essentiel à connaître pour l'examen}}}
% Figure encadrée avec légende
\newtcolorbox{popfigure}[1][]{enhanced,breakable=false,colback=white,colframe=popline,boxrule=1.4pt,arc=8pt,
  left=10pt,right=10pt,top=10pt,bottom=8pt,before skip=10pt,after skip=12pt,halign=center,
  lower separated=false,halign lower=center,
  fontlower=\popsemi\fontsize{9}{11.5}\selectfont\color{popmuted},
  #1}
\newcommand{\legende}[1]{\par\vspace{6pt}{\popsemi\fontsize{9}{11.5}\selectfont\color{popmuted}#1}}

% ============================================================
% Signature OnBuch+
% ============================================================
\newcommand{\poplogoinline}[1]{{\poplogo\fontsize{#1}{#1}\selectfont\color{popink}OnBuch\color{poporange}+}}
\newcommand{\signatureonbuch}{%
  \begin{tcolorbox}[enhanced,colback=popink,colframe=popink,boxrule=2.2pt,arc=10pt,
      drop shadow=poporange,left=14pt,right=14pt,top=10pt,bottom=10pt,before skip=0pt,after skip=0pt]
    \tikz[baseline=-0.7ex]\node[circle,fill=popgreen,draw=popcream,line width=1.3pt,minimum size=22pt,inner sep=0]{\popcheck{white}};%
    \hspace{9pt}\begin{minipage}[c]{0.62\linewidth}
      {\popxbold\fontsize{12}{13}\selectfont\color{popcream}Signé\ \textbullet\ {\poplogo OnBuch\color{poporange}+}}\par\vspace{2pt}
      {\raggedright\popsemi\fontsize{8}{10.5}\selectfont\color{popline}\MakeUppercase{Équipe pédagogique OnBuch+}\\\MakeUppercase{Conforme au programme officiel MINESEC}\par}
    \end{minipage}\hfill
    {\poplogo\fontsize{22}{22}\selectfont\color{popcream}OnBuch\color{poporange}+}
  \end{tcolorbox}%
}

% ============================================================
% Page de garde
% #1 titre · #2 matière · #3 niveau · #4 module · #5 n° leçon · #6 durée ·
% #7 description / objectifs (texte ou itemize)
% ============================================================
\newcommand{\pagegarde}[7]{%
  \thispagestyle{empty}
  \newgeometry{a4paper,margin=1.7cm,top=1.6cm,bottom=1.4cm}%
  \begin{tikzpicture}[remember picture,overlay]
    \fill[poporange!18] ([xshift=-3.2cm,yshift=-3.6cm]current page.north east) circle (4.6cm);
    \fill[poppurple!14] ([xshift=2.4cm,yshift=7.6cm]current page.south west) circle (3.8cm);
  \end{tikzpicture}%
  {\poplogo\fontsize{30}{30}\selectfont\color{popink}OnBuch\color{poporange}+}\hfill
  \raisebox{0.6ex}{\poppill{Cours signé \textbullet\ Premium}{popink}{popcream}}\par
  \vspace{1pt}\popsquiggle{poppurple}\par
  \vspace{8pt}
  \begin{tcolorbox}[enhanced,colback=poporange,colframe=popink,boxrule=2.8pt,arc=13pt,
      drop shadow=popink,left=18pt,right=18pt,top=12pt,bottom=13pt,after skip=10pt]
    \poppill{#2 \textbullet\ #3}{popink}{popcream}\hspace{5pt}\poppill{#4}{white}{popink}\par\vspace{8pt}
    {\popdisplay\fontsize{14}{14}\selectfont\color{popink}LEÇON #5}\par\vspace{4pt}
    {\raggedright\hyphenpenalty=10000\exhyphenpenalty=10000\popdisplay\fontsize{25}{27}\selectfont\color{popcream}\MakeUppercase{#1}\par}
    \vspace{8pt}
    {\popxbold\fontsize{9.5}{9.5}\selectfont\color{popink}\MakeUppercase{Durée conseillée : #6}}
  \end{tcolorbox}
  \begin{tcolorbox}[enhanced,colback=white,colframe=popink,boxrule=2.2pt,arc=10pt,
      drop shadow=popink,left=16pt,right=16pt,top=10pt,bottom=10pt,after skip=8pt]
    \poppill{Dans cette leçon}{poppurple}{white}\par\vspace{5pt}
    \popsemi\fontsize{9.3}{12.6}\selectfont\color{popink}#7
  \end{tcolorbox}
  \noindent\begin{minipage}[t]{0.487\linewidth}
    \begin{tcolorbox}[enhanced,colback=popgreen,colframe=popink,boxrule=2.2pt,arc=10pt,
        drop shadow=popink,left=11pt,right=11pt,top=10pt,bottom=10pt,height=3.1cm,valign=center]
      \tcbox[colback=white,colframe=popink,boxrule=1.3pt,arc=5pt,boxsep=0pt,nobeforeafter,
        left=3pt,right=3pt,top=3pt,bottom=3pt,valign=center]{\qrcode[height=1.9cm]{https://play.google.com/store/apps/details?id=com.luvvix.onbuch&hl=fr}}%
      \hspace{8pt}\begin{minipage}[c]{0.5\linewidth}
        {\popdisplay\fontsize{11}{12.5}\selectfont\color{white}\MakeUppercase{Télécharge OnBuch}}\par\vspace{4pt}
        {\raggedright\popsemi\fontsize{8.4}{11}\selectfont\color{white}Gratuit \textbullet\ Google Play\\Tuteur IA, annales, concours, résultats\par}
      \end{minipage}
    \end{tcolorbox}
  \end{minipage}\hfill
  \begin{minipage}[t]{0.487\linewidth}
    \begin{tcolorbox}[enhanced,colback=poppurple,colframe=popink,boxrule=2.2pt,arc=10pt,
        drop shadow=popink,left=11pt,right=11pt,top=10pt,bottom=10pt,height=3.1cm,valign=center]
      \tikz[baseline=-0.7ex]\node[circle,fill=white,draw=popink,line width=1.3pt,minimum size=1.6cm,inner sep=0]{\popdisplay\fontsize{24}{24}\selectfont\color{poppurple}+};%
      \hspace{8pt}\begin{minipage}[c]{0.6\linewidth}
        {\popdisplay\fontsize{11}{12.5}\selectfont\color{white}\MakeUppercase{Passe à OnBuch+}}\par\vspace{4pt}
        {\raggedright\popsemi\fontsize{8.4}{11}\selectfont\color{white}Tous les cours signés, Tuteur IA illimité, fiches PDF — onglet OnBuch+ de l'app\par}
      \end{minipage}
    \end{tcolorbox}
  \end{minipage}
  \vspace{8pt}
  \popcontenu
  \vfill
  \signatureonbuch
  \restoregeometry
  \newpage
}

% Rangée « Ce cours contient » (page de garde)
\newcommand{\poptile}[3]{% #1 couleur #2 gros texte #3 légende
  \begin{tcolorbox}[enhanced,colback=#1,colframe=popink,boxrule=2pt,arc=9pt,shadow={2.5pt}{-2.5pt}{0pt}{popink},
      left=6pt,right=6pt,top=9pt,bottom=9pt,halign=center,valign=center,height=1.75cm,width=\linewidth,nobeforeafter]
    {\popdisplay\fontsize{12.5}{13}\selectfont\color{white}\MakeUppercase{#2}}\par\vspace{3pt}
    {\centering\popsemi\fontsize{7.8}{9.5}\selectfont\color{white}#3\par}
  \end{tcolorbox}}
\newcommand{\popcontenu}{%
  \par\noindent{\popxboldls\fontsize{7.5}{7.5}\selectfont\color{popmuted}CE COURS SIGNÉ CONTIENT}\par\vspace{7pt}
  \noindent\begin{minipage}[t]{0.235\linewidth}\poptile{popblue}{Cours}{complet, illustré, pas à pas}\end{minipage}\hfill
  \begin{minipage}[t]{0.235\linewidth}\poptile{poporange}{Méthodes}{et exercices résolus}\end{minipage}\hfill
  \begin{minipage}[t]{0.235\linewidth}\poptile{poppink}{Exercices}{type examen + corrigés}\end{minipage}\hfill
  \begin{minipage}[t]{0.235\linewidth}\poptile{popgreen}{Fiche bilan}{l'essentiel à retenir}\end{minipage}\par}

% Sommaire
\newtcolorbox{sommaire}{enhanced,colback=white,colframe=popink,boxrule=2.2pt,arc=10pt,
  drop shadow=popink,left=18pt,right=18pt,top=13pt,bottom=13pt,before skip=6pt,after skip=20pt,
  before upper={\poppill{Sommaire}{poppurple}{white}\par\vspace{9pt}\popsemi\fontsize{10.6}{14.5}\selectfont\color{popink}}}

% Signature de fin de cours — poussée tout en bas de la dernière page
\newcommand{\findecours}{%
  \par\vfill\nopagebreak
  \begin{center}\popsquiggle{poporange}\end{center}\vspace{4pt}
  \signatureonbuch
}
\AtBeginDocument{\def\llbracket{\mathopen{[\mkern-3.5mu[}}\def\rrbracket{\mathclose{]\mkern-3.5mu]}}} % intervalles d'entiers (glyphe ⟦ absent de la police)
% Glyphes absents de Fira Math : redéfinis à partir de glyphes disponibles
\AtBeginDocument{%
  \def\setminus{\mathbin{\backslash}}\let\smallsetminus\setminus
  \def\vdots{\mathord{\vbox{\baselineskip3.2pt\lineskiplimit0pt\kern4pt\hbox{.}\hbox{.}\hbox{.}}}}%
  \def\ddots{\mathinner{\mkern1mu\raise6.5pt\hbox{.}\mkern2mu\raise3.6pt\hbox{.}\mkern2mu\raise0.7pt\hbox{.}\mkern1mu}}%
  \def\triangle{\mathord{\scalebox{0.9}{$\symup{\Delta}$}}}%
  \def\nabla{\mathord{\raisebox{\depth}{\scalebox{1}[-1]{$\symup{\Delta}$}}}}%
  \def\checkmark{\popcheck{popdark}}%
  \def\star{\mathbin{\ast}}\def\bigstar{\mathord{\ast}}%
  \def\diamond{\mathbin{\scalebox{0.6}{$\Diamond$}}}%
  \def\bigcup{\mathop{\mathchoice{\raisebox{-0.3em}{\scalebox{1.7}{$\cup$}}}{\scalebox{1.25}{$\cup$}}{\cup}{\cup}}\displaylimits}%
  \def\bigcap{\mathop{\mathchoice{\raisebox{-0.3em}{\scalebox{1.7}{$\cap$}}}{\scalebox{1.25}{$\cap$}}{\cap}{\cap}}\displaylimits}%
  \def\lesseqqgtr{\mathrel{\lessgtr}}\def\bigoplus{\mathop{\oplus}\displaylimits}%
}
\providecommand{\ssi}{si et seulement si} % abréviation fréquente
\AtBeginDocument{\providecommand{\wideparen}{\overparen}} % arc de cercle

\begin{document}

\pagegarde{Liberté et Responsabilité}{Philosophie}{Tle Tech}{Philosophie – classe de Terminale technique}{N15}{2 séances (fiche de progression harmonisée)}{Cette leçon interroge deux notions fondamentales de la vie morale et juridique : la liberté, pouvoir d'agir par soi-même, et la responsabilité, obligation de répondre de ses actes. À partir de situations concrètes camerounaises (justice, vie scolaire, conduite en société), elle montre que la liberté est le fondement de la responsabilité et prépare à l'épreuve de dissertation du Baccalauréat technique.
\vspace{4pt}
\begin{multicols}{2}\small
\begin{itemize}
\item À la fin de cette leçon, je sais définir les concepts de liberté et de responsabilité avec leurs distinctions essentielles
\item À la fin de cette leçon, je sais distinguer la liberté physique, la liberté politique et la liberté morale
\item À la fin de cette leçon, je sais distinguer la responsabilité morale, la responsabilité civile et la responsabilité pénale
\item À la fin de cette leçon, je sais expliquer en quoi la liberté est le fondement de la responsabilité
\item À la fin de cette leçon, je sais analyser les conditions d'attribution de la responsabilité (conscience, liberté, causalité)
\item À la fin de cette leçon, je sais appliquer les notions de liberté et de responsabilité à des situations concrètes de la vie courante au Cameroun
\end{itemize}\end{multicols}}

\begin{sommaire}
\begin{multicols}{2}
\begin{enumerate}
\item Introduction : problématisation et définition de la liberté
\item Les formes de la liberté
\item La responsabilité : définition et formes
\item La liberté, fondement de la responsabilité
\item Applications concrètes au quotidien camerounais
\item Méthode de la dissertation et synthèse
\item Activité d'intégration
\item Méthodes et exercices résolus
\item Exercices
\item Corrigés des exercices
\item Fiche bilan
\end{enumerate}
\end{multicols}
\end{sommaire}

\begin{objectifs}
\begin{itemize}
\item À la fin de cette leçon, je sais définir les concepts de liberté et de responsabilité avec leurs distinctions essentielles
\item À la fin de cette leçon, je sais distinguer la liberté physique, la liberté politique et la liberté morale
\item À la fin de cette leçon, je sais distinguer la responsabilité morale, la responsabilité civile et la responsabilité pénale
\item À la fin de cette leçon, je sais expliquer en quoi la liberté est le fondement de la responsabilité
\item À la fin de cette leçon, je sais analyser les conditions d'attribution de la responsabilité (conscience, liberté, causalité)
\item À la fin de cette leçon, je sais appliquer les notions de liberté et de responsabilité à des situations concrètes de la vie courante au Cameroun
\item À la fin de cette leçon, je sais construire une problématique et rédiger une introduction de dissertation sur ce thème
\item À la fin de cette leçon, je sais rédiger et justifier une démarche argumentée, une réponse ou une conclusion personnelle
\end{itemize}
\end{objectifs}

\begin{prerequis}
\begin{itemize}
\item Notion de devoir et de morale (classes de Première)
\item Notion de conscience et d'inconscient (leçon antérieure de Terminale)
\item Notion d'État et de droit (vie de cité, éducation civique)
\item Vocabulaire philosophique de base : concept, argument, thèse, objection
\item Méthode de l'analyse de sujet de dissertation (début d'année de Terminale)
\end{itemize}
\end{prerequis}

\newpage

\coursec{Introduction : problématisation et définition de la liberté}

\courssub{Une situation d'accroche : le conducteur de moto-taxi à Yaoundé}

Imaginez la scène, banale sur les artères de la capitale. Un jeune conducteur de moto-taxi, pressé par un client en retard pour un concours administratif, grille un feu rouge au carrefour \emph{Poste Central}. Pour gagner les précieuses secondes qui séparent le client de son centre d'examen, il accélère au lieu de freiner. Un piéton traverse, confiant dans le vert qui lui est acquis. L'impact est violent. Le piéton est grièvement blessé, la moto gît au milieu de l'asphalte brûlant.

Devant le commissariat du quartier, le conducteur tente de se justifier, la voix tremblante : \og Ce n'est pas ma faute, Monsieur l'Agent ! Le client me pressait, il m'a même donné 500 francs de pourboire pour aller vite. Et puis, les freins étaient usés, le patron ne veut pas les changer. Je n'ai pas choisi d'avoir cet accident ! \fg{}

Cette situation, tirée du quotidien camerounais, condense le problème philosophique qui nous occupe : \textbf{la responsabilité suppose-t-elle la liberté ?} Le conducteur invoque des causes extérieures (la pression du client, l'état de la machine) pour s'exonérer. Mais a-t-il \emph{réellement} agi sans liberté ? Au moment précis où le feu est passé au rouge, n'avait-il pas le \emph{pouvoir} de freiner, de klaxonner, de dévier sa trajectoire ? S'il a choisi d'accélérer, même sous la contrainte, n'est-il pas l'auteur de ce choix ?

La question s'élargit immédiatement : si un mineur de 14 ans, envoyé par son père faire une course, avait conduit la moto, le jugerait-on de la même manière ? Si le conducteur était en proie à une crise d'épilepsie soudaine, ou sous l'emprise d'une drogue administrée à son insu, la responsabilité disparaîtrait-elle ? Partout, le critère implicite est le même : \textbf{on ne tient pour responsable que celui qui a pu faire autrement, c'est-à-dire celui qui était libre.}

\courssub{Étymologie et sens premier : de l'esclave au citoyen}

Pour penser rigoureusement cette articulation, commençons par les mots. La langue latine, matrice de notre vocabulaire philosophique, nous offre deux pistes lumineuses.

\begin{savaistu}[L'origine juridique de la responsabilité]
Le mot \cle{responsabilité} dérive du latin \emph{responsabilitas}, lui-même issu du verbe \emph{respondere} : \og répondre de \fg{}, \og garantir \fg{}. Dans le droit romain archaïque, l'accusé devait \emph{respondere} : se présenter devant le préteur, répondre aux accusations, \emph{garantir} la réparation du dommage causé. Être responsable, c'est donc originellement \textbf{devoir rendre des comptes pour ses actes devant une instance de jugement}. Cela implique une capacité : celle de comprendre la charge qui pèse sur soi et d'y répondre. Un esclave, au sens strict romain, ne \emph{respondet} pas : son maître répond pour lui. La responsabilité marque ainsi l'entrée dans la personnalité juridique et morale.
\end{savaistu}

Le terme \cle{liberté} vient du latin \emph{libertas}, nom abstrait formé sur \emph{liber} (\og libre \fg{}). \emph{Liber} s'oppose radicalement à \emph{servus} (l'esclave). \emph{Libertas} désigne d'abord un \textbf{statut juridique} : la condition de celui qui n'est la propriété de personne, qui dispose de sa personne (\emph{sui juris}), qui peut aller et venir, contracter, témoigner. Être libre, c'est \emph{ne pas être esclave}. Mais très vite, la philosophie (les Stoïciens, puis les chrétiens, puis les Modernes) déplace le sens : de la \emph{liberté statutaire} (absence de maître humain), on passe à la \emph{liberté métaphysique} (absence de déterminisme absolu, capacité de s'autodéterminer).

\begin{definition}[Liberté : définition générale]
La \cle{liberté} est le \textbf{pouvoir d'agir sans contrainte extérieure et de se déterminer soi-même}. Elle comporte deux dimensions indissociables :
\begin{itemize}
    \item \textbf{Négative (liberté de faire) :} absence d'obstacles physiques, juridiques ou sociaux qui empêchent l'exécution de l'action voulue (pas de chaînes, pas de lois liberticides, pas de menace immédiate de mort).
    \item \textbf{Positive (liberté de vouloir / autonomie) :} capacité à être la source de sa propre volonté, à choisir ses fins et ses moyens en toute lucidité, sans être poussé par des pulsions irréfléchies, des conditionnements inconscients ou une contrainte psychologique irrésistible.
\end{itemize}
\end{definition}

\courssub{Le libre arbitre : la faculté de choisir entre plusieurs possibles}

Au cœur de la liberté positive se niche une notion technique, cruciale pour la responsabilité : le \cle{libre arbitre}. L'expression vient du latin \emph{liberum arbitrium} (\og libre jugement \fg{}, \og libre choix \fg{}). Elle désigne la \textbf{faculté de l'esprit à se porter vers l'un de deux contraires (ou de plusieurs possibles) sans y être déterminé par une cause antérieure nécessaire}.

\begin{definition}[Libre arbitre]
Le \cle{libre arbitre} est la \textbf{faculté de choisir délibérément entre plusieurs actions possibles}. Il ne s'agit pas de faire \emph{n'importe quoi} (ce serait du caprice), mais de pouvoir, face à une situation donnée, \emph{suspendre son jugement}, peser les raisons, et opter pour l'une des issues ouvertes. C'est ce \emph{pouvoir de faire autrement} qui fonde l'imputation morale et juridique.
\end{definition}

Pour visualiser ce mécanisme, observons l'arbre de décision qui s'offre au conducteur au moment précis où le feu passe au rouge :

\begin{popfigure}
\begin{tikzpicture}[
    scale=0.85,
    every node/.style={transform shape},
    decision/.style={diamond, draw=popblue, thick, fill=popblue!10, text width=3.2cm, align=center, inner sep=4pt},
    action/.style={rectangle, draw=popgreen, thick, fill=popgreen!10, text width=3.5cm, align=center, inner sep=4pt, rounded corners=4pt},
    consequence/.style={rectangle, draw=poporange, thick, fill=poporange!10, text width=3.5cm, align=center, inner sep=4pt, rounded corners=4pt},
    arrow/.style={-Stealth, thick, popdark},
    label/.style={midway, fill=white, inner sep=2pt, font=\small, popdark}
]
% Nœuds
\node[decision, align=center] (start){Feu rouge \\ Client pressé \\ Freins usés};
\node[action, align=center] (frein) [below left=1.8cm and 2.5cm of start]{FREINER \\ (Respecter le code,\\ risquer le mécontentement du client)};
\node[action, align=center] (klaxon) [below=1.8cm of start]{KLAXONNER \\ (Avertir le piéton,\\ tenter de passer)};
\node[action, align=center] (accel) [below right=1.8cm and 2.5cm of start]{ACCÉLÉRER \\ (Obéir au client,\\ prendre le risque maximal)};
\node[consequence, align=center] (cfrein) [below left=1.5cm and 2.5cm of frein]{Arrêt sécurisé \\ Responsabilité : nulle};
\node[consequence, align=center] (cklaxon) [below=1.5cm of klaxon]{Piéton averti \\ Risque réduit \\ Responsabilité : partagée ?};
\node[consequence, align=center] (caccel) [below right=1.5cm and 2.5cm of accel]{ACCIDENT GRAVE \\ Responsabilité : PLEINE};

% Flèches
\draw[arrow] (start) -- node[label, left] {Choix 1} (frein);
\draw[arrow] (start) -- node[label, above] {Choix 2} (klaxon);
\draw[arrow] (start) -- node[label, right] {Choix 3} (accel);

\draw[arrow] (frein) -- (cfrein);
\draw[arrow] (klaxon) -- (cklaxon);
\draw[arrow] (accel) -- (caccel);

% Annotations
\node[above=0.3cm of start, font=\bfseries\large, popdark] {Libre Arbitre : Moment de la décision};
\node[below=0.2cm of cfrein, font=\small\itshape, popgreen] {Action libre \& responsable};
\node[below=0.2cm of cklaxon, font=\small\itshape, poporange] {Action risquée};
\node[below=0.2cm of caccel, font=\small\itshape, poporange] {Action libre mais fautive};
\end{tikzpicture}
\legende{Schéma du libre arbitre : face à une situation contrainte (feu rouge, client, freins), le conducteur dispose de plusieurs possibles. Le choix effectif (ici accélérer) engage sa responsabilité car il \emph{aurait pu} choisir autrement (freiner).}
\end{popfigure}

Cet arbre montre que la contrainte (client, freins) \emph{réduit} l'éventail des possibles (freiner est plus difficile, plus risqué socialement), mais ne l'\emph{annihile} pas. Le libre arbitre, c'est cette marge de manœuvre résiduelle, cet \emph{entre-deux} où l'agent se saisit de lui-même.

\courssub{Distinction fondamentale : liberté de faire et liberté de vouloir}

C'est ici que se joue toute la subtilité de la responsabilité. La philosophie classique (Descartes, Kant, Alain) distingue deux niveaux qu'il ne faut jamais confondre.

\begin{propriete}[Les deux étages de la liberté]
\begin{enumerate}
    \item \textbf{Liberté de faire (liberté d'exécution, \emph{libertas exercitii}) :} C'est le pouvoir d'exécuter l'action que l'on a voulue. Elle dépend de conditions \emph{extérieures} : absence de chaînes, de paralysie, de lois interdisant l'acte, de force majeure physique. \emph{Exemple :} Je veux lever le bras ; si je ne suis pas attaché, je le lève. Je suis libre de faire.
    \item \textbf{Liberté de vouloir (liberté de spécification, \emph{libertas specificationis}, autonomie) :} C'est le pouvoir de \emph{déterminer sa propre volonté}, de choisir \emph{quoi} vouloir. Elle dépend de conditions \emph{intérieures} : raison lucide, absence de passion dévastatrice, de névrose compulsive, de conditionnement totalitaire. \emph{Exemple :} Je veux lever le bras \emph{pour voter} ou \emph{pour frapper}. Le sens moral de l'acte réside ici.
\end{enumerate}
\end{propriete}

\begin{attention}[Piège majeur : ne pas confondre Liberté et Licence]
C'est l'erreur la plus fréquente chez les élèves (et dans le langage courant).
\begin{itemize}
    \item \textbf{Liberté} = pouvoir choisir \emph{en conscience} et \emph{assumer} ses choix. Elle s'exerce \emph{dans} un cadre (lois, morale, respect d'autrui). \og Ma liberté s'arrête où commence celle des autres \fg{}.
    \item \textbf{Licence} (ou libertinage) = faire \emph{tout ce que l'on veut}, sans règle, sans frein, en niant les limites de la réalité et les droits d'autrui.
\end{itemize}
\emph{Le conducteur de moto-taxi a exercé une licence (il a fait ce qu'il voulait, n'importe comment), mais il a engagé sa liberté (il a choisi d'accélérer). C'est parce qu'il était libre qu'il est responsable, non parce qu'il était "sans loi".}
\end{attention}

\courssub{Exemple concret : le lycéen de Douala face à son choix}

\begin{exemplebox}[Libre arbitre et responsabilité étudiante]
\textbf{Situation :} Armand, élève en Terminale T au Lycée de New-Bell (Douala), a un devoir de mathématiques crucial pour le Probatoire le lendemain. Ses amis l'invitent au stade pour la finale de la coupe du Cameroun. Il hésite longuement. Il pèse le pour (joie, amitié, détente) et le contre (échec probable, déception parentale, avenir compromis). Il décide finalement de rester réviser. Il rate le match mais réussit son devoir.

\textbf{Analyse :}
\begin{itemize}
    \item \textbf{Liberté de faire :} Armand \emph{peut} physiquement aller au stade (pas de grille fermée, pas de maladie). Il \emph{peut} aussi physiquement ouvrir son cahier.
    \item \textbf{Liberté de vouloir / Libre arbitre :} C'est le moment de l'hésitation. Il \emph{délibère}. Il n'est pas \emph{poussé} par une pulsion irrépressible (addiction au foot) ni \emph{contraint} par une menace (pistolet sur la tempe). Il \emph{choisit} sa fin (la réussite) et son moyen (la révision).
    \item \textbf{Responsabilité :} S'il avait raté son devoir pour aller au stade, il en aurait été \emph{responsable} (il l'aurait choisi). Comme il a réussi, il en est \emph{l'auteur méritant}. Dans les deux cas, la responsabilité suit la liberté de vouloir.
\end{itemize}
\end{exemplebox}

\courssub{Problématique de la leçon : la liberté, condition nécessaire et/ou suffisante ?}

Au terme de cette introduction, nous pouvons formuler rigoureusement le problème qui guidera toute notre réflexion :

\begin{aretenir}[Problématique centrale]
\begin{enumerate}
    \item \textbf{La liberté est-elle la condition \emph{nécessaire} de la responsabilité ?} (Pas de liberté $\Rightarrow$ pas de responsabilité. Si je suis déterminé par le destin, mes gènes, la société, l'inconscient, suis-je encore responsable ?)
    \item \textbf{La liberté est-elle la condition \emph{suffisante} de la responsabilité ?} (Liberté $\Rightarrow$ responsabilité pleine et entière ? Si je choisis librement le mal, suis-je \emph{toujours} 100\% responsable ? La contrainte atténue-t-elle ou supprime-t-elle la responsabilité ?)
\end{enumerate}
\end{aretenir}

Notre conducteur de moto-taxi nous a montré que la situation réelle est un \emph{mélange} : il y a de la contrainte (client, freins, pauvreté du patron) et il y a de la liberté (le choix final d'accélérer). La philosophie, et le droit, devront trancher : \textbf{comment doser la responsabilité quand la liberté est imparfaite ?} C'est l'objet des sections suivantes.

\begin{popfigure}
\begin{tikzpicture}[
    mindmap,
    concept color=popblue,
    level 1/.style={sibling angle=90, level distance=4.5cm},
    level 2/.style={sibling angle=45, level distance=3.5cm},
    every node/.style={concept, font=\small},
    grow cyclic
]
\node[concept, scale=1.2, align=center]{LIBERTÉ \\ \& \\ RESPONSABILITÉ}
    child[concept color=popgreen] { node[concept] {LIBERTÉ}
        child { node[concept, scale=0.7, align=center]{Liberté de faire \\ (absence d'obstacles)} }
        child { node[concept, scale=0.7, align=center]{Liberté de vouloir \\ (autonomie, libre arbitre)} }
        child { node[concept, scale=0.7, align=center]{Choix entre possibles \\ (pouvoir faire autrement)} }
    }
    child[concept color=poporange] { node[concept] {RESPONSABILITÉ}
        child { node[concept, scale=0.7, align=center]{Imputabilité \\ (auteur de l'acte)} }
        child { node[concept, scale=0.7, align=center]{Obligation de répondre \\ (réparer, assumer)} }
        child { node[concept, scale=0.7, align=center]{Sanction morale/juridique \\ (louange ou blâme)} }
    }
    child[concept color=poppurple] { node[concept] {LIAISON}
        child { node[concept, scale=0.7, align=center]{Condition nécessaire : \\ pas de liberté = pas de resp.} }
        child { node[concept, scale=0.7, align=center]{Condition suffisante ? \\ (cas limites, contrainte)} }
        child { node[concept, scale=0.7, align=center]{Degrés de responsabilité \\ selon degrés de liberté} }
    }
    child[concept color=poppink] { node[concept] {ENJEUX CONCRETS}
        child { node[concept, scale=0.7, align=center]{Droit pénal \\ (irresponsabilité, atténuation)} }
        child { node[concept, scale=0.7, align=center]{Vie quotidienne \\ (éducation, travail, famille)} }
        child { node[concept, scale=0.7, align=center]{Citoyenneté \\ (vote, engagement, éthique)} }
    };
\end{tikzpicture}
\legende{Carte mentale : structure conceptuelle de l'unité. La liberté (verte) et la responsabilité (orange) sont liées (violet) par un rapport de conditionnalité nécessaire et de proportionnalité, avec des enjeux concrets (rose) au Cameroun.}
\end{popfigure}

\begin{methode}[Comment définir la liberté en dissertation (Bac/Probatoire)]
\begin{enumerate}
    \item \textbf{Étymologie rapide :} \emph{Libertas} (non-esclavage) $\rightarrow$ \emph{Respondere} (répondre de ses actes).
    \item \textbf{Distinction opératoire :} Liberté \emph{négative} (pouvoir faire, absence de contrainte externe) \textbf{vs} Liberté \emph{positive} (pouvoir vouloir, autonomie, libre arbitre).
    \item \textbf{Définition synthétique :} \og La liberté est le pouvoir de s'autodéterminer par un choix délibéré entre plusieurs possibles, en l'absence de contrainte extérieure insurmontable. \fg{}
    \item \textbf{Lien avec la responsabilité :} Poser d'emblée : \og C'est parce que l'homme est libre qu'il est responsable \fg{} (thèse classique), tout en annonçant la discussion : \og Mais la liberté est-elle totale ? La responsabilité est-elle proportionnelle ? \fg{}
\end{enumerate}
\end{methode}

\coursec{Les formes de la liberté}

\courssub{Le déterminisme : la liberté est-elle une illusion ?}

Le \cle{déterminisme} affirme que tout événement, y compris les actions humaines, est l'effet nécessaire de causes antérieures. Si tout est causé, le "choix" n'est qu'une apparence.

\begin{propriete}[Les trois visages du déterminisme]
\begin{enumerate}
    \item \textbf{Déterminisme physique / naturel :} Les lois de la nature (physique, chimie, biologie, neurosciences) régissent nos cerveaux et nos corps. \emph{Exemple :} Mes gènes, mon taux de dopamine, mes connexions neuronales "décident" avant ma conscience (expériences de Libet).
    \item \textbf{Déterminisme psychologique / psychanalytique :} L'inconscient (pulsions, refoulements, complexes) guide nos choix à notre insu. \emph{Freud :} \og Le moi n'est pas maître en sa propre maison. \fg{}
    \item \textbf{Déterminisme social / sociologique :} Le milieu, la classe, l'éducation, la culture, la langue structurent nos désirs et nos possibilités. \emph{Durkheim, Marx, Bourdieu :} L'habitus dispose, l'individu ne fait qu'exécuter.
\end{enumerate}
\end{propriete}

\begin{attention}[Déterminisme $\neq$ Fatalisme]
\begin{itemize}
    \item \textbf{Déterminisme :} Thèse scientifique/philosophique : \og Tout a une cause. \fg{} Compatible avec l'action (je suis cause de mes actes).
    \item \textbf{Fatalisme :} Attitude résignée : \og Quoi que je fasse, le sort en est jeté. \fg{} Conduit à l'inaction.
\end{itemize}
\emph{Au Bac : ne pas écrire "Le déterminisme prouve qu'on ne peut rien faire." Faux. Le déterminisme dit que nos actes ont des causes, dont nos propres délibérations.}
\end{attention}

\courssub{Le libertarisme (indéterminisme) : la liberté comme rupture causale}

Le \cle{libertarisme} (ou indéterminisme métaphysique) soutient que l'agent est une \emph{causa sui} (cause de soi), capable d'initier une chaîne causale nouvelle sans être déterminé par le passé.

\begin{definition}[Libertarisme]
Thèse selon laquelle \textbf{la volonté humaine peut s'affranchir du déterminisme universel}. L'agent est une source première d'action (\emph{agent-causation}). \emph{Descartes} (le libre arbitre est infini, image de Dieu), \emph{Kant} (la liberté comme postulat de la raison pratique, monde nouménal), \emph{Sartre} (l'homme est condamné à être libre, l'existence précède l'essence).
\end{definition}

\courssub{Le compatibilisme (déterminisme mou) : la liberté dans la causalité}

Le \cle{compatibilisme} (Hume, Stoïciens, Spinoza, auteurs contemporains comme Frankfurt ou Dennett) réconcilie les deux : \textbf{la liberté n'est pas l'absence de cause, mais l'absence de contrainte}.

\begin{propriete}[Thèse compatibiliste]
\begin{itemize}
    \item Je suis libre \textbf{si} mes actions émanent de \emph{mes propres désirs, raisons, caractère}, même si ceux-ci ont des causes antérieures.
    \item \textbf{Contrainte} = cause externe qui me force contre ma volonté (pistolet, chaînes, loi injuste).
    \item \textbf{Détermination interne} = mes motifs, ma raison, mon caractère. Agir selon sa nature rationnelle, c'est être libre (\emph{Spinoza} : \og La liberté est la nécessité comprise \fg{}).
\end{itemize}
\end{propriete}

\begin{exemplebox}[Compatibilisme au quotidien : le fonctionnaire de Yaoundé]
Un agent de la fonction publique traite un dossier. Il \emph{peut} le traiter vite (son devoir, son désir de bien faire) ou le bloquer pour réclamer un "motif" (corruption).
\begin{itemize}
    \item \textbf{Déterministe dur :} Son salaire bas, la culture du "motif", son éducation l'obligent à demander l'argent. Pas de liberté.
    \item \textbf{Libertarien :} Il crée ex nihilo son choix. Liberté totale.
    \item \textbf{Compatibiliste :} Il \emph{veut} l'argent (désir déterminé par son contexte), mais \emph{personne ne le menace de mort} s'il refuse. Il agit selon son désir propre $\rightarrow$ Il est libre \emph{et} responsable. C'est la position du droit camerounais (Code pénal, art. 72 : la contrainte doit être \emph{irrésistible} pour exonérer).
\end{itemize}
\end{exemplebox}

\begin{center}
\begin{tabularx}{\linewidth}{|l|X|X|X|}\hline
\rowcolor{popblueL}
\textbf{Thèse} & \textbf{Liberté} & \textbf{Responsabilité} & \textbf{Auteurs clés} \\ \hline
\textbf{Déterminisme dur} & Illusion (épiphénomène) & Impossible (sauf utilitaire : protection société) & Spinoza (partiel), Holbach, Freud, Marx, Neurosciences \\ \hline
\textbf{Libertarisme} & Absolue (causa sui) & Pleine, entière, métaphysique & Descartes, Kant, Sartre, Leibniz \\ \hline
\textbf{Compatibilisme} & Réelle (absence contrainte externe) & Proportionnelle au degré d'autonomie/raison & Hume, Stoïciens, Hobbes, Droit positif (Cameroun) \\ \hline
\end{tabularx}
\end{center}

\coursec{La responsabilité : définition et formes}

\courssub{Définition générale : répondre de ses actes}

\begin{definition}[Responsabilité]
La \cle{responsabilité} est l'\textbf{obligation pour un agent moral ou juridique de répondre des conséquences de ses actes}, d'en assumer la charge (réparation, sanction, louange) devant une instance (conscience, tribunal, opinion publique). Elle suppose l'\emph{imputabilité} : la capacité d'être regardé comme l'auteur propre de l'acte.
\end{definition}

\courssub{Responsabilité morale vs responsabilité juridique}

\begin{center}
\begin{tabularx}{\linewidth}{|l|X|X|}\hline
\rowcolor{popblueL}
\textbf{Critère} & \textbf{Responsabilité Morale} & \textbf{Responsabilité Juridique} \\ \hline
\textbf{Fondement} & Conscience, devoir, dignité humaine & Loi, Code pénal, Code civil, contrat \\ \hline
\textbf{Juge} & Soi-même (conscience), opinion publique, Dieu & Tribunaux (pénal, civil, administratif) \\ \hline
\textbf{Sanction} & Blâme, remords, honte, estime de soi (interne) & Peine (prison, amende), dommages-intérêts (externe) \\ \hline
\textbf{Condition} & Liberté de vouloir (intention, négligence) & Imputabilité juridique (discernement, contrainte) \\ \hline
\textbf{Exemple Cameroun} & Refuser un pot-de-vin par honneur (même si personne ne voit) & Être condamné pour corruption (Loi 2006/011, Code pénal art. 134) \\ \hline
\end{tabularx}
\end{center}

\courssub{Les formes de la responsabilité juridique au Cameroun}

\begin{propriete}[Typologie juridique (Droit camerounais : Code pénal Loi 2016/007, Code Civil)]
\begin{enumerate}
    \item \textbf{Responsabilité pénale (personnelle, subjective) :} Sanctionne une infraction (crime, délit, contravention). Exige \textbf{l'élément moral} (intention ou imprudence) et \textbf{l'imputabilité} (discernement, absence de contrainte irrésistible).
        \begin{itemize}
            \item \emph{Irresponsabilité pénale (Art. 71 Code pénal) :} Trouble psychique ou neuropsychique ayant aboli le discernement ou le contrôle des actes (aliénation mentale).
            \item \emph{Atténuation / Exonération (Art. 72) :} Contrainte physique ou morale irrésistible (pistolet sur la tempe, enlèvement).
            \item \emph{Minorité :} Moins de 10 ans = irresponsable. 10-18 ans = responsabilité atténuée (mesures de protection/éducation).
        \end{itemize}
    \item \textbf{Responsabilité civile (réparatrice) :} Obligation de réparer le dommage causé à autrui.
        \begin{itemize}
            \item \emph{Pour faute (Art. 1382-1383 Code Civil) :} Faute (liberté) + Dommage + Lien causal. \textbf{Liberté requise.}
            \item \emph{Sans faute (Responsabilité objective) :} \textbf{Garde de la chose} (Art. 1384 al. 1 : \og On est responsable du dommage causé par les choses que l'on a sous sa garde \fg{}), \textbf{Responsabilité du fait d'autrui} (Art. 1384 al. 4-5 : parents/enfants, employeurs/employés), \textbf{Risque} (activités dangereuses). \textbf{Liberté NON requise}, seulement lien causal et garde/bénéfice.
        \end{itemize}
    \item \textbf{Responsabilité collective / solidaire :} Responsabilité de l'État (service public), des personnes morales (entreprises), des groupements.
\end{enumerate}
\end{propriete}

\courssub{Responsabilité objective : le paradoxe apparent}

\begin{exemplebox}[Responsabilité sans liberté ? Le cas du "bâtiment qui s'effondre" à Douala]
Un mur de clôture mal entretenu s'effondre sur la voie publique (Rond-point Deïdo) et blesse un passant. Le propriétaire est en voyage à l'étranger, il n'a \emph{pas choisi} l'effondrement, il n'avait \emph{aucune liberté} sur l'événement immédiat.
\textbf{Pourtant}, le Code civil (Art. 1384, 1386) le tient pour responsable : il a la \emph{garde} de la chose, il en tire profit (sécurité de sa concession), il doit en assumer les risques.
\textbf{Analyse :} C'est une responsabilité \emph{réparatrice} (civile), non \emph{punitif} (pénale). Le propriétaire ne va pas en prison, il indemnise. La liberté n'est pas la source, mais la \emph{garde} et le \emph{bénéfice du risque}.
\end{exemplebox}

\coursec{La liberté, fondement de la responsabilité}

\courssub{L'argument kantien : "Devoir implique Pouvoir"}

\begin{savaistu}[Kant, \emph{Fondements de la métaphysique des mœurs} (1785)]
Pour Kant, la loi morale (\og Agis uniquement d'après la maxime qui fait que tu puisses vouloir en même temps qu'elle devienne loi universelle \fg{}) s'adresse à une volonté \emph{libre}. \textbf{Si je n'étais pas libre, je ne pourrais pas "devoir".} \og Devoir \fg{} implique \og Pouvoir \fg{}. La liberté est le \emph{ratio essendi} (raison d'être) de la responsabilité morale. Sans liberté, l'homme est un automate, un rouage de la nature, non une fin en soi.
\end{savaistu}

\courssub{L'imputabilité : le lien technique entre liberté et responsabilité}

L'\cle{imputabilité} est la capacité juridique et morale d'être l'auteur de ses actes. Elle gradue la responsabilité selon le degré de liberté effective.

\begin{propriete}[Les degrés d'imputabilité (échelle de la responsabilité)]
\begin{enumerate}
    \item \textbf{Imputabilité pleine (Liberté pleine) :} Adulte sain d'esprit, lucide, non contraint. $\rightarrow$ Responsabilité pénale maximale + civile.
    \item \textbf{Imputabilité altérée (Liberté diminuée) :} Troubles psychiques partiels (Art. 71 al. 2 Code pénal : \og trouble ayant altéré le discernement \fg{}), ivresse non volontaire, contrainte morale forte (menace grave), passion soudaine (crime passionnel). $\rightarrow$ Peine réduite, responsabilité civile maintenue.
    \item \textbf{Imputabilité abolie (Liberté nulle) :} Trouble abolissant le discernement (Art. 71 al. 1), contrainte physique irrésistible (Art. 72), enfant < 10 ans. $\rightarrow$ Irresponsabilité pénale. Responsabilité civile possible (garde, risque) mais souvent transférée aux ayants droit (parents, tuteurs).
\end{enumerate}
\end{propriete}

\courssub{Ricœur : la responsabilité comme attestation d'identité}

\begin{savaistu}[Paul Ricœur, \emph{Soi-même comme un autre} (1990)]
Ricœur dépasse l'opposition liberté/déterminisme. La responsabilité n'est pas d'abord un problème métaphysique (suis-je cause de moi ?), mais \textbf{éthique et narrative}. \textbf{Être responsable, c'est "répondre" de sa vie comme d'une histoire cohérente.} Je suis responsable parce que je me \emph{reconnais} comme l'auteur de mes actes (identité narrative) et que je \emph{réponds} devant autrui (attestation). La liberté n'est pas le point de départ, mais le \emph{corrélat} de la responsabilité : je deviens libre \emph{en} devenant responsable.
\end{savaistu}

\coursec{Applications concrètes au quotidien camerounais}

\courssub{La corruption : liberté, contrainte et responsabilité partagée}

\begin{exemplebox}[Le "motif" au guichet administratif]
Un usager paie 5 000 FCFA à un agent pour accélérer un passeport.
\begin{itemize}
    \item \textbf{Agent :} \og Je suis libre, je choisis de prendre l'argent. \fg{} $\rightarrow$ Responsable pénalement (corruption passive, Loi 2006/011).
    \item \textbf{Usager :} \og Je n'étais pas libre, il me forçait la main (contrainte), j'avais l'avion. \fg{} $\rightarrow$ Responsable aussi (corruption active), car la contrainte n'était pas \emph{irrésistible} (il pouvait porter plainte, refuser, saisir le CONAC). La jurisprudence camerounaise sanctionne les deux.
    \item \textbf{Analyse :} La \emph{pauvreté} ou l'\emph{urgence} expliquent le choix (déterminisme social), mais ne le \emph{justifient} pas juridiquement (compatibilisme légal).
\end{itemize}
\end{exemplebox}

\courssub{L'engagement citoyen : le vote comme exercice de liberté responsable}

\begin{exemplebox}[Élections locales ou présidentielles]
\begin{itemize}
    \item \textbf{Liberté négative :} Nul ne m'empêche d'aller voter (pas de barrière physique, carte d'électeur valide).
    \item \textbf{Liberté positive :} Me suis-je informé sur les programmes ? Vote-je par conviction (autonomie) ou par clientélisme (t-shirt, 2 000 FCFA, pression chef de village) ?
    \item \textbf{Responsabilité :} Si j'élis un incompétent par clientélisme, je suis \emph{responsable} de la mauvaise gestion qui suit. \og Peuple souverain \fg{} = peuple \emph{responsable}.
\end{itemize}
\end{exemplebox}

\courssub{L'éthique professionnelle : l'ingénieur technicien face aux normes}

\begin{exemplebox}[Chantier BTP à Kribi : béton dosé à 250 kg/m³ au lieu de 350]
Le chef de chantier ordonne : \og On dose léger, le ciment est cher, le patron ne verra rien. \fg{}
\begin{itemize}
    \item \textbf{Technicien (vous) :} Liberté de faire (vous pouvez refuser, alerter le bureau de contrôle). Liberté de vouloir (votre conscience pro, votre diplôme, la sécurité des usagers).
    \item \textbf{Responsabilité :} Si l'immeuble s'effondre dans 5 ans, le Code pénal (homicide involontaire, mise en danger d'autrui) et le Code civil (responsabilité décennale Art. 1792) vous rattraperont. \og J'obéissais aux ordres \fg{} n'est pas une cause d'irresponsabilité (sauf contrainte irrésistible = menace de mort immédiate).
\end{itemize}
\end{exemplebox}

\courssub{Responsabilité environnementale : le cas des plastiques à Yaoundé}

\begin{exemplebox}[Interdiction des sachets plastiques non biodégradables (Arrêté 2012/0041/MINEPDED)]
\begin{itemize}
    \item \textbf{Producteur/Importateur :} Responsabilité civile (réparation pollution) + pénale (amendes). Liberté d'entreprendre $\neq$ liberté de polluer.
    \item \textbf{Commerçant (Marché Mokolo) :} Responsable s'il vend les sachets interdits. Liberté de choisir ses fournisseurs.
    \item \textbf{Citoyen :} Responsabilité morale/civique : refuser le sachet, utiliser son panier. \emph{Liberté de consommer responsable.}
\end{itemize}
\end{exemplebox}

\coursec{Méthode de la dissertation et synthèse}

\begin{methode}[Structure type dissertation : « La liberté est-elle le fondement de la responsabilité ? »]
\begin{enumerate}
    \item \textbf{Introduction :}
        \begin{itemize}
            \item \textbf{Accroche :} Exemple concret (moto-taxi, vote, corruption, chantier).
            \item \textbf{Définitions :} Liberté (pouvoir s'autodéterminer, libre arbitre) ; Responsabilité (devoir répondre, imputabilité). Distinction morale/juridique.
            \item \textbf{Problématique :} La liberté est-elle la \emph{seule} source (condition nécessaire et suffisante) de la responsabilité, ou existe-t-il des responsabilités sans liberté (solidarité, risque, garde) ?
            \item \textbf{Annonce du plan :} Thèse (liberté = fondement nécessaire) / Antithèse (responsabilités sans liberté / liberté altérée) / Synthèse (responsabilité analogique : morale vs juridique, proportionnalité).
        \end{itemize}
    \item \textbf{Développement (3 parties, 2-3 arguments chacune) :}
        \begin{itemize}
            \item \textbf{I. Thèse : La liberté, condition \emph{sine qua non} de la responsabilité morale et pénale.}
                \begin{itemize}
                    \item A. Kant : Devoir $\Rightarrow$ Pouvoir. L'agent moral est une cause libre.
                    \item B. Droit pénal camerounais (Art. 71, 72) : Irresponsabilité si discernement aboli ou contrainte irrésistible. Le mineur, l'aliéné ne sont pas coupables \emph{parce que} non libres.
                    \item C. Le libre arbitre comme "pouvoir faire autrement" (Descartes) : fondement du blâme et de la louange.
                \end{itemize}
            \item \textbf{II. Antithèse : Des responsabilités sans liberté (ou avec liberté altérée).}
                \begin{itemize}
                    \item A. Responsabilité civile objective (Code Civil Art. 1384 ss) : Garde de la chose, fait d'autrui, risque. Pas de faute (liberté) requise, seulement causalité + garde.
                    \item B. Responsabilité solidaire / vicariée : Employeur/employé, parent/enfant, État/service public. L'un répond pour l'autre sans avoir choisi l'acte.
                    \item C. Responsabilité collective / historique / écologique : Dette écologique, devoir de mémoire. Je réponds d'actes que je n'ai pas commis (solidarité intergénérationnelle).
                \end{itemize}
            \item \textbf{III. Synthèse : La responsabilité, notion analogique à degrés.}
                \begin{itemize}
                    \item A. Distinction nette : Responsabilité \emph{d'imputation} (morale/pénale) $\leftarrow$ exige liberté. Responsabilité \emph{d'attribution/réparation} (civile/objective) $\leftarrow$ repose sur solidarité, risque, garde.
                    \item B. Proportionnalité (Compatibilisme juridique) : Plus la liberté est grande (intention), plus la responsabilité est lourde (crime vs délit vs contravention ; circonstance aggravante/atténuante).
                    \item C. Ricœur : La responsabilité comme \emph{compétence} éthique (répondre de soi devant autrui). La liberté se \emph{construit} dans l'exercice de la responsabilité.
                \end{itemize}
        \end{itemize}
    \item \textbf{Conclusion :}
        \begin{itemize}
            \item \textbf{Réponse synthétique :} \og La liberté est le fondement \emph{nécessaire} de la responsabilité \emph{pleine} (morale, pénale), mais non \emph{suffisant} pour épuiser le champ de la responsabilité \emph{juridique} (civile objective, solidaire) qui répond à d'autres impératifs (réparation, solidarité). \fg{}
            \item \textbf{Ouverture :} L'intelligence artificielle et la responsabilité algorithmique : qui répond quand la "machine apprend" et décide sans libre arbitre ? (Projet de règlement UE sur l'IA, responsabilité du fabricant/déployeur).
        \end{itemize}
\end{enumerate}
\end{methode}

\begin{exoresolu}[Sujet type Probatoire : « Y a-t-il une responsabilité sans liberté ? »]
\textbf{Énoncé :} \og Y a-t-il une responsabilité sans liberté ? \fg{}

\textbf{Analyse du sujet :}
\begin{itemize}
    \item \textbf{Termes clés :} Responsabilité (devoir répondre, imputabilité, réparation), Liberté (pouvoir choisir, autonomie, absence de contrainte), \emph{Sans} (absence, privation, indépendance).
    \item \textbf{Type de sujet :} Question fermée (Oui/Non) $\rightarrow$ Exige une réponse nuancée (Thèse/Antithèse/Synthèse). \textbf{Piège :} Répondre simplement "Non" (thèse classique) ou "Oui" (droit civil) sans articuler les deux sens.
    \item \textbf{Problématique :} La liberté est-elle le fondement \emph{unique} et \emph{indépassable} de toute responsabilité, ou la responsabilité peut-elle s'étendre à des cas où la liberté fait défaut (solidarité, risque, causalité objective) ?
\end{itemize}

\textbf{Plan détaillé :}

\textbf{Introduction}
\begin{itemize}
    \item Accroche : Exemple du conducteur de moto-taxi (Yaoundé) / du propriétaire du mur (Douala).
    \item Définitions : Liberté (pouvoir s'autodéterminer, libre arbitre) ; Responsabilité (obligation de répondre/réparer, imputabilité).
    \item Problématique : La liberté est-elle la condition \emph{sine qua non} de la responsabilité, ou la responsabilité peut-elle exister sur d'autres fondements (solidarité, risque, garde) ?
    \item Annonce du plan.
\end{itemize}

\textbf{I. Thèse : La liberté est le fondement nécessaire de la responsabilité morale et pénale (classique).}
\begin{itemize}
    \item A. Responsabilité morale : Kant, \emph{Fondements de la métaphysique des mœurs}. \og Devoir \fg{} implique \og Pouvoir \fg{}. Nul n'est tenu à l'impossible. Si je ne suis pas libre, je suis un automate, non un agent moral. Le blâme/la louange supposent le libre arbitre (Descartes, \emph{Méditations IV} : \og Je peux faire autrement \fg{}).
    \item B. Responsabilité pénale : Code pénal camerounais (Loi 2016/007). Art. 71 : Irresponsabilité pour trouble psychique ayant \emph{aboli} le discernement ou le contrôle (liberté abolie). Art. 72 : Exonération pour contrainte \emph{irrésistible} (liberté annihilée). Le mineur (<10 ans) irresponsable car liberté naissante. \textbf{Principle :} \emph{Nullum crimen, nulla poena sine culpa} (pas de crime, pas de peine sans faute/liberté).
    \item C. Le libre arbitre comme condition de l'imputabilité : Je suis responsable \emph{parce que} j'aurais pu faire autrement (pouvoir de contre-factuel).
\end{itemize}

\textbf{II. Antithèse : Des formes de responsabilité juridiques effectives sans liberté (ou avec liberté altérée).}
\begin{itemize}
    \item A. Responsabilité civile objective / sans faute (Code Civil applicable au Cameroun, Art. 1384 ss).
        \begin{itemize}
            \item \emph{Garde de la chose (Art. 1384 al. 1) :} \og On est responsable du dommage causé par les choses que l'on a sous sa garde \fg{}. Ex: Mur qui s'effondre (Douala), animal qui blesse (Art. 1385), voiture prêtée. Pas de liberté requise, seulement \emph{garde} + \emph{dommage} + \emph{lien causal}.
            \item \emph{Fait d'autrui (Art. 1384 al. 4-5) :} Parents responsables des enfants (mineurs), employeurs des employés (fonction de leurs fonctions). L'employeur n'a pas \emph{choisi} la faute, mais il en répond (bénéfice du risque, organisation).
            \item \emph{Responsabilité du fait des bâtiments (Art. 1386) :} Propriétaire responsable du défaut d'entretien, même absent.
        \end{itemize}
    \item B. Responsabilité pour risque / solidarité. Celui qui crée un risque profitable (usine, barrage, transport) doit en assumer les charges (réparation), libre ou non. Responsabilité de l'État (service public, égalité devant les charges publiques).
    \item C. Responsabilité collective / historique. Devoir de mémoire, réparation esclavagisme/colonisation. Je réponds pour des actes que je n'ai pas commis, par appartenance à une communauté politique.
\end{itemize}

\textbf{III. Synthèse : La responsabilité comme notion analogique (plusieurs sens selon le régime).}
\begin{itemize}
    \item A. Distinction cruciale : \textbf{Responsabilité d'imputation (morale/pénale)} vs \textbf{Responsabilité d'attribution/réparation (civile objective/solidaire)}.
        \begin{itemize}
            \item La première touche à la \emph{dignité} de l'agent (blâme, peine, honneur) $\rightarrow$ \textbf{Exige la liberté} (condition nécessaire et proportionnelle).
            \item La seconde touche à la \emph{réparation} du dommage (justice commutative, solidarité) $\rightarrow$ \textbf{Peut se passer de liberté} (garde, risque, lien de droit).
        \end{itemize}
    \item B. Ricœur (\emph{Sois-toi même comme un autre}) : La responsabilité comme \emph{attestation} d'une identité narrative. Je réponds de mes actes \emph{parce que} je me reconnais comme leur auteur (liberté), mais aussi \emph{pour} les autres (solidarité). La liberté n'est pas seulement le point de départ, elle est le \emph{corrélat} de la responsabilité : je deviens libre \emph{en} assumant mes responsabilités.
    \item C. Au Cameroun : Le Code pénal (Loi 2016/007) maintient l'irresponsabilité pour trouble mental (liberté abolie) mais prévoit la responsabilité civile sans faute (garde des choses, animaux, bâtiments, fait d'autrui). La responsabilité \emph{pleine} (sanction pénale, blâme moral) reste attachée à la liberté ; la responsabilité \emph{réparatrice} (civile) peut s'en détacher pour protéger la victime.
\end{itemize}

\textbf{Conclusion}
\begin{itemize}
    \item Réponse : \og Oui, il existe une responsabilité \emph{juridique/réparatrice} sans liberté (solidarité, risque, garde), mais la responsabilité \emph{morale/pénale} (celle qui touche à la dignité de l'agent et à la culpabilité) \emph{exige} la liberté. \fg{}
    \item Ouverture : L'IA, les algorithmes décisionnels (notation sociale, voitures autonomes, diagnostic médical) : qui est responsable quand la machine « choisit » sans liberté ni conscience ? (Responsabilité du fabricant, du déployeur, fonds de garantie ?). Le droit camerounais devra légiférer (protection données, cybercriminalité).
\end{itemize}
\end{exoresolu}

\begin{aretenir}[Synthèse de l'unité : Liberté et Responsabilité]
\begin{itemize}
    \item \textbf{Liberté} = Pouvoir s'autodéterminer (Libre arbitre). Deux étages : \emph{Liberté de faire} (externe) et \emph{Liberté de vouloir} (interne, autonomie).
    \item \textbf{Responsabilité} = Devoir répondre de ses actes. Deux régimes : \emph{Morale/Pénale} (subjective, sanction, exige liberté) et \emph{Civile/Objective} (réparatrice, solidarité/risque, peut dispenser de liberté).
    \item \textbf{Lien :} \textbf{Condition nécessaire} pour la responsabilité d'imputation (culpabilité). \textbf{Condition non suffisante} pour épuiser la responsabilité d'attribution (réparation).
    \item \textbf{Degrés :} Responsabilité proportionnelle au degré de liberté (intention > imprudence > contrainte > irresponsabilité).
    \item \textbf{Enjeu citoyen (Cameroun) :} Lutter contre la corruption (assumer sa liberté), voter en conscience (exercer sa responsabilité), respecter les normes professionnelles (responsabilité technique), protéger l'environnement (responsabilité solidaire).
\end{itemize}
\end{aretenir}

\begin{exercice}[Entraînement Probatoire - Sujets probables]
\begin{enumerate}
    \item \og La responsabilité suppose-t-elle la liberté ? \fg{} (Classique, thèse/antithèse/synthèse).
    \item \og Est-on responsable seulement de ce que l'on veut ? \fg{} (Focus intention/négligence/responsabilité objective).
    \item \og Peut-on être responsable sans être coupable ? \fg{} (Distinction responsabilité civile/pénale, responsabilité sans faute).
    \item \og La contrainte supprime-t-elle la responsabilité ? \fg{} (Degrés de contrainte : physique/psychologique, irrésistible/résistible, Code pénal art. 72).
    \item \og Commenter : "L'homme est condamné à être libre" (Sartre) / "La liberté est la nécessité comprise" (Spinoza). \fg{}
\end{enumerate}
\end{exercice}

\begin{corrige}[Exercice n°1 : « La responsabilité suppose-t-elle la liberté ? »]
\textbf{Plan type :}
\textbf{Intro :} Accroche (ex: moto-taxi / mur qui tombe). Déf : Liberté (pouvoir choisir), Responsabilité (devoir répondre). Problématique : Liberté = condition nécessaire ? unique ? Annonce plan.
\textbf{I. Oui, pour la responsabilité morale et pénale (imputation).} A. Kant : Devoir $\Rightarrow$ Pouvoir. B. Droit pénal (Cameroun Art 71, 72) : Irresponsabilité si pas de discernement/contrainte. C. Libre arbitre = pouvoir faire autrement.
\textbf{II. Non, pour la responsabilité civile objective / solidaire (réparation).} A. Garde de la chose (Art 1384), fait d'autrui, risque. B. Responsabilité employeur/parent/État. C. Solidarité nationale/intergénérationnelle.
\textbf{III. Synthèse : Analogie de la responsabilité.} A. Distinction imputation (liberté requise) / attribution (solidarité/risque). B. Proportionnalité (compatibilisme) : plus de liberté = plus de responsabilité. C. Ricœur : Responsabilité comme compétence éthique (répondre de soi).
\textbf{Conclusion :} Réponse nuancée. Ouverture IA / responsabilité algorithmique.
\end{corrige}

\coursec{Les formes de la liberté}

Dans la section précédente, nous avons problématisé la notion de liberté et proposé une définition générale : être libre, c'est pouvoir agir selon sa volonté, sans contrainte extérieure qui nous empêche de faire ce que nous voulons. Mais cette définition reste abstraite. Dans la vie réelle, la liberté se présente sous des visages très différents : le prisonnier derrière les barreaux, le journaliste qui publie son article, le citoyen qui résiste à une mauvaise tentation, le sage qui reste serein dans l'épreuve. Tous sont concernés par la liberté, mais pas de la même manière. Il faut donc distinguer les \cle{formes de la liberté} : la liberté \cle{physique}, la liberté \cle{politique}, la liberté \cle{morale} et la liberté \cle{intérieure}. Cette distinction est indispensable pour la dissertation : un sujet sur la liberté exige presque toujours de préciser de quelle liberté l'on parle.

\courssub{La liberté physique : l'absence d'obstacles matériels}

Commençons par l'intuition la plus simple. Un oiseau dans une cage ne peut pas voler ; ouvrez la porte, et le voilà libre. De même, un homme enchaîné ne peut pas se déplacer ; brisez ses chaînes, et il recouvre sa liberté de mouvement. Cette première forme de liberté est la plus immédiate, la plus visible : elle concerne le \emph{corps}.

\begin{definition}[La liberté physique]
La \cle{liberté physique} est l'absence d'obstacles matériels qui empêcheraient le corps d'accomplir les mouvements voulus par la volonté. Elle consiste dans le pouvoir effectif d'aller, de venir et d'agir sans entrave corporelle.
\end{definition}

L'exemple paradigmatique est celui du \textbf{prisonnier} : enfermé dans une cellule, il peut \emph{vouloir} sortir, mais les murs et les verrous rendent sa volonté impuissante. Il en va de même de l'\textbf{otage} retenu contre son gré, ou du malade paralysé cloué dans son lit. Dans tous ces cas, l'obstacle est \emph{matériel} : murs, chaînes, barreaux, ou même le corps lui-même lorsqu'il est malade.

Thomas Hobbes (1588-1679) a donné de cette liberté une définition célèbre dans le \emph{Léviathan} (1651) : l'homme libre est « celui qui, dans les choses que ses forces et son intelligence lui permettent d'accomplir, n'est pas empêché de faire ce qu'il a la volonté de faire ». La liberté physique est donc une liberté \emph{négative} : elle se définit par l'absence d'empêchement, non par un pouvoir positif.

\begin{exemplebox}[Situations camerounaises]
Un commerçant de Douala qui circule librement entre le marché Mboppi et son quartier jouit de sa liberté physique. En revanche, un individu placé en garde à vue ou détenu à la prison centrale de New Bell en est privé — privation qui, lorsqu'elle est prononcée par un jugement régulier, est justement la forme que prend la peine : la société punit en retirant la liberté physique, ce qui montre combien celle-ci est précieuse.
\end{exemplebox}

Mais attention : la liberté physique est la forme la plus \emph{superficielle} de la liberté. Un homme peut disposer de tout son corps et se comporter pourtant comme un esclave — de ses passions, de l'opinion, de la peur. C'est pourquoi les philosophes distinguent d'autres formes de liberté, plus profondes.

\courssub{La liberté politique : les droits garantis par l'État de droit}

Passons du corps à la cité. Un citoyen peut être physiquement libre de ses mouvements et pourtant ne pas oser exprimer une opinion, craindre d'être arrêté sans motif, ou se voir interdire de se réunir avec d'autres. La liberté physique ne suffit donc pas : il faut que la société politique \emph{reconnaisse} et \emph{protège} certaines libertés.

\begin{definition}[La liberté politique]
La \cle{liberté politique} est l'ensemble des droits reconnus et protégés par les lois d'un État. Elle comprend les libertés publiques ou fondamentales : liberté d'expression, liberté de la presse, liberté de circulation, liberté de réunion et d'association, liberté de conscience et de religion.
\end{definition}

L'idée essentielle est celle d'\textbf{État de droit} : un État où le pouvoir lui-même est soumis à la loi, et où les libertés des citoyens sont inscrites dans des textes juridiques que nul, pas même le gouvernant, ne peut violer impunément. Montesquieu (1689-1755), dans \emph{De l'esprit des lois} (1748), formule le principe décisif : « La liberté est le droit de faire tout ce que les lois permettent. » La liberté politique n'est donc pas le pouvoir de faire n'importe quoi : c'est le pouvoir de faire ce que la loi n'interdit pas, dans un cadre qui protège chacun.

\begin{exemplebox}[Les libertés fondamentales dans la Constitution camerounaise]
Le préambule de la Constitution de la République du Cameroun du 18 janvier 1996 affirme solennellement l'attachement du peuple camerounais aux libertés fondamentales inscrites dans la Déclaration universelle des droits de l'homme et la Charte africaine des droits de l'homme et des peuples. Il proclame notamment que « la liberté de religion, de culte et la liberté de conscience sont garanties », que « la liberté de communication, d'expression, de la presse, de réunion et d'association » sont reconnues, et que nul ne peut être poursuivi, arrêté ou détenu arbitrairement. Ces principes ont valeur constitutionnelle : ils s'imposent à tous.
\end{exemplebox}

\begin{exemplebox}[La liberté de la presse au Cameroun]
La liberté de la presse au Cameroun est encadrée par la loi n° 90/052 du 19 décembre 1990 portant régime de la communication sociale, texte fondateur de la libéralisation des médias au Cameroun. Grâce à elle, des journaux privés, des radios et des télévisions peuvent publier et diffuser librement. Mais cette liberté a des limites légales : la diffamation, l'injure, l'appel à la haine ou la divulgation de fausses nouvelles peuvent être poursuivis. La loi n'écrase donc pas la liberté de la presse : elle la \emph{garantit} en la délimitant.
\end{exemplebox}

\begin{attention}[Erreur fréquente : confondre liberté et absence de loi]
Beaucoup d'élèves écrivent que « la liberté politique supprime toute loi » ou que « être libre, c'est faire ce qu'on veut ». C'est exactement le contraire ! Sans loi, c'est la loi du plus fort qui règne : le plus puissant écrase la liberté des plus faibles. La loi \emph{protège} la liberté : elle interdit à chacun de nuire à autrui, et garantit ainsi à tous un espace de liberté. Comme le dit Montesquieu, la liberté consiste à pouvoir faire ce qu'on doit vouloir et à n'être point contraint de faire ce qu'on ne doit pas vouloir. Dans une dissertation, écrire « la liberté, c'est l'absence de lois » est une faute grave.
\end{attention}

\courssub{La liberté morale : choisir en conscience entre le bien et le mal}

Il reste une forme de liberté plus intime encore. Un homme peut être physiquement libre et vivre dans un État de droit, et pourtant se trouver devant un choix : dois-je tricher à l'examen ? Dois-je détourner l'argent qui m'est confié ? Dois-je dénoncer une injustice au risque de me faire des ennemis ? Ici, aucun obstacle matériel, aucune loi ne décide à ma place : c'est ma \emph{conscience} qui doit trancher.

\begin{definition}[La liberté morale]
La \cle{liberté morale} est le pouvoir de se déterminer soi-même selon sa conscience et sa raison, c'est-à-dire la capacité de choisir librement entre le bien et le mal, et d'assumer ce choix.
\end{definition}

Cette liberté suppose le \cle{libre arbitre} : le pouvoir de choisir sans être totalement déterminé par des causes extérieures. Lorsque je résiste à la tentation de copier sur mon voisin au Baccalauréat technique, je fais usage de ma liberté morale : je pourrais tricher (rien ne m'en empêche physiquement), mais je choisis l'honnêteté parce que ma raison et ma conscience me disent que c'est le bien. Emmanuel Kant (1724-1804) voit dans cette capacité le propre de l'homme : être libre, c'est pouvoir obéir à la loi morale que l'on se donne à soi-même (l'\emph{autonomie}), au lieu de suivre aveuglément ses désirs (l'\emph{hétéronomie}).

La liberté morale est le cœur de notre leçon, car c'est elle qui fonde la \emph{responsabilité} : on ne peut reprocher à quelqu'un ses actes que s'il aurait pu faire autrement. Nous y reviendrons dans la section 4.

\courssub{La liberté intérieure selon Épicure et les stoïciens}

Les philosophes de l'Antiquité grecque et romaine ont poussé la réflexion jusqu'au bout : peut-on être libre \emph{même enchaîné} ? Leur réponse est oui, à condition de maîtriser ses désirs.

Pour \textbf{Épicure} (341-270 av. J.-C.), le bonheur et la liberté consistent à ne désirer que ce qui est naturel et nécessaire (manger, boire, l'amitié), et à renoncer aux désirs vains et illimités (richesses, gloire, plaisirs excessifs) qui nous rendent dépendants et malheureux. Celui qui désire peu dépend de peu : il est libre.

Pour les \textbf{stoïciens} (Épictète, Sénèque, Marc Aurèle), il faut distinguer ce qui dépend de nous (nos jugements, nos volontés, nos désirs) et ce qui n'en dépend pas (le corps, la réputation, la fortune, la mort). La liberté véritable consiste à vouloir ce qui dépend de nous et à accepter sereinement le reste. Ainsi, même en prison, même esclave, l'homme reste libre \emph{intérieurement}, car nul ne peut l'empêcher de penser et de juger comme il veut.

\begin{savaistu}[Épictète, l'esclave devenu philosophe de la liberté]
Épictète (vers 50-vers 130 ap. J.-C.) naquit esclave à Hiérapolis, en Asie Mineure. Boiteux, propriété d'un maître à Rome, il fut affranchi et devint l'un des plus grands maîtres du stoïcisme. Lui qui avait connu les chaînes enseignait : « Ce qui trouble les hommes, ce ne sont pas les choses, mais les opinions qu'ils en ont. » La vraie liberté, disait-il, est intérieure : le tyran peut enchaîner mon corps, mais non ma volonté. Voilà un exemple éclatant que la liberté ne se réduit pas à la liberté physique.
\end{savaistu}

\courssub{Les limites de la liberté : la liberté d'autrui}

Si chacun est libre, alors ma liberté rencontre nécessairement celle des autres. D'où le principe célèbre, popularisé par John Stuart Mill (1806-1873) dans \emph{De la liberté} (1859) : « ma liberté s'arrête où commence celle des autres ». Ma liberté d'expression ne m'autorise pas à diffamer mon voisin ; ma liberté de circulation ne m'autorise pas à m'installer chez autrui ; ma liberté de faire du bruit s'arrête là où commence le droit de mes voisins au repos.

Ce principe n'est pas une négation de la liberté, mais sa condition : sans limites réciproques, la liberté de chacun serait détruite par celle de tous. C'est le fondement du droit et de la vie en société.

\begin{popfigure}
\begin{center}
\begin{tikzpicture}[scale=1]
\draw[very thick, popblue, fill=popblueL] (-1.6,0) circle (1.6);
\draw[very thick, poporange, fill=poporangeL, fill opacity=0.7] (1.6,0) circle (1.6);
\node[popdark, font=\bfseries] at (-1.6,0.5) {Ma liberté};
\node[popdark, font=\small] at (-1.6,-0.2) {(expression,};
\node[popdark, font=\small] at (-1.6,-0.65) {circulation, bruit...)};
\node[popdark, font=\bfseries] at (1.6,0.5) {La liberté d'autrui};
\node[popdark, font=\small] at (1.6,-0.2) {(réputation,};
\node[popdark, font=\small] at (1.6,-0.65) {propriété, repos...)};
\draw[very thick, poppink] (0,-1.6) -- (0,1.6);
\node[below, poppink, font=\bfseries\small] at (0,-1.7) {La loi : frontière commune};
\end{tikzpicture}
\end{center}
\legende{Schéma du principe « ma liberté s'arrête où commence celle d'autrui » : les deux sphères de liberté se touchent sans se recouvrir ; la loi (trait vertical) marque la frontière qui protège chacun.}
\end{popfigure}

\courssub{Tableau récapitulatif des formes de la liberté}

\begin{center}
\begin{tabular}{|p{3cm}|p{5.2cm}|p{5.2cm}|}
\hline
\rowcolor{popblueL} \textbf{Forme de liberté} & \textbf{Définition} & \textbf{Exemple camerounais} \\
\hline
Liberté physique & Absence d'obstacles matériels empêchant le corps d'agir & Circuler librement entre Yaoundé et Douala ; le détenu en est privé \\
\hline
Liberté politique & Droits reconnus et protégés par les lois d'un État de droit & Liberté de la presse garantie par la loi de 1990 et le préambule de la Constitution de 1996 \\
\hline
Liberté morale & Pouvoir de se déterminer selon sa conscience et sa raison & Refuser de tricher à l'examen ou de détourner des fonds publics \\
\hline
Liberté intérieure & Maîtrise de ses désirs et de ses jugements (Épicure, stoïciens) & Rester digne et serein dans l'épreuve, comme Épictète esclave \\
\hline
\end{tabular}
\end{center}

\courssub{Complément Bac : la question du déterminisme}

\begin{attention}[Approfondissement pour le Bac — hors cœur du programme, mais valorisé en dissertation]
La liberté, surtout morale, semble contestée par le \cle{déterminisme} : la thèse selon laquelle tout événement, y compris nos choix, est l'effet nécessaire de causes antérieures. Trois types de causes semblent nier notre liberté : les causes \textbf{naturelles} (l'hérédité, le tempérament, les lois de la nature), les causes \textbf{sociales} (le milieu familial, l'éducation, les coutumes — un enfant né dans un quartier pauvre n'a-t-il pas sa conduite déjà orientée ?) et les causes \textbf{psychologiques} (les habitudes, les passions, l'inconscient). Si tout est déterminé, le libre arbitre serait une illusion, et la responsabilité s'effondrerait. Les philosophes répondent que connaître les causes qui nous influencent ne supprime pas notre capacité de juger et de choisir : le déterminisme explique nos penchants, mais ne décide pas à notre place. Retenez cette objection pour vos dissertations : elle montre que vous savez problématiser.
\end{attention}

\begin{exoresolu}[Identifier les formes de la liberté]
\textbf{Énoncé.} Pour chacune des situations suivantes, indiquez la forme de liberté concernée et justifiez : (a) un journaliste camerounais publie un article critiquant la gestion municipale ; (b) un détenu médite et pardonne intérieurement à ceux qui l'ont accusé injustement ; (c) un élève refuse de participer à une fraude organisée à l'examen ; (d) un otage est libéré par les forces de sécurité.
\tcblower
\textbf{Solution détaillée.}
\begin{enumerate}
\item[(a)] Il s'agit de la \textbf{liberté politique}, et plus précisément de la liberté de la presse : un droit reconnu et protégé par les lois de l'État (préambule de la Constitution de 1996, loi de 1990 sur la communication sociale).
\item[(b)] Il s'agit de la \textbf{liberté intérieure} au sens stoïcien : le détenu est privé de liberté physique, mais reste maître de ses jugements et de ses sentiments ; nul ne peut l'empêcher de pardonner.
\item[(c)] Il s'agit de la \textbf{liberté morale} : l'élève choisit en conscience entre le bien (l'honnêteté) et le mal (la fraude), et se détermine lui-même selon sa raison.
\item[(d)] Il s'agit de la \textbf{liberté physique} : la libération supprime l'obstacle matériel (les chaînes, l'enfermement) qui empêchait l'otage de se mouvoir.
\end{enumerate}
\textbf{Conseil méthodologique :} ce type d'exercice entraîne le savoir-faire « appliquer les notions à des situations concrètes ». Au Bac, procédez toujours en deux temps : nommer la notion, puis justifier en rappelant sa définition.
\end{exoresolu}

\begin{aretenir}[L'essentiel de la section]
La liberté se décline en quatre formes : la \cle{liberté physique} (absence d'obstacles matériels : le prisonnier, l'otage), la \cle{liberté politique} (droits garantis par l'État de droit, affirmés dans le préambule de la Constitution camerounaise de 1996), la \cle{liberté morale} (choix en conscience entre le bien et le mal) et la \cle{liberté intérieure} (maîtrise de ses désirs selon Épicure et les stoïciens). La liberté de chacun est limitée par celle d'autrui : la loi ne supprime pas la liberté, elle la protège. Enfin, le déterminisme (causes naturelles, sociales, psychologiques) constitue l'objection majeure contre la liberté — objection à mobiliser pour problématiser une dissertation.
\end{aretenir}

Ayant ainsi distingué les formes de la liberté, nous pouvons aborder la notion symétrique qui lui est indissolublement liée : la responsabilité. Car s'il existe plusieurs manières d'être libre, il existe aussi plusieurs manières de répondre de ses actes. C'est l'objet de la section suivante.

\coursec{La responsabilité : définition et formes}

Dans la section précédente, nous avons montré que la liberté se décline en plusieurs formes : liberté physique (faire ce qu'on veut), liberté morale (choisir le bien), liberté politique (participer à la vie de la cité). Mais une liberté qui n'engagerait à rien serait un simple caprice. Or, dans la vie quotidienne, chaque fois que quelqu'un agit, on lui demande des comptes : l'élève qui casse la vitre de la salle de classe, le chauffeur qui provoque un accident sur l'axe Yaoundé--Douala, le ministre qui gère mal les deniers publics. Tous doivent \emph{répondre} de leurs actes. C'est précisément ce qu'exprime le concept de \cle{responsabilité}, inséparable de celui de liberté : on ne peut demander des comptes qu'à celui qui a agi librement. Cette section définit la responsabilité, en distingue les formes principales et en analyse les conditions.

\courssub{Définition générale de la responsabilité}

Commençons par une intuition simple. Un motocycliste (« benskineur ») à Douala percute un piéton parce qu'il roulait trop vite. Tout le monde s'accorde à dire qu'il doit « payer » : payer les soins du blessé, éventuellement répondre devant un tribunal. Pourquoi lui ? Parce que c'est \emph{lui} qui a agi, qu'il savait ce qu'il faisait et qu'il aurait pu faire autrement. La responsabilité naît de ce lien entre un auteur, son acte et les conséquences de cet acte.

\begin{definition}[La responsabilité]
La \cle{responsabilité} est l'obligation morale ou juridique de \textbf{répondre de ses actes} et d'en \textbf{assumer les conséquences}. Être responsable, c'est pouvoir être tenu pour l'auteur de ses actes et devoir en réparer les effets (réparation, peine) ou en recevoir le mérite (louange, récompense).
\end{definition}

L'étymologie éclaire cette définition : responsabilité vient du latin \emph{respondere}, « répondre ». L'être responsable est celui qui \emph{répond} : il répond de lui-même devant les autres, il répond à la question « qui a fait cela ? » par « c'est moi ». La responsabilité comporte donc trois éléments :

\begin{itemize}
\item un \textbf{auteur} : celui qui a accompli l'acte ;
\item un \textbf{acte} et ses \textbf{conséquences} (dommage, faute, mais aussi réussite) ;
\item une \textbf{instance} devant laquelle on répond : sa conscience, la victime, la société, la loi.
\end{itemize}

Selon l'instance devant laquelle on répond, on distingue plusieurs formes de responsabilité : morale, civile, pénale, politique et professionnelle.

\courssub{La responsabilité morale : répondre devant sa conscience}

\begin{exemplebox}[Le remords de l'élève]
Un élève de Terminale triche à une évaluation en recopiant le téléphone de son voisin. Personne ne s'en aperçoit : ni le surveillant, ni le correcteur. Il obtient une bonne note. Pourtant, le soir, il ne dort pas tranquille : il éprouve du \emph{remords}. Aucune sanction extérieure ne le menace, mais sa propre conscience le juge.
\end{exemplebox}

Cet exemple montre qu'il existe un tribunal intérieur : la \cle{conscience morale}. La \textbf{responsabilité morale} est l'obligation de répondre de ses actes devant sa propre conscience. Sa sanction n'est pas une amende ni une prison, mais un sentiment pénible : le \textbf{remords} (douleur d'avoir mal agi) ou la \textbf{culpabilité} (sentiment d'être fautif). Inversement, celui qui a bien agi éprouve la satisfaction du devoir accompli.

Pour le philosophe Emmanuel \textbf{Kant}, la conscience morale est la voix du devoir en nous : même si nul ne nous voit, nous savons ce que nous devions faire, et nous nous jugeons nous-mêmes. La responsabilité morale est donc la forme la plus intérieure et la plus universelle de la responsabilité : elle concerne tous nos actes, même les plus secrets.

\courssub{La responsabilité civile : réparer le dommage causé à autrui}

Lorsque la faute cause un dommage à autrui, la conscience ne suffit plus : la loi intervient pour obliger l'auteur à \emph{réparer}.

\begin{definition}[Responsabilité civile]
La \cle{responsabilité civile} est l'obligation, imposée par le droit, de \textbf{réparer le dommage} (matériel, corporel ou moral) causé à autrui par sa faute, sa négligence ou son imprudence. Sa sanction est la \textbf{réparation}, le plus souvent par le versement de dommages et intérêts.
\end{definition}

Le principe est formulé de manière célèbre par l'article 1382 du Code civil (applicable dans la tradition juridique héritée du droit français au Cameroun) : « Tout fait quelconque de l'homme, qui cause à autrui un dommage, oblige celui par la faute duquel il est arrivé à le réparer. »

\begin{exemplebox}[La vitrine brisée]
Lors d'une partie de football dans la cour d'un quartier de Bafoussam, un adolescent envoie le ballon dans la vitrine d'une boutique. Il n'avait pas l'intention de casser, mais il a commis une imprudence. La responsabilité civile joue : le dommage doit être réparé, c'est-à-dire la vitre remboursée. Comme il est mineur, ce sont ses parents qui paieront (voir le « Le savais-tu ? » plus bas).
\end{exemplebox}

L'idée directrice de la responsabilité civile n'est pas de punir, mais de \textbf{réparer} : il s'agit d'effacer, autant que possible, le dommage subi par la victime.

\courssub{La responsabilité pénale : répondre devant la société}

Certaines fautes sont si graves qu'elles ne blessent pas seulement une victime : elles menacent l'ordre social tout entier. Ce sont les \textbf{infractions} (contraventions, délits, crimes). La société, par l'intermédiaire des tribunaux, punit alors leur auteur.

\begin{definition}[Responsabilité pénale]
La \cle{responsabilité pénale} est l'obligation de répondre d'une \textbf{infraction} devant la société et de \textbf{subir la peine} prévue par la loi (amende, emprisonnement, etc.). Elle suppose une infraction définie par le Code pénal et un auteur capable de discernement.
\end{definition}

Contrairement à la responsabilité civile, qui vise la réparation du dommage, la responsabilité pénale vise la \textbf{sanction} de la faute et la protection de la société. Les deux peuvent se cumuler : le chauffard qui a blessé un piéton sera condamné à une peine (pénal) et à payer des dommages et intérêts à la victime (civil).

\begin{exemplebox}[Majorité pénale au Cameroun]
Au Cameroun, le Code pénal fixe la \textbf{majorité pénale à 18 ans} (Loi n°2016/007 du 12 juillet 2016, art. 80). En dessous de cet âge, le mineur délinquant n'est pas jugé comme un adulte : des \textbf{mesures éducatives} (avertissement, placement dans un centre d'éducation, liberté surveillée) remplacent ou encadrent les peines d'adulte. Le droit reconnaît ainsi que le discernement et la maturité du choix se développent avec l'âge : la responsabilité pénale suit la liberté.
\end{exemplebox}

\begin{popfigure}
\begin{center}
\begin{tikzpicture}[
  font=\small,
  racine/.style={rectangle, rounded corners, draw=popdark, fill=popblue!25, thick, text width=4.2cm, align=center},
  branche/.style={rectangle, rounded corners, draw=popdark, thick, text width=3.6cm, align=center},
  sanction/.style={rectangle, rounded corners, draw=popmuted, fill=popcream, text width=3.6cm, align=center, font=\footnotesize},
  lien/.style={-{Stealth[length=2.5mm]}, thick, popdark}
]
\node[racine, align=center] (R) at (0,0){\textbf{La responsabilité}\\ répondre de ses actes};
\node[branche, fill=popgreen!25, align=center] (M) at (-5.4,-2.4){\textbf{Morale}\\ devant sa conscience};
\node[branche, fill=poporange!25, align=center] (C) at (0,-2.4){\textbf{Civile}\\ devant la victime};
\node[branche, fill=poppink!30, align=center] (P) at (5.4,-2.4){\textbf{Pénale}\\ devant la société};
\node[sanction] (SM) at (-5.4,-4.6) {Sanction : remords, culpabilité};
\node[sanction] (SC) at (0,-4.6) {Sanction : réparation du dommage (art. 1382 C. civ.)};
\node[sanction] (SP) at (5.4,-4.6) {Sanction : peine prévue par la loi};
\draw[lien] (R) -- (M);
\draw[lien] (R) -- (C);
\draw[lien] (R) -- (P);
\draw[lien, popmuted] (M) -- (SM);
\draw[lien, popmuted] (C) -- (SC);
\draw[lien, popmuted] (P) -- (SP);
\end{tikzpicture}
\end{center}
\legende{Les trois grandes formes de responsabilité et leurs sanctions associées : la responsabilité morale (sanction intérieure : le remords), la responsabilité civile (réparation du dommage causé à autrui) et la responsabilité pénale (peine infligée au nom de la société).}
\end{popfigure}

\courssub{La responsabilité politique et professionnelle}

Au-delà des rapports entre individus, certaines personnes répondent de leurs actes en raison de la \textbf{fonction} qu'elles occupent.

\begin{itemize}
\item La \textbf{responsabilité politique} : le ministre, le maire, le député répondent de leurs décisions devant le peuple et les institutions. Un ministre dont la politique échoue peut être démis de ses fonctions ; un élu malhonnête peut être poursuivi pour détournement de deniers publics. Dans une démocratie, les élections sont le moment où les citoyens font « rendre des comptes » aux gouvernants.
\item La \textbf{responsabilité professionnelle} : le chef d'entreprise répond de la sécurité de ses employés et de la qualité de ses produits ; le médecin répond de ses soins ; l'enseignant répond de la formation et de la sécurité des élèves qui lui sont confiés. L'enseignant qui abandonne sa classe sans motif engage sa responsabilité devant sa hiérarchie et devant les parents.
\end{itemize}

Ces formes montrent que la responsabilité grandit avec le \textbf{pouvoir} : plus une fonction donne de pouvoir d'agir sur autrui, plus elle charge celui qui l'exerce de comptes à rendre. Comme le dit la sagesse populaire, « à grands pouvoirs, grandes responsabilités ».

\courssub{Les conditions de la responsabilité}

On ne peut pas rendre n'importe qui responsable de n'importe quoi. Pour attribuer légitimement une responsabilité à quelqu'un, trois conditions doivent être réunies.

\begin{propriete}[Les trois conditions de la responsabilité]
Un agent n'est responsable de son acte que si :
\begin{enumerate}
\item la \textbf{conscience de l'acte} : il savait ce qu'il faisait (discernement) ;
\item la \textbf{liberté du choix} : il aurait pu faire autrement, il n'était pas contraint ;
\item le \textbf{lien de causalité} : c'est bien son acte qui a causé le dommage.
\end{enumerate}
Si l'une de ces conditions fait défaut, la responsabilité s'atténue ou disparaît.
\end{propriete}

Vérifions ces conditions sur un cas. Un élève, bousculé violemment dans un couloir, renverse par accident le sac d'un camarade et casse sa calculatrice. A-t-il « causé » le dommage ? Matériellement oui, mais il n'a \emph{pas choisi} d'être bousculé : la condition de liberté fait défaut, sa responsabilité s'efface au profit de celui qui l'a bousculé. De même, un enfant de cinq ans qui met le feu en jouant n'a pas le discernement d'un adulte : la condition de conscience fait défaut.

\begin{popfigure}
\begin{center}
\begin{tikzpicture}[
  font=\small,
  cond/.style={rectangle, rounded corners, draw=popdark, fill=popblue!20, thick, text width=3.4cm, align=center},
  plus/.style={font=\Large\bfseries, text=popdark},
  res/.style={rectangle, rounded corners, draw=popdark, fill=popgreen!30, thick, text width=6.5cm, align=center},
  lien/.style={-{Stealth[length=2.5mm]}, thick, popdark}
]
\node[cond, align=center] (C1) at (-5.2,0){\textbf{Conscience}\\ je savais ce que je faisais};
\node[plus] at (-2.6,0) {$+$};
\node[cond, align=center] (C2) at (0,0){\textbf{Liberté}\\ je pouvais faire autrement};
\node[plus] at (2.6,0) {$+$};
\node[cond, align=center] (C3) at (5.2,0){\textbf{Causalité}\\ mon acte a causé le dommage};
\node[res, align=center] (R) at (0,-2.6){\textbf{Responsabilité}\\ l'agent répond de son acte};
\draw[lien] (C1.south) -- ($(R.north)+(-2.2,0)$);
\draw[lien] (C2.south) -- (R.north);
\draw[lien] (C3.south) -- ($(R.north)+(2.2,0)$);
\end{tikzpicture}
\end{center}
\legende{Organigramme des conditions d'attribution de la responsabilité : la conscience de l'acte, la liberté du choix et le lien de causalité doivent être réunis pour que l'agent soit tenu pour responsable.}
\end{popfigure}

\courssub{Les causes d'atténuation ou d'irresponsabilité}

Lorsque les conditions précédentes ne sont pas remplies, la responsabilité diminue (atténuation) ou s'annule (irresponsabilité). Le droit et la morale reconnaissent plusieurs cas :

\begin{itemize}
\item \textbf{La contrainte} : celui qui agit sous la menace (par exemple un employé forcé par un bandit d'ouvrir le coffre de sa boutique) n'a pas choisi librement ; sa responsabilité est effacée ou fortement atténuée.
\item \textbf{La folie} (trouble mental) : le malade mental qui commet une infraction sans discernement est déclaré pénalement irresponsable ; on le soigne plutôt qu'on ne le punit.
\item \textbf{La minorité} : l'enfant n'a pas encore la pleine maturité du jugement ; c'est pourquoi le droit camerounais prévoit des mesures éducatives pour les mineurs de moins de 18 ans.
\item \textbf{L'ignorance invincible} : celui qui ne pouvait pas savoir (par exemple le chauffeur dont les freins lâchent sans aucun signe annonciateur) n'est pas responsable du dommage imprévisible qui en résulte.
\end{itemize}

\begin{attention}[Responsabilité et culpabilité : ne pas confondre]
Erreur fréquente dans les copies : confondre la \textbf{responsabilité} et la \textbf{culpabilité}. La responsabilité est le fait de \emph{répondre de ses actes} et d'en assumer les conséquences (c'est un statut moral ou juridique). La culpabilité est le \emph{sentiment} d'avoir mal agi (c'est un vécu psychologique). On peut être responsable sans se sentir coupable (le chauffard qui refuse de reconnaître sa faute), et se sentir coupable sans être responsable (le survivant d'un accident qui se reproche d'avoir survécu).
\end{attention}

\begin{savaistu}[La responsabilité du fait d'autrui]
En droit, on distingue la \textbf{responsabilité du fait personnel} (je réponds de ce que j'ai fait moi-même) et la \textbf{responsabilité du fait d'autrui} : les parents répondent des dommages causés par leurs enfants mineurs, les maîtres de leurs employés, les enseignants et artisans de leurs élèves et apprentis pendant le temps où ils les surveillent. Ainsi, si un élève casse du matériel au lycée, sa famille peut être mise à contribution. Cette responsabilité repose sur le devoir de surveillance et d'éducation : on répond de ceux dont on a la charge.
\end{savaistu}

\courssub{Méthode : analyser un cas de responsabilité}

\begin{methode}[Analyser un cas concret de responsabilité]
Face à une situation (sujet de dissertation, cas pratique, fait divers), procéder en quatre étapes :
\begin{enumerate}
\item \textbf{Identifier l'acte et son auteur} : qui a fait quoi, et quel dommage ou quelle faute en résulte ?
\item \textbf{Vérifier la conscience et la liberté} : l'auteur savait-il ce qu'il faisait ? Pouvait-il faire autrement ? (Rechercher contrainte, folie, minorité, ignorance invincible.)
\item \textbf{Établir le lien de causalité} : le dommage vient-il bien de cet acte, ou d'une cause étrangère ?
\item \textbf{Conclure sur le type de responsabilité} : morale (remords), civile (réparation), pénale (peine), politique ou professionnelle, avec atténuation ou irresponsabilité éventuelle.
\end{enumerate}
\end{methode}

\begin{exoresolu}[Le chauffeur de bus et l'accident]
\textbf{Énoncé.} Un chauffeur de bus de la ligne Yaoundé--Bafoussam, après avoir bu de l'alcool au cours d'une pause, perd le contrôle de son véhicule et blesse plusieurs passagers. Analysez sa responsabilité en suivant la méthode en quatre étapes.
\tcblower
\textbf{Solution détaillée.}
\begin{enumerate}
\item \emph{Acte et auteur} : le chauffeur a conduit en état d'ivresse et a provoqué un accident ; le dommage est corporel (passagers blessés) et matériel (bus endommagé).
\item \emph{Conscience et liberté} : le chauffeur est un adulte sain d'esprit ; il savait que boire avant de conduire est dangereux et interdit, et il pouvait choisir de ne pas boire. Aucune cause d'irresponsabilité ne s'applique : il est pleinement conscient et libre.
\item \emph{Causalité} : la perte de contrôle du véhicule résulte directement de son état d'ivresse ; le lien de causalité entre sa faute et le dommage est établi.
\item \emph{Conclusion} : sa responsabilité est pleine et entière, et elle est \textbf{double} : \textbf{pénale} (conduite en état d'ivresse et blessures involontaires, punies par le Code pénal : amende et/ou emprisonnement) et \textbf{civile} (il devra réparer les dommages causés aux passagers et à l'employeur). S'y ajoute une responsabilité \textbf{professionnelle} (perte possible de son emploi) et \textbf{morale} (le remords devant sa conscience).
\end{enumerate}
\end{exoresolu}

\begin{exercice}[À vous de juger]
Dans un atelier de menuiserie à Ngaoundéré, un apprenti de 15 ans, laissé seul sans surveillance, blesse un client en manipulant une machine. Analysez la situation en quatre étapes et déterminez qui répond du dommage : l'apprenti, le patron, ou les deux ? (Pensez à la minorité et à la responsabilité du fait d'autrui.)
\end{exercice}

\begin{corrige}[Exercice : l'apprenti et la machine]
(1) \emph{Acte et auteur} : l'apprenti a manipulé la machine et blessé le client. (2) \emph{Conscience et liberté} : à 15 ans, son discernement est réel mais incomplet (minorité) ; de plus, il a été laissé seul, ce qui relève d'une faute de surveillance du patron. (3) \emph{Causalité} : le dommage vient de la manipulation, mais aussi de l'absence de surveillance. (4) \emph{Conclusion} : la responsabilité pénale de l'apprenti est atténuée par sa minorité (mesures éducatives plutôt que peine d'adulte). La responsabilité civile du dommage incombe principalement au patron, qui répond du fait de son apprenti placé sous sa surveillance (responsabilité du fait d'autrui) et qui a lui-même commis une faute professionnelle en le laissant seul. Les deux sont donc concernés, mais à des titres différents.
\end{corrige}

\begin{aretenir}[L'essentiel de la section]
\begin{itemize}
\item La \cle{responsabilité} est l'obligation morale ou juridique de répondre de ses actes et d'en assumer les conséquences.
\item Elle est \textbf{morale} (devant la conscience : remords), \textbf{civile} (réparation du dommage, art. 1382 du Code civil), \textbf{pénale} (peine pour une infraction), \textbf{politique} et \textbf{professionnelle} (comptes liés à une fonction).
\item Trois conditions : \textbf{conscience}, \textbf{liberté}, \textbf{causalité}. Leur absence entraîne atténuation ou irresponsabilité (contrainte, folie, minorité, ignorance invincible).
\item Ne pas confondre responsabilité (répondre de ses actes) et culpabilité (sentiment d'avoir mal agi).
\end{itemize}
\end{aretenir}

Ainsi définie, la responsabilité apparaît comme l'envers nécessaire de la liberté : on ne répond que de ce qu'on a fait librement. Reste à établir rigoureusement ce lien, ce qui fera l'objet de la section suivante : la liberté comme fondement de la responsabilité.

\coursec{La liberté, fondement de la responsabilité}

\courssub{Transition : de la définition à la fondation}

Dans la section précédente, nous avons défini la responsabilité comme l'obligation de répondre de ses actes devant autrui et devant sa propre conscience, et nous en avons distingué les formes (morale, juridique, civile, pénale). Une question demeure cependant en suspens : \emph{sur quoi repose cette exigence de répondre de ses actes ?} Pourquoi l'homme, et non la pierre qui tombe ou l'animal qui suit son instinct, est-il tenu pour responsable ? La réponse philosophique classique, que nous allons approfondir, est que la \cle{responsabilité} n'est pas un attribut accidentel : elle est la \cle{contrepartie nécessaire de la liberté}. Sans liberté, point de responsabilité ; plus la liberté est grande, plus la responsabilité est lourde. C'est ce lien de fondement que nous allons démontrer.

\courssub{La thèse centrale : on n'est responsable que de ce que l'on a fait librement}

\begin{propriete}[Thèse fondatrice : la liberté, condition \textit{sine qua non} de la responsabilité]
Nul n'est responsable d'un acte qu'il n'a pas librement accompli. La responsabilité suppose un sujet capable de \cle{se déterminer lui-même} (autonomie), c'est-à-dire de choisir entre plusieurs possibles sans y être contraint par une force extérieure ou une nécessité intérieure aveugle. L'acte responsable est un acte \emph{imputable} à un auteur qui en est la \cle{source}.
\end{propriete}

Cette thèse structure tout notre droit et notre morale. Si je suis poussé dans le dos et que je renverse un vase, je suis la \emph{cause} physique du bris, mais non l'\emph{auteur} moral. La responsabilité exige une \cle{volonté} qui s'engage. Analysons les deux figures limites qui éclairent cette condition : la contrainte et l'ignorance.

\courssub{Argument 1 : La contrainte extériorise la cause de l'acte}

Lorsque j'agis sous la menace d'une arme, sous la torture, ou par force majeure (tempête, effondrement), mon corps exécute un mouvement, mais ma \cle{volonté} ne s'y associe pas. Je suis un \emph{instrument}, non un \emph{agent}.

\begin{exemplebox}[Cas de force majeure : le chauffeur de bus de Bafoussam]
Un chauffeur de bus assure la liaison Bafoussam--Douala. Soudain, un enfant traverse la route sans regarder. Le chauffeur freine brutalement, évitant l'enfant mais projetant deux passagers l'un contre l'autre : l'un se fracture le bras, l'autre se blesse au front.
\textbf{Analyse juridique et morale :} Le chauffeur n'est \textbf{pas responsable} des blessures. Il a subi une \cle{contrainte} (le danger imminent pour l'enfant) qui a dicté son acte. Son intention était de sauver une vie, non de blesser. La justice camerounaise (Code pénal, art. 122 al. 1 sur la contrainte irrésistible) ne le poursuivra pas pour coups et blessures involontaires dans ce cas de \emph{force majeure}. L'acte vient de l'événement extérieur, non de sa liberté.
\end{exemplebox}

\begin{attention}[Piège : confondre contrainte physique et pression sociale]
La pression de la famille, la pauvreté, l'opinion publique sont des \emph{influences}, pas des \emph{contraintes irrésistibles} au sens strict. Un jeune qui vole pour « nourrir sa famille » subit une forte pression, mais il conserve une marge de choix (demander de l'aide, travailler, etc.). La justice tiendra compte de la \cle{responsabilité diminuée} (circonstances atténuantes), mais ne l'absoudra pas totalement. Seule une contrainte qui \emph{anéantit} toute possibilité de choisir (pistolet sur la tempe, cataclysme) supprime la responsabilité.
\end{attention}

\courssub{Argument 2 : L'ignorance empêche l'imputation}

Pour être responsable, il faut \emph{savoir} ce que l'on fait et \emph{vouloir} le faire. L'ignorance peut porter sur la nature de l'acte (je donne du poison croyant donner un remède) ou sur sa portée. Deux figures classiques illustrent l'absence de responsabilité par défaut de liberté éclairée :

\begin{itemize}
    \item \textbf{L'enfant :} Il n'a pas encore atteint l'\cle{âge de la raison} (discernement). Au Cameroun, la majorité pénale est fixée à 18 ans (Code de procédure pénale, art. 388 ; Ordonnance n°62/OF/18 du 12 mars 1962 modifiée). En dessous, on parle de \emph{mesures de protection, d'assistance et d'éducation}, non de peine. L'enfant agit, mais ne \emph{choisit} pas pleinement en conscience.
    \item \textbf{Le malade mental / l'aliéné :} Au moment de l'acte, il a perdu l'usage de sa raison (bouffée délirante, schizophrénie aiguë). L'article 64 du Code pénal camerounais dispose : « N'est pas responsable pénalement la personne qui, au moment des faits, était atteinte d'un trouble psychique ou neuropsychique ayant aboli son discernement ou le contrôle de ses actes. » La liberté étant abolie, la responsabilité l'est aussi.
\end{itemize}

\begin{experience}[Expérience de pensée : le somnambule et le criminel]
\textbf{Scénario :} Un homme tue sa femme en plein sommeil (somnambulisme avéré par expertise médicale). Un autre homme tue sa femme après l'avoir planifié pendant des semaines.
\textbf{Observation :} L'acte matériel est identique (homicide). La conséquence est identique (mort de la victime).
\textbf{Interprétation :} Le premier n'a \textbf{pas choisi} ; sa conscience était absente. Il n'est pas \emph{coupable} (irresponsable pénalement), bien qu'il soit l'auteur matériel. Le second a \textbf{choisi} librement ; il est \emph{coupable} et responsable.
\textbf{Conclusion :} La responsabilité ne s'attache pas à l'acte brut, mais à l'\cle{intention libre} qui le précède.
\end{experience}

\courssub{Analyse de texte guidée : Sartre -- « Nous sommes condamnés à être libres »}

\begin{exemplebox}[Extrait : Jean-Paul Sartre, \textit{L'Existentialisme est un humanisme} (1946)]
« L'homme est condamné à être libre. Condamné, parce qu'il ne s'est pas créé lui-même, et qu'il est néanmoins libre, parce qu'une fois jeté dans le monde, il est responsable de tout ce qu'il fait. [...] L'homme n'est rien d'autre que ce qu'il se fait. C'est le premier principe de l'existentialisme. »
\end{exemplebox}

\textbf{Démarche d'analyse (méthode Probatoire / Baccalauréat technique) :}
\begin{enumerate}
    \item \textbf{Thèse :} La liberté n'est pas un don agréable, mais une \cle{condamnation} : nous n'avons pas choisi d'exister, mais nous sommes forcés de choisir à chaque instant.
    \item \textbf{Argument :} « L'homme n'est rien d'autre que ce qu'il se fait. » Il n'y a pas de \cle{nature humaine} prédéfinie (essence) qui dicterait nos actes. L'\cle{existence précède l'essence}. Je me crée par mes choix.
    \item \textbf{Conséquence radicale :} Puisque je me choisis \emph{totalement}, je suis responsable \emph{de tout ce que je suis et de tout ce que je fais}. Je ne peux invoquer ni mon tempérament, ni mon milieu, ni mes gènes, ni Dieu. « L'existentialisme [...] déclare que l'homme est responsable de sa propre passion. »
    \item \textbf{Portée :} Cette responsabilité est \emph{absolue} et \emph{universelle}. En choisissant pour moi, je choisis pour l'humanité entière (je légifère pour tous).
\end{enumerate}

\begin{aretenir}[Sartre et la responsabilité totale]
Pour Sartre, la liberté est \emph{ontologique} : elle est la structure même de l'être humain. Fuir sa liberté (mauvaise foi), c'est mentir à soi-même en prétendant « je n'avais pas le choix ». Mais « ne pas choisir, c'est encore choisir ». La responsabilité est le \emph{fardeau} indissociable de cette liberté radicale.
\end{aretenir}

\courssub{Analyse de texte guidée : Kant -- L'impératif catégorique suppose la liberté}

\begin{exemplebox}[Extrait : Emmanuel Kant, \textit{Fondements de la métaphysique des mœurs} (1785)]
« Un libre arbitre et un arbitre soumis à des lois morales sont une seule et même chose. [...] L'impératif catégorique n'est possible que sous la présupposition de la liberté. »
\end{exemplebox}

\textbf{Démarche d'analyse :}
\begin{enumerate}
    \item \textbf{Contexte :} Kant cherche le fondement \emph{a priori} de la morale. Il distingue les \emph{impératifs hypothétiques} (si tu veux X, fais Y -- conditionnels, hétéronomes) de l'\emph{impératif catégorique} (fais Y \emph{inconditionnellement} -- autonome).
    \item \textbf{Le lien Liberté / Loi morale :} Agir moralement, c'est agir \emph{par devoir}, c'est-à-dire par respect pour la loi que \emph{je me donne à moi-même} (autonomie : \textit{autos} = soi, \textit{nomos} = loi).
    \item \textbf{La liberté comme \cle{présupposition} :} Je ne peux \emph{prouver} la liberté théoriquement (phénomène vs noumène), mais je \emph{dois} la postuler pour que la morale ait un sens. Si je n'étais pas libre, je ne ferais qu'obéir à des causes naturelles (désirs, penchants, déterminismes) : mon acte serait un \emph{fait naturel}, non un \emph{acte moral}.
    \item \textbf{Réciprocité :} « Un libre arbitre et un arbitre soumis à des lois morales sont une seule et même chose. » Être libre, c'est \emph{nécessairement} se soumettre à la loi morale (universalisable) ; se soumettre à la loi morale, c'est \emph{exercer} sa liberté.
\end{enumerate}

\begin{propriete}[Kant : La liberté est le \emph{ratio essendi} de la responsabilité morale]
Sans liberté (autonomie), l'impératif catégorique est vide. La responsabilité morale (le devoir, la dignité, le respect de l'humanité en moi et en autrui) \emph{suppose} et \emph{exige} la liberté comme condition de possibilité. Je suis responsable \emph{parce que} je suis libre, et je suis libre \emph{parce que} je suis un être rationnel capable de se donner la loi.
\end{propriete}

\courssub{Objection majeure : Le déterminisme -- La liberté est-elle une illusion ?}

\begin{definition}[Déterminisme (universel / causal)]
Doctrine selon laquelle \textbf{tous les événements}, y compris les actes humains, les décisions et les pensées, sont les \emph{effets nécessaires} de causes antérieures (physiques, biologiques, psychologiques, sociales). Rien ne « pourrait être autrement » que ce qui est. Il n'y a pas de place pour un « premier moteur immobile » (libre arbitre) dans la chaîne causale.
\end{definition}

\textbf{Les formes du déterminisme :}
\begin{itemize}
    \item \textbf{Physique / Neuroscientifique :} Mes neurones s'activent selon des lois physico-chimiques ; ma « décision » n'est que l'épiphanie consciente d'un processus cérébral déjà enclenché (expériences de Libet, IRMf).
    \item \textbf{Psychologique / Psychanalytique :} Mes choix sont dictés par mon inconscient, mes pulsions, mon histoire infantile (Freud).
    \item \textbf{Social / Sociologique :} Je suis le produit de ma classe, de ma culture, de mon époque (Marx, Bourdieu, Durkheim). « L'individu est un carrefour de déterminismes. »
    \item \textbf{Théologique / Fatalisme :} Tout est écrit d'avance par la volonté divine ou le Destin.
\end{itemize}

\begin{attention}[Erreur fréquente au Probatoire / Baccalauréat technique : « Puisque tout a une cause, je ne suis pas responsable »]
C'est l'argument du \emph{déterministe vulgaire} pour se dédouaner. \textbf{La société (le droit, la morale courante) ne juge pas les causes lointaines, elle juge l'acte et l'intention présente.} Même si mes gènes et mon enfance m'ont \emph{prédisposé} à la violence, au moment de l'acte, \emph{je} tiens le couteau, \emph{je} vise, \emph{je} frappe. La \cle{conscience du choix} (même brève, même contrainte) suffit à fonder le jugement moral et juridique. Le droit sépare la \emph{culpabilité} (subjective, liée à la volonté) de la \emph{dangerosité} (objective, liée aux causes).
\end{attention}

\begin{propriete}[Réponse philosophique : Le compatibilisme (ou « liberté dans la nécessité »)]
On peut admettre le déterminisme \emph{explicatif} (mes actes ont des causes) sans renoncer à la liberté \emph{pratique} (je délibère, je choisis, je pourrais faire autrement \emph{si je le voulais}).
\begin{itemize}
    \item \textbf{Hume / Stoïciens :} La liberté n'est pas l'absence de cause, mais l'absence de \emph{contrainte extérieure}. Je suis libre si mon acte procède de \emph{mes} désirs et \emph{ma} raison, non d'une arme sur la tempe.
    \item \textbf{Ricoeur :} L'homme est un \emph{être capable}. Capable de dire « je », capable d'agir, capable de se raconter (identité narrative), capable d'imputabilité. La responsabilité est l'\emph{attribution} d'un acte à un agent capable de répondre.
\end{itemize}
La \cle{conscience du choix} (le « je peux ») est le fait premier, indéniable, sur lequel repose l'édifice moral et juridique, quelles que soient les causes scientifiques ultérieures.
\end{propriete}

\courssub{Conséquence pratique : La justice camerounaise gradue la peine selon la liberté}

Le Code pénal camerounais (Loi n°2016/007 du 12 juillet 2016) et le Code de procédure pénale incarnaient concrètement le lien \textbf{Liberté $\rightarrow$ Responsabilité}. La peine n'est pas une vengeance, mais une réponse proportionnée à la \cle{culpabilité} (degré de liberté engagée).

\begin{center}
\begin{tabularx}{\linewidth}{|l|X|X|}
\hline
\rowcolor{popblueL}
\textbf{Degré de liberté dans l'acte} & \textbf{Qualification juridique (Code pénal camerounais)} & \textbf{Conséquence sur la peine (Exemples)} \\
\hline
\cle{Liberté pleine} (Preméditation, guet-apens) & Assassinat (Art. 276) & Peine maximale : \textbf{Peine de mort} (encourue, moratoire de fait depuis 1997) ou \textbf{Réclusion criminelle à perpétuité}. Responsabilité \textbf{maximale}. \\
\hline
\cle{Liberté intentionnelle simple} (Coup porté volontairement sans préméditation) & Meurtre (Art. 275) ou Coups et blessures volontaires (Art. 277+) & Peine de \textbf{réclusion} (10--20 ans pour meurtre). Responsabilité \textbf{entière}. \\
\hline
\cle{Liberté diminuée} (Provocation, état passionnel, troubles psychiques n'abolissant pas le discernement) & Circonstances \textbf{atténuantes} (Art. 65) & Réduction de la peine (ex: 10 ans au lieu de 20). Responsabilité \textbf{partielle}. \\
\hline
\cle{Liberté nulle} (Force majeure, contrainte physique irrésistible, aliénation mentale abolissant le discernement) & \textbf{Irresponsabilité pénale} (Art. 64, Art. 122 al. 1) & \textbf{Aucune peine}. Mesures de sûreté (soins, surveillance) si dangerosité. Responsabilité \textbf{nulle}. \\
\hline
\cle{Non-liberté} (Accident pur, cas fortuit, légitime défense) & \textbf{Non-lieu} / Absolution / Justification (Art. 122 al. 6 Légitime défense) & Aucune sanction. L'acte n'est pas imputable à une volonté coupable. \\
\hline
\end{tabularx}
\end{center}

\begin{exemplebox}[Cas concret : Accident de la route à Yaoundé]
\textbf{Cas A :} Un conducteur roule à 100 km/h en zone urbaine (limite 60), téléphone au volant, percute un piéton sur un passage clouté. \textbf{Analyse :} Il a \emph{choisi} la vitesse, le téléphone. Liberté pleine engagée $\rightarrow$ \emph{Homicide involontaire avec circonstances aggravantes} (mise en danger d'autrui, Art. 279-1). Peine lourde.
\textbf{Cas B :} Un conducteur respecte le code, un enfant surgit soudainement derrière un bus arrêté. Freinage impossible. \textbf{Analyse :} Aucune faute, aucune liberté de choisir l'évitement. $\rightarrow$ \emph{Non-lieu} (cas fortuit / force majeure, Art. 122). Pas de responsabilité pénale.
\textbf{Cas C :} Conducteur ivre (choix de boire) $\rightarrow$ accident. \textbf{Analyse :} La liberté s'est exercée \emph{avant} (boire), la responsabilité rejaillit sur la conséquence. \emph{Ivresse volontaire = circonstance aggravante} (Art. 66), non excusante.
\end{exemplebox}

\courssub{Illustration : La chaîne causale de la responsabilité}

\begin{popfigure}
\begin{tikzpicture}[
    node distance=1.8cm and 1.2cm,
    box/.style={rectangle, draw=popdark, thick, fill=popcream, rounded corners=6pt, minimum height=1.2cm, minimum width=2.8cm, align=center, font=\small\bfseries},
    arrow/.style={-Stealth, thick, popblue, shorten >=2pt, shorten <=2pt},
    labelstyle/.style={font=\footnotesize\itshape, text=popmuted, midway, above, sloped}
]
% Nœuds
\node[box, align=center] (lib){LIBERTÉ\\(Autonomie,\\Capacité de choisir)};
\node[box, right=of lib, align=center] (choix){CHOIX\\ (Délibération,\\Intention, Volonté)};
\node[box, right=of choix, align=center] (acte){ACTE\\ (Comportement\\extériorisé)};
\node[box, right=of acte, align=center] (cons){CONSÉQUENCES\\ (Effets réels\\sur autrui/monde)};
\node[box, right=of cons, fill=poporangeL, draw=poporange, align=center] (resp){RESPONSABILITÉ\\ (Imputabilité,\\Obligation de répondre)};

% Flèches principales
\draw[arrow] (lib) -- node[labelstyle] {permet} (choix);
\draw[arrow] (choix) -- node[labelstyle] {déclenche} (acte);
\draw[arrow] (acte) -- node[labelstyle] {produit} (cons);
\draw[arrow] (cons) -- node[labelstyle] {fonde} (resp);

% Boucle de rétroaction (anticipation)
\draw[arrow, poppurple, bend left=40] (resp.north) to node[labelstyle, pos=0.5, above, align=center]{Anticipation\\éthique (Kant)} (lib.north);

% Annotations latérales
\node[left=0.3cm of lib, font=\footnotesize, text=popmuted, align=right] {Condition\\de possibilité};
\node[right=0.3cm of resp, font=\footnotesize, text=popmuted, align=left] {Exigence\\de réponse};
\node[below=0.2cm of choix, font=\footnotesize, text=popgreen] {\checkmark\ Conscience};
\node[below=0.2cm of acte, font=\footnotesize, text=popgreen] {\checkmark\ Volonté};
\node[below=0.2cm of cons, font=\footnotesize, text=poporange] {\checkmark\ Causalité};
\node[below=0.2cm of resp, font=\footnotesize, text=popred] {\checkmark\ Imputabilité};
\end{tikzpicture}
\legende{Schéma fonctionnel : La liberté (source) rend possible le choix (intention), qui génère l'acte (cause efficiente), produisant des conséquences (effets), qui fondent l'imputabilité responsable (jugement moral/juridique). La flèche violette montre le retour réflexif : j'anticipe ma responsabilité \emph{avant} d'agir (impératif catégorique).}
\end{popfigure}

\courssub{Illustration : Tableau de correspondance -- Degrés de liberté et de responsabilité}

\begin{popfigure}
\begin{center}
\begin{tabularx}{\linewidth}{|>{\bfseries\centering}m{2.8cm}|X|>{\centering}m{2.5cm}|X|}
\hline
\rowcolor{popblue}
\textcolor{white}{\textbf{DEGRÉ DE LIBERTÉ}} & \textcolor{white}{\textbf{CARACTÉRISTIQUES DE L'ACTE}} & \textcolor{white}{\textbf{DEGRÉ DE RESPONSABILITÉ}} & \textcolor{white}{\textbf{SANCTION / RÉPONSE SOCIALE}} \\
\hline
\rowcolor{popblueL}
PLEINE & Choix réfléchi, intention claire, absence de contrainte, discernement intact & PLEINE ET ENTIÈRE & Peine maximale (réclusion, amende max), réparation intégrale, blâme moral fort \\
\hline
DIMINUÉE & Provocation, émotion violente, troubles psychiques partiels, contrainte morale forte & PARTIELLE / ATTÉNUÉE & Peine réduite (circonstances atténuantes Art. 65), réparation partielle, suivi socio-éducatif \\
\hline
\rowcolor{popblueL}
NULLE / ABOLIE & Force majeure, contrainte physique irrésistible, aliénation mentale (discernement aboli Art. 64), enfant sans discernement & NULLE (Irresponsabilité) & Aucune peine pénale. Mesures de protection / soins / éducation (mineurs), non-lieu \\
\hline
NON CONCERNÉE (Acte non humain) & Cas fortuit, accident pur, légitime défense, acte réflexe, somnambulisme & AUCUNE (Non-imputabilité) & Non-lieu, absolution. Indemnisation possible (assurance, solidarité) sans faute \\
\hline
\end{tabularx}
\end{center}
\legende{Tableau de correspondance synthétique : Plus la liberté effective (capacité de choisir en conscience) est grande, plus la responsabilité (imputabilité) est lourde, et plus la sanction sociale (peine, réparation, blâme) est sévère. Ce tableau est un outil direct pour l'analyse de cas pratiques (Probatoire / Baccalauréat technique).}
\end{popfigure}

\courssub{Synthèse : Liberté et responsabilité, un couple indissociable}

\begin{aretenir}[Synthèse finale de la section]
\begin{enumerate}
    \item \textbf{Liberté condition nécessaire :} On ne juge que ce qui dépend de nous. \textbf{Point de responsabilité sans liberté.} (Contrainte, ignorance, aliénation $\rightarrow$ irresponsabilité).
    \item \textbf{Liberté condition suffisante (pratique) :} Dès lors qu'un sujet \emph{choisit} en conscience (même sous influence), il devient \emph{auteur} et \emph{répondant}. La société ne peut pas ne pas juger.
    \item \textbf{Proportionnalité :} \textbf{Plus on est libre, plus on est responsable.} L'adulte sain d'esprit, informé, non contraint, porte la plénitude de la responsabilité. L'enfant, le malade, la victime de force majeure en portent une part décroissante jusqu'à zéro.
    \item \textbf{Fondement ontologique (Sartre) :} L'homme \emph{est} sa liberté ; il ne peut s'en déposséder. La responsabilité est l'ombre inséparable de cette liberté.
    \item \textbf{Fondement normatif (Kant) :} La loi morale (devoir) ne s'adresse qu'à des êtres libres. La responsabilité (dignité) est le corrélat de l'autonomie.
    \item \textbf{Réalité camerounaise :} Le Code pénal traduit cette philosophie : il module la peine (mort, réclusion, amende, mesure éducative, non-lieu) \emph{exclusivement} selon le degré de liberté avéré dans l'acte (discernement, intention, contrainte).
\end{enumerate}
\end{aretenir}

\begin{methode}[Pour la dissertation / le commentaire : Argumenter le lien Liberté--Responsabilité]
\begin{enumerate}
    \item \textbf{Définir les termes} : Liberté (pouvoir de s'autodéterminer), Responsabilité (obligation de répondre de ses actes).
    \item \textbf{Poser la thèse} : La liberté est le \emph{fondement} (condition de possibilité) de la responsabilité. « Nul n'est responsable que de ce qu'il a fait librement. »
    \item \textbf{Argumenter par l'analyse des contre-exemples} : Montrer que contrainte (force majeure) et ignorance (enfant, fou, erreur) suppriment la responsabilité $\rightarrow$ preuve par l'absurde que la liberté est nécessaire.
    \item \textbf{Apporter les références majeures} :
        \begin{itemize}
            \item \textbf{Sartre :} Responsabilité absolue (condamnés à être libres, existence précède essence).
            \item \textbf{Kant :} Autonomie = liberté + loi morale ; responsabilité = dignité de l'agent rationnel.
            \item \textbf{Droit positif (Cameroun) :} Art. 64, 65, 66, 122, 275, 276 C. pén. (graduation peine selon discernement/liberté).
        \end{itemize}
    \item \textbf{Traiter l'objection déterministe} : Distinguer causes explicatives (sciences) et raisons justificatives (morale/droit). La conscience du choix \emph{ici-maintenant} fonde l'imputabilité.
    \item \textbf{Conclure} : Liberté et responsabilité sont les deux faces d'une même médaille : la \cle{subjectivité morale}. Être homme, c'est être \emph{capable de répondre}.
\end{enumerate}
\end{methode}

\begin{exoresolu}[Sujet type Probatoire / Baccalauréat technique : « Peut-on être responsable sans être libre ? »]
\textbf{Plan détaillé :}
\textbf{Introduction :} Accroche (ex: procès d'un enfant soldat / chauffeur Bafoussam). Définitions. Problématique : Le lien est-il nécessaire ou contingent ? Thèse : La liberté est la condition \textit{sine qua non} de la responsabilité, mais le déterminisme complexifie l'exercice concret de cette liberté.
\textbf{I. La liberté, condition nécessaire de la responsabilité (Preuve par les cas limites).}
    A. Contrainte physique/force majeure $\rightarrow$ irresponsabilité (ex: chauffeur, pistolet sur tempe). Art 122 C. pén.
    B. Absence de discernement (enfant, aliéné, somnambule) $\rightarrow$ irresponsabilité. Art 64 C. pén.
    C. Conclusion partielle : Sans liberté (choix conscient), point d'auteur moral, point de responsabilité.
\textbf{II. La liberté comme fondement de la responsabilité pleine (Philosophie et Droit).}
    A. Sartre : Liberté ontologique $\rightarrow$ responsabilité totale (mauvaise foi = fuite).
    B. Kant : Autonomie (loi morale auto-donnée) $\rightarrow$ Dignité $\rightarrow$ Responsabilité comme attribution.
    C. Droit camerounais : Preméditation (liberté max) = peine max. Intention simple = peine standard.
\textbf{III. L'objection déterministe et la responsabilité « compatibiliste ».}
    A. Déterminisme (neuro, psycho, socio) : causes totales $\rightarrow$ liberté illusion ?
    B. Réponse : Distinction niveau d'explication (causes) / niveau d'imputation (raisons). Hume/Ricoeur : liberté = absence de contrainte \emph{extérieure} + action selon \emph{ses} motifs.
    C. Pratique : La justice juge l'acte intentionnel, pas le génome. Responsabilité maintenue pour la vie sociale.
\textbf{Conclusion :} Oui, la liberté est le fondement. Mais « être libre » n'est pas binaire (tout/rien) : c'est un \emph{degré}. La responsabilité suit ce degré. Être responsable, c'est assumer d'être l'auteur de ses choix, même contraints par l'histoire.
\tcblower
\textbf{Points de vigilance (Attention) :}
\begin{itemize}
    \item Ne pas confondre \emph{responsabilité civile} (réparation du dommage, parfois sans faute) et \emph{responsabilité pénale/morale} (faute, liberté requise). Le sujet porte sur la responsabilité \emph{morale/pénale}.
    \item Ne pas dire « le déterminisme prouve qu'il n'y a pas de liberté ». Dire : « Le déterminisme est une thèse scientifique ; la responsabilité est une exigence pratique/juridique. Elles opèrent sur des plans différents (fait / valeur). »
    \item Citer \textbf{précisément} les articles du Code pénal camerounais (64, 65, 66, 122, 275, 276) : c'est du concret attendu au Bac Tech.
\end{itemize}
\end{exoresolu}

\begin{savaistu}[Sartre et l'impossible fuite]
Sartre écrit dans \textit{L'Être et le Néant} : « Je suis mes choix. Je ne peux pas \emph{ne pas} choisir. Si je ne choisis pas, c'est encore un choix. » Même le refus de décider (passivité, soumission à l'autorité, « j'obéissais aux ordres ») est un \emph{choix} de ne pas exercer sa liberté critique. C'est pourquoi les procès de Nuremberg (1945-46) ont condamné des nazis qui invoquaient l'obéissance : ils avaient \emph{choisi} d'obéir. Au Cameroun, un militaire qui exécute un « ordre manifestement illégal » (ex: tirer sur des civils désarmés) est \textbf{pénalement responsable} (Code de justice militaire). Il ne peut pas dire « je n'avais pas le choix ». Il \emph{avait} le choix de désobéir, même au prix de sa vie. C'est le prix de la liberté : la responsabilité est \emph{inescapable}.
\end{savaistu}

\coursec{Applications concrètes au quotidien camerounais}

La section précédente a établi que la \cle{liberté} est le fondement ontologique et juridique de la \cle{responsabilité}~: nul ne peut être tenu pour responsable d'un acte qu'il n'a pas accompli librement. Mais cette articulation théorique ne prend tout son sens que lorsqu'elle est confrontée à la complexité du réel. Au Cameroun, comme ailleurs, les situations concrètes — tricherie à l'examen, corruption, accidents de la route, rumeurs numériques, pressions traditionnelles — mettent à l'épreuve la distinction entre liberté formelle et liberté effective, entre responsabilité légale et responsabilité morale. C'est à cette épreuve de la vie quotidienne que nous allons maintenant soumettre nos concepts.

\courssub{Étude de cas 1 : L'élève et la tricherie au Probatoire — Liberté du choix et responsabilité devant sa conscience}

\begin{exemplebox}[Cas concret : La tentation de la « feuille »]
Imagine un élève de Terminale technique, appelons-le \emph{Jean}, assis dans la salle d'examen du Probatoire. Il bloque sur un sujet de dissertation en philosophie. Son voisin de table, plus à l'aise, a déjà rédigé deux pages. Jean a une \emph{feuille de brouillon} préparée secrètement avec des arguments clés. Il hésite : la copier discrètement ou rester fidèle à son travail, au risque d'échouer.
\end{exemplebox}

\noindent\textbf{Analyse philosophique.}
Ce cas illustre la \cle{liberté d'indifférence} (choix entre deux possibles : tricher ou ne pas tricher) et la \cle{liberté comme autonomie} (agir selon sa propre loi morale).
\begin{itemize}
    \item \textbf{L'acte libre :} Jean n'est pas physiquement contraint. La contrainte est psychologique (peur de l'échec, pression familiale). Au sens strict, il \emph{peut} ne pas tricher. Sa liberté est intacte, même si elle est « située » (Sartre) dans un contexte anxiogène.
    \item \textbf{La responsabilité :} Si Jean triche, il est \emph{auteur} de l'acte. Il en répond devant le jury (sanction disciplinaire : annulation de l'épreuve, exclusion) et devant sa \emph{conscience} (sentiment de culpabilité, atteinte à l'estime de soi).
    \item \textbf{Le paradoxe :} En trichant pour « réussir », Jean nie sa liberté réelle (il se fait passer pour ce qu'il n'est pas) et détruit la valeur du diplôme qu'il convoite. Il devient \emph{irresponsable} vis-à-vis de lui-même.
\end{itemize}

\begin{attention}[Piège à éviter : La victimisation]
Ne pas confondre \emph{explication} et \emph{justification}. Dire « J'ai triché parce que la pression était trop forte » explique le mobile, mais ne justifie pas l'acte. La responsabilité suppose que l'agent \emph{aurait pu} faire autrement. L'excuse « c'est plus fort que moi » relève de la \cle{mauvaise foi} (Sartre) : on se ment à soi-même pour fuir l'angoisse de la liberté.
\end{attention}

\courssub{Étude de cas 2 : Le fonctionnaire et le détournement de fonds publics — Responsabilité pénale et morale}

\begin{exemplebox}[Cas concret : La « gestion » du budget de fonctionnement]
Un chef de service dans une administration publique à Yaoundé perçoit des fonds destinés à l'entretien des locaux. Il décide d'en détourner une partie pour financer les études de ses enfants à l'étranger, justifiant cela par un salaire « insuffisant » et par la pratique prétendument généralisée (« Tout le monde le fait »).
\end{exemplebox}

\noindent\textbf{Analyse philosophique et juridique.}
\begin{itemize}
    \item \textbf{Responsabilité pénale (juridique) :} Le Code pénal camerounais (Loi n°2016/007 du 12 juillet 2016) sanctionne le \emph{détournement de deniers publics} (articles 134 et suivants). Le fonctionnaire est un \emph{agent public} investi d'une mission. Les sanctions prévues sont l'emprisonnement, l'amende et la confiscation des biens. La loi ne tient compte ni de la « coutume » ni du bas salaire : la liberté de gérer les fonds impliquait l'obligation de les gérer loyalement.
    \item \textbf{Responsabilité morale :} Au-delà du droit, il y a trahison de la \emph{confiance} (contrat social implicite entre l'État et l'agent) et atteinte au \emph{bien commun}. L'argument « tout le monde le fait » est une \cle{généralisation fallacieuse} : la multiplicité des coupables n'efface pas la faute individuelle.
    \item \textbf{Liberté et occasion :} Avoir l'\emph{occasion} de voler (l'accès à la caisse) n'est pas une \emph{obligation} de voler. La liberté du fonctionnaire s'est exercée \emph{contre} son devoir.
\end{itemize}

\begin{savaistu}[Lutte contre la corruption au Cameroun]
La \cle{CONAC} (Commission Nationale Anti-Corruption, créée en 2006) et le \cle{Tribunal Criminel Spécial (TCS)}, créé par la loi n°2011/028 du 14 décembre 2011, reposent sur le principe que \textbf{chaque agent public est personnellement responsable} de la gestion des deniers publics. L'opération dite « Épervier », lancée en 2006, vise à briser le sentiment d'impunité en montrant que la responsabilité pénale peut atteindre tout agent, quel que soit son rang. Cela illustre concrètement que la liberté de gérer implique la responsabilité de rendre compte.
\end{savaistu}

\courssub{Étude de cas 3 : Le conducteur ivre sur l'axe Yaoundé-Douala — L'alcool diminue-t-il la responsabilité ?}

\begin{exemplebox}[Cas concret : Accident mortel à la sortie de Yaoundé]
Un conducteur de poids lourd, ayant consommé de l'alcool dans une buvette, prend le volant sur l'axe Yaoundé-Douala. Son taux d'alcoolémie est supérieur à la limite légale. Il perd le contrôle de son véhicule et percute un car de transport en commun : 3 morts, 12 blessés graves. Au procès, son avocat plaide l'\emph{altération du discernement} due à l'alcool pour demander l'irresponsabilité pénale de son client (article 23 du Code pénal : n'est pénalement responsable celui qui était atteint, au moment des faits, d'un trouble psychique ayant aboli son discernement).
\end{exemplebox}

\noindent\textbf{Analyse philosophique et juridique.}
\begin{itemize}
    \item \textbf{La cause de la cause :} L'état d'ivresse n'est pas survenu par magie. Le conducteur a \emph{librement choisi} de boire \emph{avant} de conduire. Il a \emph{volontairement} mis sa lucidité « entre parenthèses » pour un temps, en acceptant le risque de perdre le contrôle. C'est ce que le droit appelle la \cle{responsabilité par anticipation} : on répond de la faute commise en se mettant soi-même hors d'état de se maîtriser.
    \item \textbf{Atténuation ou suppression :} L'alcool peut, dans certains cas, \emph{atténuer} la responsabilité (trouble psychique passager reconnu par expertise), mais il ne la \emph{supprime} presque jamais lorsqu'il s'agit de conduite en état d'ivresse, car la mise en danger d'autrui était prévisible et acceptée.
    \item \textbf{Responsabilité civile :} Au-delà du pénal, le conducteur (et son assureur) doit \emph{réparer le dommage} causé (dommages-intérêts aux victimes et à leurs familles). La responsabilité civile repose sur le principe que tout fait de l'homme qui cause un dommage à autrui oblige celui par la faute duquel il est arrivé à le réparer.
\end{itemize}

\begin{attention}[Erreur fréquente : « C'est l'alcool qui a conduit »]
Croire que l'alcool ou la colère \textbf{suppriment} la responsabilité est une erreur majeure. Ils l'\textbf{atténuent} parfois (si un trouble psychique est médicalement prouvé), mais ne l'effacent pas. Le juge sanctionne le choix initial de s'enivrer en sachant qu'on devra conduire. \textbf{La liberté de boire engage la responsabilité de conduire.}
\end{attention}

\courssub{Étude de cas 4 : La rumeur WhatsApp et les violences — Responsabilité de l'auteur et des relais}

\begin{exemplebox}[Cas concret : « Enlèvement d'enfants à Douala »]
Un message vocal circule sur WhatsApp : « Des hommes en voiture blanche enlèvent des enfants pour la traite d'organes ! » Le message est faux (fausse information, ou \emph{fake news}). Paniqués, des habitants lynchent un passant innocent, pris pour un ravisseur. L'auteur du message est anonyme ; les relais (ceux qui ont transféré le message) sont identifiables.
\end{exemplebox}

\noindent\textbf{Analyse philosophique et juridique.}
\begin{itemize}
    \item \textbf{Liberté d'expression et abus :} La liberté d'expression (article 19 de la Déclaration universelle des droits de l'homme ; préambule de la Constitution camerounaise du 18 janvier 1996) n'est pas absolue. Elle s'arrête où commence l'atteinte à l'ordre public, à l'honneur d'autrui ou l'incitation à la violence.
    \item \textbf{Responsabilité de l'auteur :} Créer sciemment une fausse alerte est un délit (diffamation, propagation de fausses nouvelles, mise en danger de la vie d'autrui). L'anonymat numérique ne garantit pas l'impunité : la loi n°2010/012 du 21 décembre 2010 relative à la cybersécurité et à la cybercriminalité au Cameroun permet la traçabilité des communications électroniques.
    \item \textbf{Responsabilité des relais (partage) :} Transmettre sans vérifier (\emph{partage viral}) engage la \cle{responsabilité morale} et parfois pénale du relais. « Je n'ai fait que transmettre » n'est pas une excuse valable si le contenu est manifestement diffamatoire ou incendiaire : on est responsable de ce qu'on \emph{amplifie}.
    \item \textbf{Responsabilité collective :} La violence de la foule (lynchage) montre une \emph{aliénation de la liberté} : l'individu semble disparaître dans le « nous » grégaire (Gustave Le Bon, \emph{Psychologie des foules}, 1895). Chaque participant reste pourtant \emph{personnellement} responsable de ses actes.
\end{itemize}

\begin{exemplebox}[Exemple : Le droit camerounais face aux réseaux sociaux]
Le Code pénal camerounais punit la \emph{diffamation} et la propagation de fausses nouvelles « par quelque moyen que ce soit », y compris les réseaux sociaux. La loi n°2010/012 sur la cybersécurité et la cybercriminalité sanctionne l'utilisation des technologies de l'information pour nuire. \textbf{La liberté d'expression a des limites légales :} on est responsable de ses mots, même — et surtout — s'ils sont numériques et viraux.
\end{exemplebox}

\courssub{Débat dirigé : Tradition, pression familiale et liberté individuelle}

\begin{methode}[Conduite d'un débat philosophique dirigé]
\begin{enumerate}
    \item \textbf{Poser la thèse :} « La tradition (mariage forcé, choix d'orientation imposé) annule la liberté individuelle. »
    \item \textbf{Poser l'antithèse :} « La tradition structure la liberté ; sans cadre culturel, l'individu est perdu. Le devoir filial prime. »
    \item \textbf{Distinctions conceptuelles :}
    \begin{itemize}
        \item \cle{Liberté négative} (absence de contrainte extérieure) et \cle{liberté positive} (capacité effective d'autodétermination, d'autonomie).
        \item \cle{Contrainte physique} (enfermement, menaces) et \cle{pression morale} (honte, exclusion, « qu'en-dira-t-on »).
    \end{itemize}
    \item \textbf{Argumenter :}
    \begin{itemize}
        \item \emph{Pour la thèse :} Le mariage forcé est une violation d'un droit fondamental ; le droit camerounais exige le consentement libre des époux. Imposer une filière (par exemple la médecine au lieu des arts) brise l'autonomie et engendre l'échec ou le ressentiment.
        \item \emph{Pour l'antithèse :} La famille finance les études ; elle estime avoir un « droit de regard ». La tradition protège l'individu de l'isolement. Le respect des aînés est une vertu (piété filiale), non une servitude.
    \end{itemize}
    \item \textbf{Synthèse :} La \cle{liberté effective} nécessite un \emph{seuil d'autonomie matérielle et psychique}. La tradition légitime devient oppression quand elle nie la personnalité et le libre consentement de l'individu. La responsabilité du jeune adulte est de \emph{négocier} son émancipation, non de la subir ni de la rejeter violemment.
\end{enumerate}
\end{methode}

\courssub{TP guidé : Construction d'un tableau d'analyse de cas réels}

\begin{methode}[Méthode pour traiter une étude de cas en philosophie]
\begin{enumerate}
    \item \textbf{Décrire les faits objectivement :} Qui ? Quoi ? Où ? Quand ? Comment ? (Sans jugement moral.)
    \item \textbf{Dégager le problème philosophique :} Quelle notion de liberté ou de responsabilité est en jeu ? (Contrainte, ignorance, passion, devoir, droit.)
    \item \textbf{Appliquer les notions du cours :} Identifier le type de liberté (indifférence, autonomie, spontanéité), le type de responsabilité (pénale, civile, morale), le lien de causalité et l'imputabilité.
    \item \textbf{Argumenter un jugement nuancé :} Distinguer les degrés de responsabilité (pleine, atténuée, partagée). Proposer une sanction juste (réparatrice, éducative, punitive).
    \item \textbf{Conclure :} Quel éclairage ce cas apporte-t-il sur la condition humaine ?
\end{enumerate}
\end{methode}

\noindent\textbf{Consigne du TP :} Complétez le tableau ci-dessous pour les 4 études de cas traitées (et ajoutez-en un 5\textsuperscript{e} de votre choix : grève sauvage, pollution industrielle, harcèlement scolaire, etc.).

\begin{popfigure}
\begin{center}
\begin{tabularx}{\linewidth}{|l|X|X|X|}
\hline
\rowcolor{popblueL} \textbf{Acte (faits)} & \textbf{Liberté en jeu (type / degré)} & \textbf{Type de responsabilité} & \textbf{Sanction / réparation méritée} \\
\hline
Élève qui triche au Probatoire (copie sur le voisin, antisèche) & Liberté d'indifférence (choix binaire). Pression psychologique (peur de l'échec), mais pas de contrainte physique. & Morale (conscience) et disciplinaire (règlement de l'examen). & Annulation de l'épreuve, exclusion temporaire ; travail sur l'intégrité. \\
\hline
Fonctionnaire qui détourne des fonds publics (budget de fonctionnement) & Liberté d'autonomie détournée : choix libre de trahir la mission. L'occasion n'est pas une obligation. & Pénale (Code pénal, art. 134 et s.), administrative (révocation) et morale (trahison de la confiance). & Emprisonnement, amende, restitution, confiscation des biens ; révocation. \\
\hline
Conducteur ivre qui cause un accident mortel (axe Yaoundé-Douala) & Liberté anticipée : choix libre de boire puis de conduire ; altération du discernement volontairement acceptée. & Pénale (homicide involontaire, circonstance aggravante) et civile (réparation aux victimes). & Peine de prison, retrait du permis, dommages-intérêts importants. \\
\hline
Auteur et relais d'une fausse nouvelle WhatsApp conduisant à un lynchage & Liberté d'expression abusée. Auteur : intention de nuire ou insouciance. Relais : négligence (devoir de vérification). & Pénale (cybercriminalité, diffamation, complicité) et morale (atteinte à la dignité). & Amende, emprisonnement (auteur et relais actifs) ; réparation à la victime ; éducation aux médias. \\
\hline
\textbf{Votre cas 5 :} \emph{(grève sauvage, pollution, harcèlement...)} & \emph{À compléter} & \emph{À compléter} & \emph{À compléter} \\
\hline
\end{tabularx}
\end{center}
\legende{Tableau d'analyse de cas (TP guidé) : à compléter par l'élève pour le 5\textsuperscript{e} cas. Colonnes : acte, liberté en jeu, type de responsabilité, sanction ou réparation.}
\end{popfigure}

\courssub{Illustration quantitative : La part de la responsabilité individuelle dans les accidents de la route}

Le cas du conducteur ivre soulève une question empirique : quelle est la part de la responsabilité individuelle (le comportement) par rapport aux facteurs externes (route, véhicule) dans les accidents au Cameroun ? Le diagramme circulaire ci-dessous présente une répartition \emph{indicative et simplifiée}, cohérente avec les constats réguliers des services de la Gendarmerie Nationale et du Ministère des Transports : la grande majorité des accidents mortels est liée au comportement des usagers.

\begin{popfigure}
\begin{center}
\begin{tikzpicture}[scale=1]
% Camembert : 35% vitesse, 25% alcool, 20% véhicule, 12% chaussée, 8% autres
% Angles cumulés (départ 90°, sens horaire) : 35% -> 126° ; 25% -> 90° ; 20% -> 72° ; 12% -> 43.2° ; 8% -> 28.8°
\fill[poppink!80] (0,0) -- (90:3) arc (90:-36:3) -- cycle;
\fill[poporange!85] (0,0) -- (-36:3) arc (-36:-126:3) -- cycle;
\fill[popgold!85] (0,0) -- (-126:3) arc (-126:-198:3) -- cycle;
\fill[popgreen!80] (0,0) -- (-198:3) arc (-198:-241.2:3) -- cycle;
\fill[popmuted!70] (0,0) -- (-241.2:3) arc (-241.2:-270:3) -- cycle;
\draw[popdark, thick] (0,0) circle (3);
% Pourcentages
\node[white, font=\bfseries\small] at (27:1.8) {35 \%};
\node[white, font=\bfseries\small] at (-81:1.8) {25 \%};
\node[popdark, font=\bfseries\small] at (-162:1.8) {20 \%};
\node[white, font=\bfseries\footnotesize] at (-219.6:1.9) {12 \%};
\node[white, font=\bfseries\footnotesize] at (-255.6:1.9) {8 \%};
% Légende
\node[anchor=west, font=\small] at (3.6,2.2) {\textcolor{poppink!80}{\rule{0.35cm}{0.35cm}} Vitesse excessive : 35 \%};
\node[anchor=west, font=\small] at (3.6,1.4) {\textcolor{poporange!85}{\rule{0.35cm}{0.35cm}} Alcool / stupéfiants : 25 \%};
\node[anchor=west, font=\small] at (3.6,0.6) {\textcolor{popgold!85}{\rule{0.35cm}{0.35cm}} État du véhicule (freins, pneus) : 20 \%};
\node[anchor=west, font=\small] at (3.6,-0.2) {\textcolor{popgreen!80}{\rule{0.35cm}{0.35cm}} Chaussée / signalisation : 12 \%};
\node[anchor=west, font=\small] at (3.6,-1.0) {\textcolor{popmuted!70}{\rule{0.35cm}{0.35cm}} Autres (météo, piétons) : 8 \%};
\end{tikzpicture}
\end{center}
\legende{Répartition indicative et simplifiée des causes principales d'accidents mortels de la route au Cameroun (ordres de grandeur cohérents avec les bilans de la Gendarmerie Nationale et du Ministère des Transports ; données pédagogiques synthétiques). \textbf{Interprétation :} 60 \% des causes (vitesse 35 \% + alcool 25 \%) relèvent \textbf{directement de la liberté et de la responsabilité du conducteur}. L'état du véhicule (20 \%) relève de la responsabilité du propriétaire ou du gestionnaire. Seuls 20 \% environ (chaussée 12 \% + autres 8 \%) échappent largement à la volonté immédiate de l'agent. Vérification : $35 + 25 + 20 + 12 + 8 = 100$ \%. La responsabilité individuelle est donc le levier principal de la prévention routière.}
\end{popfigure}

\courssub{Conclusion des applications : Être libre, c'est accepter de répondre de ce qu'on fait de sa liberté}

Ces quatre études de cas, le débat et le TP convergent vers une même vérité philosophique et juridique : \textbf{la responsabilité n'est pas une charge qu'on subit, c'est la rançon de la liberté.}

\begin{aretenir}[Synthèse finale]
\begin{itemize}
    \item \textbf{L'élève} apprend que tricher, c'est se mentir à soi-même : on ne construit pas son avenir sur un mensonge.
    \item \textbf{Le fonctionnaire} comprend que gérer le bien public n'est pas un droit mais un devoir ; la corruption n'est pas une « astuce », c'est un vol fait à la nation.
    \item \textbf{Le conducteur} réalise que le volant peut devenir une arme ; boire avant de conduire, c'est choisir de devenir dangereux.
    \item \textbf{L'utilisateur de WhatsApp} découvre que le pouce qui transfère est une main qui peut tuer ; la viralité n'efface pas la responsabilité personnelle.
    \item \textbf{Le jeune face à la tradition} mesure que l'autonomie se conquiert par le dialogue et l'assomption de ses choix, non par la soumission ni par la rupture stérile.
\end{itemize}
\end{aretenir}

\noindent En définitive, comme l'écrit \textbf{Jean-Paul Sartre} dans \emph{L'Être et le Néant} (1943) : \og L'homme est condamné à être libre \fg. Condamné, parce qu'il ne s'est pas créé lui-même ; mais libre, parce qu'une fois jeté dans le monde, il est responsable de tout ce qu'il fait. Au Cameroun aujourd'hui, face aux défis de l'intégrité, de la sécurité routière, de la paix sociale et de l'émancipation de la jeunesse, cette sentence n'est pas une abstraction : c'est un programme d'action éthique et civique.

\begin{exoresolu}[Exercice type Bac : Dissertation]
\textbf{Sujet :} « La liberté est-elle la condition de la responsabilité ou son fardeau ? »

\textbf{Analyse du sujet :}
\begin{itemize}
    \item \textbf{Termes clés :} liberté (pouvoir de choisir), condition (préalable nécessaire), fardeau (poids, charge pénible), responsabilité (obligation de répondre de ses actes).
    \item \textbf{Problématique :} Si la liberté fonde la responsabilité (sans liberté, pas de responsable), elle en est aussi le fardeau (le libre choix engage l'avenir, expose au risque, à la faute, au remords). Faut-il pour autant vouloir s'en décharger ?
\end{itemize}

\textbf{Plan détaillé :}
\begin{enumerate}
    \item \textbf{La liberté, condition \emph{sine qua non} de la responsabilité (fondation).}
    \begin{itemize}
        \item Pas de responsabilité sans imputabilité (exemples : l'enfant, le malade mental, la contrainte physique).
        \item Droit pénal : principe de légalité et imputabilité personnelle (art. 23 du Code pénal camerounais).
        \item Exemple : l'élève qui triche (cas 1) est responsable \emph{parce que} libre de ne pas tricher.
    \end{itemize}
    \item \textbf{La liberté comme fardeau : l'angoisse de choisir et le poids des conséquences.}
    \begin{itemize}
        \item Sartre : l'\emph{angoisse} — je suis le seul auteur de mes valeurs.
        \item Responsabilité élargie : je réponds non seulement de mes actes, mais de ce que je \emph{fais de moi} (mon projet de vie).
        \item Exemples : le fonctionnaire (choix entre corruption et intégrité), le conducteur (choix entre le verre et le volant). Le fardeau, c'est le risque d'échec, la culpabilité, la sanction.
    \end{itemize}
    \item \textbf{Assumer le fardeau : la responsabilité comme accomplissement de la liberté (dépassement).}
    \begin{itemize}
        \item La responsabilité n'est pas une punition, c'est la \emph{dignité} de l'agent moral (Kant : l'autonomie, c'est se donner à soi-même sa propre loi).
        \item Citoyenneté active : refuser la corruption, vérifier l'information avant de la partager, négocier son orientation.
        \item La sanction juste n'est pas une vengeance : elle \emph{rétablit} l'ordre rompu et \emph{éduque} le sujet (fonction réparatrice).
    \end{itemize}
\end{enumerate}

\textbf{Conclusion :} La liberté n'est pas \emph{soit} condition \emph{soit} fardeau ; elle est \emph{les deux inséparablement}. Vouloir la liberté sans la responsabilité, c'est vouloir le pouvoir sans le devoir : c'est l'arbitraire (tyrannie ou infantilisme). Être un homme libre, et un citoyen camerounais responsable, c'est \emph{porter} le fardeau de ses choix avec lucidité et courage.
\tcblower
\textbf{Corrigé type (structure attendue) :}
L'introduction doit accrocher le lecteur (citation de Sartre, ou sagesse populaire : « Celui qui porte l'enfant sur le dos en sent seul le poids »), définir les termes, dégager la problématique et annoncer le plan. Le développement suit le plan dialectique (thèse / antithèse / synthèse) en s'appuyant \textbf{systématiquement} sur les quatre études de cas du cours et sur les références au programme (Kant, Sartre, le Code pénal camerounais, la Déclaration universelle des droits de l'homme). La conclusion répond nettement à la question posée et peut s'ouvrir sur l'éthique contemporaine de la responsabilité : être responsable, c'est aussi répondre de l'avenir (environnement, jeunesse, paix sociale).
\end{exoresolu}

\begin{exercice}[Entraînement : Commentaire de texte]
\textbf{Texte :} \og L'homme n'est pas né pour être libre, il est né pour être responsable. \fg{} (formule attribuée à la réflexion existentialiste sur la condition humaine.)

\textbf{Consigne :} Expliquez et discutez cette thèse en vous appuyant sur les notions de liberté, de responsabilité et de condition humaine, et sur les exemples camerounais étudiés dans cette leçon.
\end{exercice}

\begin{corrige}[Exercice n°1]
\textbf{Explication :} Cette formule inverse le lieu commun. D'ordinaire, on dit avec Rousseau : « L'homme est né libre, et partout il est dans les fers ». Ici, au contraire, la liberté n'est pas un don de la nature, mais une \emph{conquête} accomplie par l'assomption de la responsabilité. \textbf{Discussion :} 1) \emph{Vérité de la thèse :} le nourrisson n'est pas libre (dépendance totale) ; il devient libre en devenant progressivement responsable (autonomie). 2) \emph{Limite :} la liberté est aussi une \emph{donnée ontologique} (la capacité de choisir) avant d'être un acquis ; sans elle, aucune responsabilité ne serait même concevable. 3) \emph{Exemples :} le fonctionnaire corrompu \emph{refuse} sa responsabilité et perd sa liberté (dépendance à l'argent, peur d'être pris) ; l'élève honnête \emph{assume} sa responsabilité et gagne sa liberté (diplôme mérité, estime de soi). \textbf{Conclusion :} liberté et responsabilité ne s'opposent pas : elles se fondent mutuellement, et c'est en répondant de ses actes que l'homme s'accomplit comme être libre.
\end{corrige}

\coursec{Méthode de la dissertation et synthèse}

Après avoir exploré les applications concrètes de la liberté et de la responsabilité dans le quotidien camerounais — du choix professionnel au respect du code de la route, en passant par l'engagement associatif —, nous disposons maintenant de la matière vivante pour affronter l'examen. Cette section ne se contente pas de réciter une méthode : elle vous fait \emph{vivre} la construction d'une dissertation de A à Z, sur un sujet classique du Baccalauréat technique, pour que le jour de l'épreuve, chaque geste intellectuel soit un réflexe maîtrisé.

\courssub{Rappel méthodologique : les quatre temps de la dissertation}

Une dissertation philosophique n'est pas un exposé de connaissances, mais une \emph{réponse argumentée à un problème}. Elle obéit à une architecture rigoureuse en quatre temps, que tout correcteur attend de retrouver explicitement.

\begin{methode}[Les quatre temps de la dissertation]
\begin{enumerate}
    \item \textbf{Analyser le sujet} : identifier les termes clés, leurs sens usuels et philosophiques, les présupposés, la tension ou la question cachée.
    \item \textbf{Problématiser} : transformer le sujet en une question précise qui met en tension deux thèses opposées. La problématique n'est pas le sujet reformulé, c'est le \emph{moteur} de la démonstration.
    \item \textbf{Construire le plan dialectique} : organiser la réponse en trois moments logiques (Thèse / Antithèse / Synthèse). Chaque partie comporte deux à trois arguments, chacun illustré par un exemple précis ou une référence d'auteur.
    \item \textbf{Rédiger} : Introduction (accroche, définitions, problématique, annonce du plan), Développement (paragraphes structurés : thèse - argument - exemple - lien), Conclusion (réponse synthétique, ouverture).
\end{enumerate}
\end{methode}

\begin{attention}[Erreur fréquente au Bac : le « catalogue d'auteurs »]
Ne citez pas Descartes, Kant ou Sartre pour « faire savant ». Une citation sans analyse ne vaut rien. Chaque référence doit \emph{servir} votre argumentation : « \emph{Kant montre que... donc...} ». Le correcteur préfère une dissertation sans grands noms mais bien problématisée, structurée et illustrée d'exemples concrets (comme ceux vus en section 5), à un défilé d'auteurs mal compris.
\end{attention}

\courssub{Entraînement guidé : analyse du sujet « Peut-on être responsable sans être libre ? »}

Appliquons la méthode pas à pas.

\textbf{1. Analyse des termes}\par
\begin{itemize}
    \item \cle{Peut-on} : question de possibilité (physique, métaphysique, juridique, morale), pas seulement de fait.
    \item \cle{Être responsable} : devoir répondre de ses actes devant autrui, la loi, sa conscience ; assumer les conséquences.
    \item \cle{Sans être libre} : absence de liberté (contrainte physique, déterminisme psychologique/social, ignorance, aliénation).
\end{itemize}

\textbf{2. Présupposés et tension}\par
Le sujet suppose habituellement que responsabilité \emph{implique} liberté (pas de sanction juste sans agent libre). Il invite à tester cette évidence : y a-t-il des cas où l'on tient quelqu'un pour responsable \emph{malgré} un déterminisme ? La responsabilité peut-elle exister \emph{autrement} que comme corollaire de la liberté métaphysique ?

\textbf{3. Problématique}\par
\cle{Problématique} : \emph{La responsabilité est-elle l'apanage exclusif de l'agent libre, ou peut-elle s'étendre à des êtres déterminés, dès lors qu'on la conçoit non comme imputation rétrospective mais comme exigence éducative, sociale ou existentielle ?}

\courssub{Construction collective du plan dialectique}

Voici le plan type, validé par les jurys, appliqué à notre sujet. Le schéma ci-dessous en donne la logique visuelle.

\begin{popfigure}
\begin{tikzpicture}[
    node distance=1.8cm and 2.5cm,
    box/.style={rectangle, draw=popblue, thick, rounded corners, fill=popblue!5, text width=4.5cm, align=center, minimum height=2.2cm, font=\small},
    arrow/.style={-Stealth, thick, popdark},
    labelbox/.style={rectangle, draw=poporange, thick, rounded corners, fill=poporange!10, text width=5cm, align=center, minimum height=1.2cm, font=\small\itshape}
]
% Nodes
\node[box, align=center] (thesis){\textbf{I. La responsabilité suppose la liberté (Thèse)}\\\textit{Argument 1 :} Condition de l'imputation morale (Kant : autonomie).\\\textit{Argument 2 :} Fondement du droit et de la sanction (on ne punit pas un automate).\\\textit{Argument 3 :} Expérience vécue du « j'aurais pu faire autrement ».};
\node[box, right=of thesis, align=center] (antithesis){\textbf{II. Mais des déterminismes limitent la liberté (Antithèse)}\\\textit{Argument 1 :} Déterminismes physique, psychologique, social (Spinoza, Marx, Freud).\\\textit{Argument 2 :} Responsabilité sans liberté métaphysique : responsabilité civile, collective, éducative.\\\textit{Argument 3 :} L'illusion du libre arbitre (neurosciences, conditionnement).};
\node[box, below=2cm of thesis.south east, xshift=1.5cm, align=center] (synthesis){\textbf{III. La liberté comme conquête et la responsabilité comme assomption (Synthèse)}\\\textit{Argument 1 :} Liberté non comme don mais comme tâche (Sartre : « condamné à être libre »).\\\textit{Argument 2 :} Responsabilité comme \emph{assomption} de son passé et de son avenir (Ricoeur).\\\textit{Argument 3 :} Exigence d'une société qui \emph{forme} des sujets responsables (éducation, droit réparateur).};
% Arrows
\draw[arrow] (thesis) -- (antithesis) node[midway, above, labelbox] {Tension : \emph{Liberté condition nécessaire ?}};
\draw[arrow] (thesis.south) -- (synthesis.north west);
\draw[arrow] (antithesis.south) -- (synthesis.north east);
% Title
\node[above=0.5cm of thesis.north east, xshift=1.5cm, font=\bfseries\large, text=popdark] {Plan dialectique : « Peut-on être responsable sans être libre ? »};
\end{tikzpicture}
\legende{Schéma du plan dialectique en trois temps : la thèse établit le lien nécessaire, l'antithèse le complexifie par le déterminisme, la synthèse le dépasse par l'idée de liberté conquise et de responsabilité assumée.}
\end{popfigure}

\begin{definition}[Plan dialectique (rappel)]
Structure en trois moments : \textbf{Thèse} (réponse affirmative argumentée), \textbf{Antithèse} (objections radicales, réponse négative argumentée), \textbf{Synthèse} (dépassement de l'opposition par un concept plus riche qui intègre le vrai des deux premiers moments). Ce n'est pas un compromis, c'est une \emph{résolution} du problème.
\end{definition}

\courssub{Rédaction guidée : introduction complète}

L'introduction est la « vitrine » de votre copie. Elle doit être rédigée \emph{au brouillon} puis recopiée proprement. Voici un modèle pour notre sujet.

\begin{methode}[Rédiger une introduction de dissertation en 4 étapes]
\begin{enumerate}
    \item \textbf{Amener le sujet (accroche)} : citation, fait d'actualité, paradoxe, référence littéraire ou historique. \emph{Évitez : « Depuis la nuit des temps... »}
    \item \textbf{Définir les termes} : sens philosophique précis de « responsabilité » (imputabilité, charge) et « liberté » (pouvoir de choisir, autonomie, absence de contrainte). Distinguer liberté métaphysique / politique / morale.
    \item \textbf{Poser la problématique} : la question précise qui guide tout le devoir. Elle doit contenir la tension.
    \item \textbf{Annoncer le plan} : « Nous montrerons d'abord que... (I), puis que... (II), pour enfin montrer que... (III). »
\end{enumerate}
\end{methode}

\begin{exoresolu}[Introduction type — Sujet : « Peut-on être responsable sans être libre ? »]
\textbf{Accroche :} À la fin de \emph{L'Étranger}, Albert Camus fait dire à Meursault : « Je n'avais qu'à souhaiter qu'il y ait beaucoup de spectateurs le jour de mon exécution et qu'ils m'accueillent avec des cris de haine. » Cette acceptation lucide de son sort, face à l'absurde, illustre le paradoxe de notre sujet : la responsabilité semble exiger la liberté, pourtant l'expérience montre des hommes déterminés que la société tient pour responsables.

\textbf{Définitions :} La \cle{responsabilité} (du latin \emph{respondere}, répondre) est l'obligation de répondre de ses actes, d'en assumer les conséquences devant autrui, la loi ou sa conscience. La \cle{liberté} est la capacité d'agir par soi-même, sans contrainte extérieure (liberté politique) ni déterminisme intérieur absolu (liberté métaphysique), c'est-à-dire le pouvoir de choisir entre plusieurs possibles.

\textbf{Problématique :} Si la responsabilité suppose traditionnellement un sujet libre et autonome, ne faut-il pas distinguer la \emph{condition métaphysique} de l'imputation (liberté de l'indifférence) de la \emph{fonction sociale et existentielle} de la responsabilité ? \emph{La responsabilité est-elle l'apanage exclusif de l'agent libre, ou peut-elle s'étendre à des êtres déterminés, dès lors qu'on la conçoit comme exigence éducative, sociale ou comme assomption de soi ?}

\textbf{Annonce du plan :} Nous montrerons d'abord que la responsabilité \emph{suppose} la liberté comme condition de possibilité de l'imputation morale et juridique (I). Nous verrons ensuite que des déterminismes (physiques, psychologiques, sociaux) semblent rendre cette liberté illusoire, sans pour autant abolir les pratiques de responsabilité (II). Enfin, nous soutiendrons que la liberté n'est pas un donné mais une \emph{conquête}, et que la responsabilité, loin de n'être qu'une sanction, est l'\emph{assomption} lucide de sa condition d'agent fini mais capable de se projeter (III).
\tcblower
\textbf{Analyse de l'introduction :} L'accroche (Camus) ancre le sujet dans la littérature. Les définitions distinguent les sens (moral/juridique vs métaphysique/politique). La problématique reformule le sujet en question philosophique tensionnée. L'annonce du plan suit exactement la structure Thèse/Antithèse/Synthèse du schéma TikZ.
\end{exoresolu}

\courssub{Rédaction et justification d'une conclusion personnelle argumentée}

La conclusion n'est pas un résumé. C'est le moment où vous \emph{tranchez} en votre nom propre, au terme du parcours dialectique.

\begin{methode}[Justifier une conclusion en 3 étapes]
\begin{enumerate}
    \item \textbf{Reprendre la problématique} : reformulez-la pour montrer qu'elle a été traitée.
    \item \textbf{Trancher en argumentant} : donnez votre réponse définitive, en vous appuyant sur le chemin parcouru (la synthèse). N'ayez pas peur du « Je pense que... » si c'est étayé.
    \item \textbf{Ouvrir sur une question nouvelle} : prolongez la réflexion vers un horizon proche (autre notion, actualité, auteur) sans y répondre.
\end{enumerate}
\end{methode}

\begin{exemplebox}[Conclusion type pour le sujet traité]
\textbf{Reprise :} Au terme de cette analyse, il apparaît que la question « Peut-on être responsable sans être libre ? » ne reçoit pas de réponse binaire.

\textbf{Trancher :} La liberté métaphysique (pouvoir absolu de choisir autrement) n'est pas une condition \emph{empirique} toujours vérifiée : nos choix sont conditionnés. Pourtant, la responsabilité ne s'effondre pas. Elle se déplace : elle n'est plus seulement \emph{imputation rétrospective} d'un acte à une cause libre, mais \emph{exigence prospective} d'assomption. Être responsable, c'est \emph{devenir} libre en reconnaissant ses déterminismes pour mieux les dépasser (Sartre, Ricoeur). La responsabilité est ainsi le \emph{moteur} de la liberté, non sa simple conséquence.

\textbf{Ouverture :} Cette conception dynamique de la responsabilité comme conquête invite à repenser la justice : ne faut-il pas passer d'une justice \emph{rétributive} (punir le libre choix) à une justice \emph{restaurative} (réparer, réinsérer, rendre capable de responsabilité) ? C'est la question du \emph{devoir} et de la \emph{justice} qui se profile à l'horizon.
\end{exemplebox}

\courssub{Bilan de la leçon : carte mentale récapitulative}

Cette carte mentale synthétise l'ensemble des concepts, auteurs, distinctions et exemples abordés dans les six sections de la leçon. Elle est votre « fiche de révision » visuelle pour le Bac.

\begin{popfigure}
\begin{tikzpicture}[
    mindmap,
    concept color=popblue!80,
    level 1 concept/.append style={level distance=4.5cm, sibling angle=90, font=\small\bfseries, text=white},
    level 2 concept/.append style={level distance=3.2cm, sibling angle=45, font=\footnotesize, text=popdark},
    level 3 concept/.append style={level distance=2.5cm, sibling angle=30, font=\scriptsize, text=popmuted},
    every node/.append style={rounded corners=2pt},
    grow cyclic
]
\node[concept] {Liberté \& Responsabilité}
    child[concept color=popgreen!70!black] {
        node[concept] {Liberté}
        child { node[concept] {Formes}
            child { node {Métaphysique (libre arbitre)} }
            child { node {Politique (droits, lois)} }
            child { node {Morale (autonomie Kant)} }
            child { node {Existentielle (Sartre)} }
        }
        child { node[concept] {Déterminismes}
            child { node {Physique / Naturel} }
            child { node {Psychologique (Freud)} }
            child { node {Social / Économique (Marx)} }
        }
        child { node[concept] {Conquête}
            child { node {Éducation / Culture} }
            child { node {Réflexivité critique} }
        }
    }
    child[concept color=poporange!80!black] {
        node[concept] {Responsabilité}
        child { node[concept] {Formes}
            child { node {Moral (conscience)} }
            child { node {Juridique (peine, réparation)} }
            child { node {Civile (dommages)} }
            child { node {Collective / Solidaire} }
            child { node {Professionnelle (déontologie)} }
        }
        child { node[concept] {Fondements}
            child { node {Liberté (condition classique)} }
            child { node {Causalité (agent causal)} }
            child { node {Vulnérabilité (Ricoeur)} }
        }
        child { node[concept] {Assomption}
            child { node {Reconnaissance du passé} }
            child { node {Engagement futur} }
        }
    }
    child[concept color=poppurple!80!black] {
        node[concept] {Lien Liberté $\leftrightarrow$ Responsabilité}
        child { node[concept] {Classique : Liberté $\rightarrow$ Responsabilité} }
        child { node[concept] {Moderne : Responsabilité $\rightarrow$ Liberté (conquête)} }
        child { node[concept] {Droit : Imputabilité (art. 122-1 Code pénal CM)} }
        child { node[concept] {Éthique : Soin / Réparation} }
    }
    child[concept color=poppink!80!black] {
        node[concept] {Exemples Camerounais}
        child { node {Choix filière / métier} }
        child { node {Code de la route / Corruption} }
        child { node {Gestion déchets / Environnement} }
        child { node {Tontine / Solidarité} }
        child { node {Stage / Contrat apprentissage} }
        child { node {Vote / Devoir civique} }
    };
\end{tikzpicture}
\legende{Carte mentale récapitulative de la leçon « Liberté et Responsabilité ». Les quatre branches principales (Liberté, Responsabilité, Lien, Exemples CM) structurent l'ensemble des savoirs et savoir-faire exigibles au Probatoire/Baccalauréat technique.}
\end{popfigure}

\courssub{Ouverture : vers la leçon suivante — Le Devoir et le Bonheur}

Notre analyse s'achève sur une tension non résolue : si la responsabilité est l'assomption lucide de notre condition, elle s'exprime souvent sous la forme d'un \textbf{devoir} (« Je dois assumer », « Je dois réparer »). Or, le devoir est fréquemment vécu comme une contrainte, un « fardeau » qui s'oppose au \textbf{bonheur} (entendu comme satisfaction désirable, plaisir ou épanouissement).

\begin{aretenir}[Transition vers l'Unité 3 : Le Devoir et le Bonheur]
\begin{itemize}
    \item La responsabilité nous a menés au seuil du \cle{devoir} : être responsable, c'est se reconnaître des obligations.
    \item Mais le devoir entre-t-il en contradiction avec le bonheur ? Kant oppose devoir et inclination ; les eudémonistes (Aristote, épicuriens) les réconcilient dans la vertu.
    \item La leçon suivante problématisera : \emph{« Le devoir est-il l'ennemi du bonheur ? »} ou \emph{« Faut-il sacrifier son bonheur pour faire son devoir ? »}.
    \item Nous y mobiliserons : Kant (impératif catégorique), Aristote (éthique à Nicomaque), les stoïciens, et des cas concrets (devoir professionnel vs vie familiale, engagement militant vs confort personnel).
\end{itemize}
\end{aretenir}

\begin{savaistu}[Au Baccalauréat : structure vs érudition]
Au Baccalauréat technique, une dissertation \textbf{bien structurée} (problématique claire, plan dialectique visible, arguments enchaînés logiquement, exemples précis — comme ceux de la section 5) \textbf{vaut mieux qu'un catalogue d'auteurs mal compris}. Les correcteurs sont instruits de valoriser la \emph{rigueur du raisonnement} et la \emph{pertinence de l'illustration} (contexte camerounais bienvenu) sur le \emph{name-dropping} philosophique. Maîtrisez la méthode (boîtes \texttt{methode} ci-dessus) : c'est votre assurance-vie le jour de l'épreuve.
\end{savaistu}

\begin{exemplebox}[Sujets tombés au Baccalauréat Technique (années récentes)]
\begin{itemize}
    \item « La liberté est-elle un pouvoir ou un fardeau ? » (Session 2022) — \emph{Travailler la distinction liberté-puissance / responsabilité-poids.}
    \item « Faut-il punir pour être juste ? » (Session 2021) — \emph{Articuler responsabilité, sanction, justice réparatrice vs rétributive.}
    \item « L'homme est-il responsable de son avenir ? » (Session 2020) — \emph{Projet, liberté, déterminisme, responsabilité prospective.}
    \item « Peut-on être libre sans être responsable ? » (Session 2019) — \emph{Inverse de notre sujet d'entraînement ; même structure dialectique.}
\end{itemize}
\emph{Conseil :} Entraînez-vous à faire l'analyse de termes, la problématique et le plan pour chacun de ces sujets en 15 minutes chrono.
\end{exemplebox}

\begin{exercice}[Entraînement final — À faire au brouillon]
Traitez l'un des deux sujets suivants en appliquant \textbf{strictement} la méthode des 4 temps (analyse, problématique, plan détaillé avec arguments/exemples, introduction rédigée, conclusion rédigée). Temps imparti : 1h30 (conditions d'examen).
\begin{enumerate}
    \item « La responsabilité est-elle la limite de la liberté ? »
    \item « Obéir, est-ce renoncer à sa liberté ? »
\end{enumerate}
\end{exercice}

\begin{corrige}[Exercice 1 — Pistes de correction]
\textbf{Analyse :} « Limite » = borne, frontière, mais aussi condition de possibilité (le cadre qui rend le jeu possible). Tension : la responsabilité borne-t-elle la liberté (frein) ou la constitue-t-elle (cadre) ?
\textbf{Problématique :} La responsabilité est-elle une contrainte extérieure qui bride la liberté, ou la forme même par laquelle la liberté devient effective et humaine ?
\textbf{Plan :} I. Responsabilité comme limite extérieure (contrainte sociale, loi, sanction). II. Responsabilité comme limite intérieure (conscience morale, autonomie kantienne : se donner sa propre loi). III. Responsabilité comme \emph{condition} de la liberté réelle : sans responsabilité, liberté = caprice/licence ; la responsabilité structure l'espace de la liberté (droit, confiance, projet).
\textbf{Exemples CM :} Code de la route (limite extérieure) ; engagement tontine (limite intérieure/autonomie) ; liberté d'entreprendre encadrée par le droit des affaires (condition).
\end{corrige}

\begin{corrige}[Exercice 2 — Pistes de correction]
\textbf{Analyse :} « Obéir » = se soumettre à un ordre/loi. « Renoncer » = abandon volontaire. Tension : obéissance = hétéronomie (loi venue d'autrui) vs autonomie (loi de soi).
\textbf{Problématique :} L'obéissance à une loi extérieure est-elle nécessairement aliénation, ou peut-elle être l'acte libre d'un sujet rationnel qui reconnaît la légitimité de la norme ?
\textbf{Plan :} I. Obéissance comme renoncement à sa liberté (contrainte, soumission, minorité kantienne). II. Obéissance comme exercice de liberté (contrat social Rousseau, loi morale Kant, règles du jeu/technique). III. La distinction cruciale : \emph{obéissance servile} (peur, contrainte) vs \emph{obéissance libre} (reconnaissance de la raison de la règle). Liberté = capacité de se donner des lois justes.
\textbf{Exemples CM :} Respect code de la route (obéissance libre rationnelle) ; corruption / pots-de-vin (fausse obéissance, perte responsabilité) ; apprentissage métier (obéissance au maître pour conquérir autonomie future).
\end{corrige}

\newpage

\coursec{Activité d'intégration}

\coursec{Liberté et Responsabilité}

\courssub{Introduction : Problématisation et définitions}
Tout être humain se vit comme libre : je lève le bras, je choisis mes études, je décide d'aider ou non. Pourtant, la science, la sociologie et le droit nous rappellent sans cesse que nous sommes déterminés : par nos gènes, notre milieu social, nos pulsions, la loi. \textbf{Comment concilier l'expérience intérieure de la liberté avec la nécessité objective des causes ?} Et si la liberté n'est qu'une illusion, sur quoi fonder la \cle{responsabilité} — cette obligation de répondre de ses actes devant autrui, la conscience et la loi ?

\begin{definition}[Liberté]
Puissance d'agir ou de ne pas agir selon sa propre détermination. Distinction cruciale :
\begin{itemize}
    \item \textbf{Libre arbitre (liberté d'indifférence) :} Choix sans cause déterminante (Descartes). Fondement de la responsabilité absolue.
    \item \textbf{Liberté comme autonomie (Spinoza, Kant) :} Agir selon sa raison, se donner sa propre loi (loi morale). \cle{Autonomie} $\neq$ indépendance (absence de contrainte extérieure).
    \item \textbf{Liberté conditionnelle / relative :} L'homme est « condamné à être libre » (Sartre) mais toujours \textit{situé} (contexte social, historique, corporel). La liberté est un \textbf{devenir} (Merleau-Ponty).
\end{itemize}
\end{definition}

\begin{definition}[Responsabilité]
Obligation de répondre de ses actes devant autrui (conscience, société, loi). Trois dimensions :
\begin{itemize}
    \item \textbf{Morale :} Imputabilité à un sujet conscient et libre. Exige \textit{liberté} + \textit{connaissance} + \textit{volonté}. Sans liberté, pas de responsabilité morale (seule une responsabilité causale subsiste).
    \item \textbf{Civile :} Réparation du dommage (Art. 1382 Code Civil camerounais : « Tout fait quelconque de l'homme, qui cause à autrui un dommage, oblige celui par la faute duquel il est arrivé à le réparer »). Peut exister sans faute morale (responsabilité sans faute : fait des choses, risque créé). Mais faute = liberté + négligence/imprudence.
    \item \textbf{Pénale :} Sanction d'une infraction à la loi. Principe : \textit{Nulle peine sans loi}, \textit{personnalité de la peine}. Exige l'élément moral (intention/dol ou imprudence) $\rightarrow$ suppose une liberté minimale.
\end{itemize}
\end{definition}

\courssub{Les formes de la liberté}
La liberté n'est pas un bloc monolithique ; elle se décline selon le rapport aux déterminismes.

\begin{propriete}[Typologie des libertés (Repères pour le Bac)]
\begin{itemize}
    \item \textbf{Liberté physique / d'indifférence :} Absence de contrainte extérieure matérielle (chaînes, enfermement, force majeure). Condition nécessaire mais non suffisante.
    \item \textbf{Liberté psychologique / de choix :} Capacité à délibérer entre plusieurs motifs sans être esclave de ses passions ou pulsions (Spinoza : \textit{conatus} rationnel).
    \item \textbf{Liberté morale / autonomie (Kant) :} Agir par devoir, selon la loi que la raison se prescrit à elle-même (Impératif catégorique). C'est le fondement de la dignité humaine.
    \item \textbf{Liberté politique / citoyenne :} Garantie par l'État de droit (liberté d'expression, de mouvement, de vote). Cadre institutionnel de l'exercice des autres libertés.
    \item \textbf{Liberté situationnelle (Sartre, Merleau-Ponty) :} « Condamnés à être libres » dans une \textit{facticité} (né à Douala en 2007, pauvre, électrotechnicien). La liberté est l'interprétation et le dépassement de cette situation.
\end{itemize}
\end{propriete}

\courssub{La responsabilité : définition et formes}
La responsabilité est le corrélat nécessaire de la liberté : \textbf{pas de responsabilité sans liberté}, mais \textbf{la responsabilité révèle l'étendue de la liberté effective}.

\begin{propriete}[Les trois registres de la responsabilité]
\begin{itemize}
    \item \textbf{Responsabilité morale :} Tribunal de la conscience. Jugement sur l'intention (Kant : \textit{bonne volonté}) et non seulement le résultat. Exige \textit{imputabilité} (sujet conscient + libre). \textbf{Personnelle et non transférable.}
    \item \textbf{Responsabilité civile :} Tribunal des intérêts. Objectif : \textit{réparation} du dommage (Art. 1382 C. civ.). Faute = violation d'une obligation de prudence (liberté + négligence). \textbf{Souvent solidaire} (employeur répond pour préposé Art. 1384 al. 4 ; gardien de la chose Art. 1384 al. 1). Peut être \textit{sans faute} (risque créé, chose).
    \item \textbf{Responsabilité pénale :} Tribunal de la loi. Objectif : \textit{sanction} (peine) + protection société. Principe de \textit{légalité} (Art. 1 Cp : pas de peine sans texte). Exige \textit{élément moral} (intentionnel = dol ; non-intentionnel = imprudence/négligence Art. 2 Cp). \textbf{Personnalité de la peine} (Art. 3 Cp) : peine ne frappe que son auteur. \textbf{Seuil d'âge crucial au Cameroun :} < 13 ans = irresponsabilité pénale absolue (Art. 80 Cp) ; 13-18 ans = responsabilité atténuée (mesures éducatives prioritaires, Ordonnance 62/OF/18).
\end{itemize}
\end{propriete}

\courssub{La liberté, fondement de la responsabilité}
Le lien est structurel : \textbf{responsabilité $\propto$ liberté effective}. Mais comment mesurer cette liberté effective face aux contraintes ?

\begin{propriete}[Critères d'imputabilité (Aristote, \textit{Éthique à Nicomaque}, III, 1-5)]
L'acte est \textbf{volontaire} (source de responsabilité) si le principe du mouvement est \textit{dans l'agent} qui connaît les circonstances particulières.
L'acte est \textbf{involontaire} (exonère) s'il y a \textbf{contrainte extérieure} (vent, hommes qui portent) \textbf{ET} que l'agent n'y contribue en rien.
L'acte est \textbf{non-volontaire} (regrettable mais pas blâmable) s'il y a contrainte mais que l'agent a \textit{choisi} le moindre mal (ex: jeter la cargaison pour sauver le navire).
\emph{Application :} La contrainte doit être \textit{irrésistible} pour abolir la liberté. La peur, la faim, la pression hiérarchique sont-elles des contraintes irrésistibles ? \textbf{Non, en droit positif (sauf état de nécessité Art. 21 Cp), mais elles atténuent la responsabilité morale.}
\end{propriete}

\begin{attention}[Pièges fréquents à l'examen]
\begin{itemize}
    \item Confondre \textbf{cause} (explication déterministe : milieu social, cerveau, gènes) et \textbf{raison} (justification normative : pourquoi j'ai agi ainsi). La cause nie la liberté ; la raison la suppose.
    \item Confondre \textbf{responsabilité civile} (indemnisation, souvent solidaire : employeur paie pour salarié) et \textbf{responsabilité pénale/morale} (personnelle, non transférable).
    \item Oublier le \textbf{seuil d'âge} : < 13 ans = irresponsabilité pénale absolue (mesures de protection seulement). 13-18 ans = responsabilité pénale atténuée (mesures éducatives prioritaires, peine privative de liberté exceptionnelle).
    \item Dire « il n'avait pas le choix » sans analyser : \textit{avait-il la conscience du choix ?} \textit{Pouvait-il refuser ?} (Coût du refus : licenciement, violence fraternelle, échec scolaire).
\end{itemize}
\end{attention}

\courssub{Applications concrètes : Simulation de tribunal (Cas camerounais)}
\emph{Durée : 2 séances. Mise en situation : Tribunal correctionnel de la Cour d'Appel du Wouri (Douala), audience spéciale « Jeunesse et Citoyenneté ».}

\courssub{Objectifs de l'activité}
\begin{itemize}
    \item Mobiliser les concepts de \cle{liberté} (libre arbitre, autonomie, conditionnelle) et de \cle{responsabilité} (morale, civile, pénale) pour analyser des cas concrets.
    \item Évaluer le degré de liberté effective d'un agent face à des contraintes (physiques, psychologiques, sociales, hiérarchiques).
    \item Distinguer les types de responsabilité engagés selon la nature de l'acte et la qualité de l'auteur.
    \item Construire un raisonnement juridique et philosophique structuré aboutissant à un verdict argumenté.
    \item Développer l'esprit critique, l'argumentation orale et la rédaction synthétique dans un cadre simulé.
\end{itemize}

\courssub{Situation : Trois dossiers à la barre}
La composition du tribunal pour cette simulation est la suivante : un président (enseignant ou élève désigné), un procureur (groupe de 3-4 élèves), des avocats de la défense (groupe de 3-4 élèves par dossier), des avocats des parties civiles (groupe de 2-3 élèves par dossier), et un jury populaire (le reste de la classe) qui délibérera. Chaque groupe dispose du dossier écrit ci-après.

\bigskip
\noindent\textbf{DOSSIER N\textsuperscript{o} 1 : L'affaire « Baccalauréat Blanc »}
\medskip
\noindent\textbf{Prévenu :} \textsc{Ngonzo} Kevin, 18 ans, élève en Terminale T (Série F4 - Électrotechnique) au Lycée Technique de New-Bell.
\medskip
\noindent\textbf{Faits :} Le 12 mars 2025, lors de l'épreuve de Mathématiques du Baccalauréat blanc (coefficient 6), le surveillant principal, M. \textsc{Ebongue}, surprend Kevin en train de consulter des formules manuscrites dissimulées dans sa calculatrice programmable (modèle autorisé mais mémoire non effacée). Kevin est immédiatement exclu de la salle. Procès-verbal dressé à 10h15.
\medskip
\noindent\textbf{Déclarations du prévenu :} « Monsieur le Président, tout le monde le fait. Les formules de transformées de Laplace et de Maxwell, on ne peut pas tout retenir par cœur. Mon camarade \textsc{Tchinda} m'a passé la calculatrice en me disant : "C'est ça qui sort souvent". J'ai pas réfléchi, j'ai mis dans ma poche. C'est le système qui nous pousse à la triche : le programme est trop chargé, les coefficients trop élevés, si on rate, on perd l'année. Je suis pas un délinquant, je veux juste mon Bac pour entrer à l'ENSET. »
\medskip
\noindent\textbf{Contexte :} Kevin n'a jamais eu de sanction disciplinaire. Moyenne annuelle : 11,5/20. Parents convoqués : père chauffeur de taxi, mère vendeuse au marché Congo. Frais de scolarité payés avec difficulté.

\bigskip
\noindent\textbf{DOSSIER N\textsuperscript{o} 2 : L'affaire « Carrefour Ndogbati »}
\medskip
\noindent\textbf{Prévenu :} \textsc{Mballa} Jean-Pierre, 38 ans, chauffeur de bus de transport en commun (ligne 14 : Douala-Nord -- Bonabéri), employé de la société \textit{Trans-Union SARL}.
\medskip
\noindent\textbf{Faits :} Le 5 avril 2025, vers 06h45, au carrefour Ndogbati (route non bitumée, ornières profondes, absence de signalisation), le bus immatriculé LT 456 XY perd le contrôle dans un virage serré, percute un motocycliste (\textsc{Njock} Armand, 24 ans, livreur \textit{Jumia}, décédé sur le coup) et deux piétons (blessés graves). Test d'alcoolémie du chauffeur : 0,0 g/L. Vitesse estimée par expertise : 45 km/h en zone limitée à 30 km/h (zone scolaire non signalée). Freins arrière défectueux (rapport expertise technique).
\medskip
\noindent\textbf{Déclarations du prévenu :} « Patron m'a dit : "Jean-Pierre, ce bus il doit faire 4 allers-retours avant 10h, sinon c'est ton remplacement". La route là, c'est la catastrophe. J'ai signalé les freins ya trois semaines au chef d'atelier, il a dit "on a pas les pièces". Ce matin-là, j'ai freiné, la pédale est allée au plancher. J'ai essayé d'éviter le moto, mais l'ornière m'a piégé. Je suis pas un assassin, je suis une victime comme les autres. »
\medskip
\noindent\textbf{Contexte :} Mballa père de 3 enfants, salaire de base 60 000 FCFA + primes au trajet (200 FCFA/trajet). Pression quotidienne pour atteindre le quota de 12 trajets/jour. Société \textit{Trans-Union} mise en cause pour défaut d'entretien du parc automobile (12 bus sur 20 en état précaire selon rapport inspection travail 2024).

\bigskip
\noindent\textbf{DOSSIER N\textsuperscript{o} 3 : L'affaire « Quartier Makepe »}
\medskip
\noindent\textbf{Prévenu :} \textsc{Ekobo} Stéphane, 15 ans (né le 14/02/2010), scolarisé en classe de 4\textsuperscript{e} au CES de Makepe, résidant au quartier Makepe Missoké.
\medskip
\noindent\textbf{Faits :} Le 20 mai 2025, vers 20h30 (couvre-feu mineurs non respecté), Stéphane, muni d'un couteau de cuisine (lame 15 cm), agresse \textsc{Mme Fouda} (52 ans, revendeuse de beignets) à la sortie de son domicile. Menace à l'arme blanche, vol du sac à main (contenant 45 000 FCFA recette du jour + téléphone). Mme Fouda blessée au bras (ITT 10 jours). Stéphane interpellé 20 min plus tard par patrouille police, sac et couteau retrouvés sur lui. Aveux immédiats.
\medskip
\noindent\textbf{Déclarations du prévenu :} « C'est mon grand frère \textsc{Bebey} (19 ans, en fuite) qui m'a dit : "Va chercher l'argent de la beignets-femme, on a faim, papa a pas donné chop depuis 3 jours". J'ai pris le couteau pour faire peur seulement. J'ai pas voulu la blesser, elle a résisté, le couteau a glissé. J'ai peur de Bebey, s'il dit fait, tu fais. »
\medskip
\noindent\textbf{Contexte :} Famille monoparentale (père décédé 2022), mère malade (insuffisance rénale, dialyse 2 fois/semaine, coût 35 000 FCFA/semaine non pris en charge). 4 enfants à charge. Stéphane déscolarisé de fait (absentéisme 60\%). Frère aîné \textsc{Bebey} connu des services police pour vols à l'arraché. Procédure éducative ouverte par le Juge des Enfants (Tribunal pour Enfants de Douala).

\courssub{Consignes de travail (Séance 1 : Analyse \& Préparation ; Séance 2 : Plaidoiries \& Verdicts)}
\begin{enumerate}
    \item \textbf{Analyse des faits (10 min) :} Pour chaque dossier, lister les éléments objectifs (acte matériel, circonstances, auteur, victime) et subjectifs (mobile, intention, discours du prévenu).
    \item \textbf{Évaluation de la liberté (15 min) :} Qualifier le type de contrainte subie par l'auteur :
    \begin{itemize}
        \item Contrainte \textbf{physique} (force majeure, impossibilité matérielle) ?
        \item Contrainte \textbf{psychologique/morale} (peur, pression hiérarchique, endoctrinement, habitude) ?
        \item Contrainte \textbf{sociale/structurelle} (pauvreté, système éducatif, insécurité routière) ?
    \end{itemize}
    Déterminer si la \cle{conscience} de l'acte et de sa valeur morale/juridique était pleine, altérée ou abolie. Tenir compte de l'\textbf{âge} (majorité pénale à 18 ans au Cameroun, minorité relative 13-18 ans, irresponsabilité pénale < 13 ans -- Art. 80 Code Pénal).
    \item \textbf{Qualification de la responsabilité (10 min) :} Identifier le(s) type(s) de responsabilité engagée(s) :
    \begin{itemize}
        \item \textbf{Responsabilité morale :} Jugement de valeur sur l'intention, la conscience, le devoir éthique (Kant : \textit{devoir pour devoir} ; Aristote : \textit{volontaire / involontaire / non-volontaire}).
        \item \textbf{Responsabilité civile :} Obligation de réparer le dommage (Art. 1382 Code Civil camerounais). Responsabilité du fait personnel, du fait d'autrui (employeur/complice), du fait des choses.
        \item \textbf{Responsabilité pénale :} Infraction à la loi pénale (triche : fraude aux examens Art. 139 Cp ; accident : homicide/involontaire blessures involontaires Art. 269-270 Cp ; vol avec violence Art. 318 Cp). Recherche de l'élément moral (intentionnel / non-intentionnel) et matériel.
    \end{itemize}
    \item \textbf{Rédaction du verdict (15 min) :} Chaque groupe (Procureur / Défense / Parties civiles) rédige un \textbf{verdict argumenté (10-15 lignes)} pour son dossier, structuré ainsi :
    \begin{enumerate}
        \item Qualification juridique des faits.
        \item Appréciation de la liberté (contraintes / conscience / âge).
        \item Type(s) de responsabilité retenue(s) et fondement.
        \item Sanction / Mesure proposée (peine, réparation, mesure éducative, relaxe) et justification philosophique (ex: finalité de la peine : réparation, dissuasion, réinsertion).
    \end{enumerate}
    \item \textbf{Débat final et méta-réflexion (20 min) :} Lecture des verdicts. Confrontation des arguments. Discussion : \textit{La responsabilité est-elle proportionnelle à la liberté effective ? La contrainte sociale excuse-t-elle moralement ? Le droit positif (Code Pénal, Code Civil) coïncide-t-il avec la responsabilité morale ?}
\end{enumerate}

\courssub{Méthode de la dissertation et synthèse}

\begin{methode}[Rédiger une dissertation philosophique (Sujet type : « La responsabilité suppose-t-elle une liberté totale ? »)]
\begin{enumerate}
    \item \textbf{Analyser le sujet :} Définir « responsabilité » (répondre de ses actes : morale, civile, pénale) et « liberté totale » (absence de toute contrainte, libre arbitre radical). Repérer le verbe « suppose » (condition nécessaire ? suffisante ?).
    \item \textbf{Problématiser :} Si responsabilité $\rightarrow$ liberté, alors liberté totale = responsabilité totale ? Mais l'expérience montre des libertés entravées (contraintes). Faut-il une liberté \textit{parfaite} ou \textit{suffisante} (seuil d'imputabilité) ?
    \item \textbf{Construire le plan (Dialectique) :}
    \begin{itemize}
        \item \textbf{Thèse :} La responsabilité \textit{exige} une liberté effective (Aristote : acte volontaire = principe en l'agent + connaissance). Sans liberté (contrainte irrésistible), pas d'imputabilité morale. Le droit le confirme (irresponsabilité < 13 ans, cause d'irresponsabilité Art. 21 Cp : contrainte force majeure).
        \item \textbf{Antithèse :} La liberté \textit{totale} est un mythe (Spinoza : « les hommes se croient libres... »). Nous sommes toujours déterminés (social, psychique, biologique). Pourtant, le droit et la morale \textit{fonctionnent} : on juge Kevin, Mballa, Stéphane. La responsabilité suppose une \textbf{liberté minimale} (seuil de conscience + marge de manœuvre), non totale.
        \item \textbf{Synthèse :} La responsabilité est \textbf{graduée} selon le degré de liberté effective (Kevin > Mballa > Stéphane). Le droit positif institue des seuils (âge, intention/imprudence, causes d'atténuation) pour approximer cette graduation. La responsabilité morale pointe au-delà : la \textbf{responsabilité collective} (État, employeurs) dans la production des situations qui rétrécissent la liberté des vulnérables.
    \end{itemize}
    \item \textbf{Rédiger :} Introduction (accroche, définitions, problème, plan), Développement (arguments + exemples précis des 3 dossiers + auteurs), Conclusion (réponse nuancée, ouverture).
\end{enumerate}
\end{methode}

\begin{exoresolu}[Corrigé type : Analyse comparative des trois dossiers et verdicts argumentés]
\textbf{Énoncé :} Appliquer les notions de liberté et de responsabilité aux trois dossiers du tribunal simulé et rendre un verdict argumenté pour chacun.

\tcblower
\textbf{Solution détaillée :}

\noindent\textbf{1. DOSSIER N\textsuperscript{o} 1 (Kevin - Triche au Bac blanc)}

\medskip
\noindent\textbf{Analyse de la liberté :}
\begin{itemize}
    \item \textbf{Contraintes alléguées :} Pression du système (coefficients, charge), pression des pairs (« tout le monde le fait »), enjeu vital (Bac = avenir ENSET). $\rightarrow$ Contraintes \textbf{psychologiques et sociales}, non physiques. Kevin \textbf{connaît} la règle (interdiction de tricher) et la sanction.
    \item \textbf{Conscience :} Pleine. Il sait que c'est une fraude (dissimulation dans calculatrice). L'argument « tout le monde le fait » est une \textbf{rationalisation} (mécanisme de défense), pas une contrainte irrésistible.
    \item \textbf{Degré de liberté :} \textbf{Élevée}. Kevin est majeur (18 ans), scolarisé, sans trouble psychique avéré. Il a \textit{choisi} la facilité/triche plutôt que l'effort/risque d'échec. La situation est \textit{difficile} mais pas \textit{impossible} (moyenne 11,5 montre capacité de réussite).
\end{itemize}

\medskip
\noindent\textbf{Responsabilité engagée :}
\begin{itemize}
    \item \textbf{Morale :} \textbf{Totale.} Violation du devoir d'honnêteté (Kant : ne pas traiter l'examen comme simple moyen, respecter la règle universelle « ne pas tricher »). Aristote : acte \textit{volontaire} (principe en lui, connaissance des circonstances).
    \item \textbf{Civile :} Faible/Indirecte. Dommage à la collectivité (valeur du diplôme) et à l'établissement (image). Réparation symbolique.
    \item \textbf{Pénale :} \textbf{Fraude aux examens (Art. 139 Code Pénal).} Infraction intentionnelle. Élément matériel (documents interdits) + Élément moral (volonté de les utiliser). Majorité pénale acquise.
\end{itemize}

\medskip
\noindent\textbf{Verdict type (Procureur) :} \textit{« Le tribunal retient la qualification de fraude aux examens. La liberté de Kevin n'était pas abolie par la pression sociale ou scolaire ; elle était pleinement exercée dans le choix de la tricherie. La responsabilité morale est entière. Sur le plan pénal, relaxe au bénéfice de la dispense de peine (premiers faits, enjeu éducatif) mais inscription au casier judiciaire (B2) pour rappel à la loi. Sur le plan disciplinaire : exclusion temporaire 8 jours, note 0 à l'épreuve, avertissement. La sanction vise la réparation symbolique et la réappropriation de l'éthique scolaire. »}

\bigskip
\noindent\textbf{2. DOSSIER N\textsuperscript{o} 2 (Mballa - Accident mortel Ndogbati)}

\medskip
\noindent\textbf{Analyse de la liberté :}
\begin{itemize}
    \item \textbf{Contraintes :} \textbf{Physique} (route dégradée, freins défectueux $\rightarrow$ perte de contrôle partielle). \textbf{Hiérarchique/Économique} (ordres du patron, quota de trajets, menace de licenciement, salaire de misère). $\rightarrow$ Contrainte \textbf{forte mais pas irrésistible} (il conduit le bus, il accepte les trajets).
    \item \textbf{Conscience :} Il connaît l'état des freins (signalé 3 semaines avant), l'état de la route, la limitation de vitesse (45 km/h au lieu de 30). Il \textbf{a conscience du risque} mais \textbf{sous-estime sa gravité} (routine, pression temps).
    \item \textbf{Degré de liberté :} \textbf{Atténuée (liberté conditionnelle).} Mballa est \textit{situé} dans une structure qui réduit sa marge de manœuvre (précaire, subordonné). Cependant, il garde la \textbf{maîtrise de la vitesse} et le \textbf{choix de refuser le départ} (droit de retrait, Art. 123 Code Travail : danger grave et imminent). Il n'a pas exercé ce droit.
\end{itemize}

\medskip
\noindent\textbf{Responsabilité engagée :}
\begin{itemize}
    \item \textbf{Morale :} \textbf{Partagée.} Mballa responsable d'imprudence (vitesse, non-exercice droit de retrait). Employeur (\textit{Trans-Union}) \textbf{principalement responsable} : défaut d'entretien (chose), pression illicite (ordre de rouler avec véhicule dangereux), violation Code Travail. Responsabilité morale de l'employeur \textbf{supérieure} (créateur du risque structurel).
    \item \textbf{Civile :} \textbf{Solidaire.} Art. 1384 al. 5 (responsabilité du fait des choses) + Art. 1384 al. 4 (responsabilité du commettant pour le préposé). \textit{Trans-Union} indemnise victimes (assurance). Action récursoire possible contre Mballa si faute détachable du service (faute intentionnelle ou d'une gravité exceptionnelle -- ici simple imprudence, peu probable).
    \item \textbf{Pénale :} \textbf{Homicide et blessures involontaires (Art. 269-270 Cp).} Élément moral : \textit{non-intentionnel} (imprudence, négligence, inobservation règlements). Peine encourue : emprisonnement 1 à 5 ans + amende. Circonstances atténuantes : pression employeur, état route. Responsabilité pénale de l'employeur (personne morale) possible (Art. 32 Cp : complicité par instruction).
\end{itemize}

\medskip
\noindent\textbf{Verdict type (Parties civiles / Procureur) :} \textit{« Le tribunal déclare Mballa coupable d'homicide et blessures involontaires par imprudence (vitesse excessive, défaut maîtrise véhicule). La contrainte hiérarchique n'abolit pas la liberté de refuser un danger grave (droit de retrait). Responsabilité morale atténuée mais réelle. Responsabilité civile principale à charge de Trans-Union SARL (employeur, gardien véhicule) pour indemnisation intégrale victimes. Peine : 12 mois de prison avec sursis, 500 000 FCFA d'amende, retrait permis 3 ans. Trans-Union condamné solidairement aux dommages-intérêts (estimation 50 millions FCFA) et à 5 millions FCFA d'amende (personne morale). Mesure complémentaire : expertise parc automobile obligatoire sous 3 mois. »}

\bigskip
\noindent\textbf{3. DOSSIER N\textsuperscript{o} 3 (Stéphane - Vol avec violence Makepe)}

\medskip
\noindent\textbf{Analyse de la liberté :}
\begin{itemize}
    \item \textbf{Contraintes :} \textbf{Physique/Psychologique forte :} Menace/pression du frère aîné (autorité, violence potentielle), faim extrême (besoin vital non satisfait), mère malade/absente, pauvreté structurelle. Âge : \textbf{15 ans (mineur)}.
    \item \textbf{Conscience :} \textbf{Altérée mais existante.} Il sait que voler et blesser est interdit (« pour faire peur seulement », « le couteau a glissé »). Il distingue le bien du mal. Mais sa \textbf{capacité de résistance} à la pression fraternelle et à la misère est \textbf{fortement diminuée} par l'âge et la situation de détresse.
    \item \textbf{Degré de liberté :} \textbf{Faible (liberté entravée).} Situation de \textit{contrainte morale irrésistible} relative (Aristote : contrainte + choix du moindre mal ? Non, ici choix imposé par la terreur et la faim). Le Code Pénal (Art. 80) prévoit l'irresponsabilité < 13 ans, responsabilité atténuée 13-18 ans. La contrainte sociale (pauvreté) + contrainte psychologique (frère) + minorité $\rightarrow$ liberté \textbf{très réduite}.
\end{itemize}

\medskip
\noindent\textbf{Responsabilité engagée :}
\begin{itemize}
    \item \textbf{Morale :} \textbf{Forte atténuation, voire quasi-abolie.} Comment blâmer moralement un enfant de 15 ans affamé, manipulé par un frère délinquant, abandonné par l'État (santé mère, protection enfance) ? Responsabilité morale \textbf{reportée sur la société/État/famille défaillante}. Stéphane est \textit{plus victime que coupable}.
    \item \textbf{Civile :} Responsabilité civile du fait personnel (Art. 1382) mais \textbf{insolvable}. Responsabilité des parents (Art. 1384 al. 4 : père/mère responsables dommages causés par enfant mineur habitant avec eux) $\rightarrow$ Mère malade/insolvable. Responsabilité de l'État (défaillance protection enfance, santé) ? Voie difficile.
    \item \textbf{Pénale :} \textbf{Mineur de 15 ans (13-18 ans).} \textbf{Responsabilité pénale atténuée (Art. 80 Cp + Ordonnance 62/OF/18 sur organisation judiciaire).} Pas de peine privative de liberté en principe. \textbf{Mesures de protection, d'assistance, de surveillance et d'éducation (MPASE)} prioritaires (Art. 22 Ordonnance 62). Placemen en centre éducatif ou remise à personne digne de confiance (éducateur). Réparation à la victime (médiation, travail d'intérêt général adapté).
\end{itemize}

\medskip
\noindent\textbf{Verdict type (Juge des Enfants / Procureur) :} \textit{« Le tribunal pour enfants déclare Stéphane auteur des faits de vol avec violence (Art. 318 Cp). Compte tenu de son âge (15 ans), de la contrainte morale exercée par son frère aîné, de l'état de nécessité alimentaire et de l'absence de structure parentale protectrice, sa responsabilité pénale est atténuée au maximum. Aucune peine privative de liberté n'est prononcée. Mesure éducative : placement au Centre d'Éducation Spécialisée (CES) de Bonabéri pour 18 mois renouvelables (scolarisation, formation professionnelle menuiserie, suivi psychologique). Obligation de réparation morale (excuses écrites à Mme Fouda) et participation à un stage de citoyenneté. Signalement au Procureur pour poursuite du frère Bebey (recel, complicité, exploitation mineur). Saisine du Juge des Tutelles pour mesure de protection administrative de la fratrie (aide sociale, prise en charge dialyse mère). »}

\bigskip
\noindent\textbf{4. SYNTHÈSE COMPARATIVE (Méta-réflexion pour le débat final)}

\begin{center}
\begin{tabularx}{\linewidth}{|l|X|X|X|}
\hline
\textbf{Critère} & \textbf{Kevin (D1)} & \textbf{Mballa (D2)} & \textbf{Stéphane (D3)} \\
\hline
\textbf{Liberté effective} & Élevée (choix calculé) & Atténuée (structurelle/hiérarchique) & Faible (âge + contrainte vitale/psychologique) \\
\hline
\textbf{Conscience} & Pleine & Pleine (risque connu) & Présente mais submergée \\
\hline
\textbf{Responsabilité Morale} & Totale & Partagée (Employeur $>$ Employé) & Quasi-nulle (Victime du système) \\
\hline
\textbf{Responsabilité Civile} & Symbolique & Solidaire (Employeur paye) & Théorique (Insolvabilité) \\
\hline
\textbf{Responsabilité Pénale} & Fraude intentionnelle (Majeur) & Homicide involontaire (Imprudence) & Vol violence (Mineur $\rightarrow$ Mesures éducatives) \\
\hline
\textbf{Finalité sanction} & Réparation symbolique / Réinsertion scolaire & Réparation victimes / Dissuasion employeur / Sécurité routière & Protection / Éducation / Réinsertion sociale \\
\hline
\end{tabularx}
\end{center}

\noindent\textbf{Conclusion philosophique :} La responsabilité n'est pas binaire (coupable/innocent) mais \textbf{graduée} selon le degré de liberté effective. Le droit positif (Code Pénal, Code Civil, Ordonnance 62) tente d'approcher cette graduation (majorité/minorité, faute intentionnelle/non-intentionnelle, responsabilité solidaire). Mais la \textbf{responsabilité morale} dépasse le droit : elle pointe la \textbf{responsabilité collective} (État, employeurs, système éducatif) dans la production des situations qui \textit{rétrécissent} la liberté des plus vulnérables (Mballa, Stéphane). Kevin, lui, dispose d'une liberté quasi-pleine : sa responsabilité est le prix de son autonomie.

\end{exoresolu}

\begin{aretenir}[Bilan de la leçon : Liberté et Responsabilité]
Cette activité de tribunal philosophique a permis de :
\begin{itemize}
    \item \textbf{Ancrer les concepts} : La distinction liberté/contrainte (physique, psychologique, sociale) et responsabilité (morale, civile, pénale) n'est plus abstraite mais opératoire pour juger des cas réels camerounais.
    \item \textbf{Montrer la gradation de la responsabilité} : Elle est proportionnelle à l'espace de liberté réel (Kevin > Mballa > Stéphane). La contrainte \textit{irrésistible} (Aristote) est rare ; la contrainte \textit{structurelle} (pauvreté, hiérarchie) est la norme pour les plus fragiles.
    \item \textbf{Révéler le décalage Droit / Morale} : Le droit sanctionne l'acte (Kevin, Mballa) ou protège le mineur (Stéphane). La morale interroge la \textit{justice} de la situation : Mballa paie pour son patron ; Stéphane paie pour la faillite de l'État-providence.
    \item \textbf{Entraîner aux épreuves du Bac} : Dissertation (sujet type : « La responsabilité suppose-t-elle une liberté totale ? » ou « Sommes-nous responsables de ce que nous faisons sous la contrainte ? ») et Commentaire de texte (extrait Aristote, Kant, Sartre, ou Code Pénal/Civil).
    \item \textbf{Développer l'argumentation} : Passer de l'émotion (« c'est pas juste ») à l'argument juridico-philosophique structuré (qualification, imputabilité, finalité de la peine).
\end{itemize}
\emph{Compétence transversale validée :} « Appliquer les notions de l'unité à des situations concrètes de la vie courante » et « Rédiger et justifier une démarche, une réponse ou une conclusion ».
\end{aretenir}

\newpage

\coursec{Méthodes et exercices résolus}

\coursec{Liberté et Responsabilité}

\courssub{Méthode 1 : Appliquer les notions de liberté et de responsabilité à une situation concrète}

\begin{methode}[Analyser une situation concrète avec les concepts de liberté et de responsabilité]
Pour traiter une situation de la vie courante (accident de circulation, vol à l'école, fraude à l'examen, propos injurieux sur les réseaux sociaux...), suivre une démarche rigoureuse en cinq étapes :
\begin{enumerate}
\item \textbf{Décrire les faits} avec précision : qui a fait quoi, dans quelles circonstances, avec quelles conséquences ?
\item \textbf{Identifier le degré de liberté de l'agent} : a-t-il agi volontairement ? Connaissait-il les conséquences possibles de son acte ? A-t-il été contraint (menace, force majeure, ignorance) ?
\item \textbf{Repérer les obstacles éventuels à la liberté} : contrainte extérieure, passion, ignorance, pression sociale. Ces obstacles peuvent atténuer ou supprimer la responsabilité.
\item \textbf{Déterminer le type de responsabilité engagée} : morale (devant la conscience et la communauté), juridique (devant la loi et les tribunaux), civile (réparer le dommage) ou pénale (subir une peine).
\item \textbf{Conclure par un jugement justifié} : dire dans quelle mesure l'agent est responsable, en liant explicitement le degré de liberté constaté au degré de responsabilité retenu.
\end{enumerate}
\textbf{Réflexe à retenir} : pas de responsabilité sans liberté. Toute analyse doit donc commencer par la question : \emph{l'agent pouvait-il faire autrement ?}
\end{methode}

\begin{exoresolu}[Le chauffeur de taxi et l'accident de la rue]
\textbf{Énoncé.} À Douala, un chauffeur de taxi, pressé de finir sa journée, roule à vive allure dans une rue commerçante. Il renverse un élève qui traversait. Un passager affirme : « Ce n'est pas de sa faute, c'est le métier qui l'oblige à se presser. » Analysez cette situation à l'aide des concepts de liberté et de responsabilité, puis dites si l'argument du passager est recevable.
\tcblower
\textbf{Solution détaillée.}
\begin{enumerate}
\item \textbf{Description des faits} : le chauffeur a adopté une vitesse excessive dans une zone fréquentée et a blessé un piéton. Le dommage est réel et directement lié à sa conduite.
\item \textbf{Degré de liberté} : le chauffeur connaît le code de la route et sait qu'une vitesse excessive en ville expose autrui au danger. Rien ne l'y contraignait physiquement : il \emph{pouvait faire autrement} (ralentir, refuser une course supplémentaire). Son acte est donc volontaire et accompli en connaissance de cause.
\item \textbf{Obstacles éventuels} : la pression économique (boucler sa journée) est une contrainte sociale réelle, mais elle ne supprime pas la liberté de choix : elle l'incline, elle ne le force pas. Ce n'est ni la force majeure, ni l'ignorance, ni la menace.
\item \textbf{Type de responsabilité} : responsabilité \textbf{morale} (il a manqué à son devoir de prudence envers autrui), responsabilité \textbf{civile} (il doit réparer le dommage : soins de l'élève) et responsabilité \textbf{pénale} (blessures par imprudence, passibles de sanctions).
\item \textbf{Jugement sur l'argument du passager} : l'argument est non recevable. Il transforme une circonstance atténuante (la pression du métier) en excuse totale. Or, selon le principe que la responsabilité suppose la liberté, seule une contrainte abolissant le choix déresponsabiliserait le chauffeur. Ici, le choix de ralentir existait.
\end{enumerate}
\textbf{Conclusion} : le chauffeur est pleinement responsable, car il a agi librement et en connaissance de cause ; la pression professionnelle peut tout au plus atténuer légèrement l'appréciation de sa faute, sans l'effacer.
\end{exoresolu}

\courssub{Méthode 2 : Définir correctement les concepts en introduction de dissertation}

\begin{methode}[Construire une définition philosophique rigoureuse]
Une bonne définition en introduction de copie obéit à quatre règles :
\begin{enumerate}
\item \textbf{Définir le terme par lui-même}, sans employer le mot défini dans la définition (éviter « la liberté, c'est le fait d'être libre »).
\item \textbf{Dégager l'essence du concept} : la \cle{liberté}, c'est le pouvoir de l'être humain de se déterminer par lui-même, de choisir et d'agir sans contrainte extérieure ni détermination fatale ; la \cle{responsabilité}, c'est l'obligation de répondre de ses actes et d'en assumer les conséquences devant soi-même, autrui et la loi.
\item \textbf{Distinguer les acceptions} : liberté physique (faire ce qu'on veut), liberté morale (obéir à la loi qu'on s'est prescrite), liberté politique (droits du citoyen) ; responsabilité morale, civile, pénale.
\item \textbf{Problématiser immédiatement} : montrer que la définition soulève une difficulté (si tout est déterminé, la liberté existe-t-elle ? peut-on être responsable de ce qu'on n'a pas choisi ?).
\end{enumerate}
\textbf{Formule à retenir} : une définition = l'essence du concept + ses distinctions + la difficulté qu'elle cache.
\end{methode}

\begin{exoresolu}[Rédiger les définitions d'une introduction]
\textbf{Énoncé.} Sujet : « Est-on responsable de tout ce qu'on fait ? » Rédigez le paragraphe de définitions de l'introduction (8 à 12 lignes), en distinguant les acceptions des deux concepts et en dégageant la tension entre eux.
\tcblower
\textbf{Solution détaillée (modèle rédigé).}

« La \textbf{liberté} désigne le pouvoir propre à l'être humain de se déterminer par lui-même : choisir entre plusieurs possibles et agir sans qu'une contrainte extérieure ou une nécessité fatale ne lui impose sa conduite. On distingue la liberté de faire (agir selon ses désirs), la liberté morale (agir selon la raison et le devoir) et la liberté politique (jouir de droits garantis). La \textbf{responsabilité}, quant à elle, est l'obligation pour un agent de répondre de ses actes, c'est-à-dire d'en endosser les conséquences devant sa conscience (responsabilité morale), devant la société et les tribunaux (responsabilité juridique, civile ou pénale). Ces deux notions semblent liées : on ne demande des comptes qu'à celui qui a agi librement. Pourtant, une tension apparaît : nos actes ne sont-ils pas souvent le produit de causes qui nous échappent — l'éducation, le milieu, les passions, la contrainte ? Si tel est le cas, la responsabilité peut-elle s'étendre à tout ce que nous faisons ? »

\textbf{Justification de la démarche} : les définitions sont positives (aucun cercle vicieux), les acceptions sont distinguées, et le paragraphe se clôt sur la difficulté qui annonce la problématique. C'est exactement ce qu'attend le correcteur du Probatoire / Baccalauréat technique.
\end{exoresolu}

\courssub{Méthode 3 : Rédiger et justifier une démarche dialectique (thèse, antithèse, synthèse)}

\begin{methode}[Construire le développement d'une dissertation]
Le développement se construit en trois parties progressives :
\begin{enumerate}
\item \textbf{Thèse} : première réponse au problème, la plus spontanée (ex. : « Oui, on est responsable de tout ce qu'on fait, car on agit librement »). Appuyer sur un auteur (Sartre : l'homme est « condamné à être libre » et responsable de tout ce qu'il fait de lui-même) et un exemple.
\item \textbf{Antithèse} : critique de la thèse (ex. : « Non, car de nombreux actes nous échappent : déterminismes naturels, sociaux, psychologiques »). Appuyer sur un auteur (Spinoza : les hommes se croient libres parce qu'ils ignorent les causes qui les déterminent) et un contre-exemple.
\item \textbf{Synthèse} : dépassement (ex. : « On n'est responsable que dans la mesure où l'on est libre : la responsabilité est proportionnée au degré de liberté et de connaissance de l'agent »). Montrer ce que chaque thèse avait de vrai et de faux.
\end{enumerate}
\textbf{Réflexes} : une idée = un paragraphe ; chaque argument est illustré par un exemple et, si possible, une référence ; chaque partie se termine par une transition critique (« Mais cette position se heurte à... »).
\end{methode}

\begin{exoresolu}[Plan détaillé d'une dissertation]
\textbf{Énoncé.} Sujet : « La liberté est-elle le fondement de la responsabilité ? » Élaborez un plan dialectique détaillé : pour chaque partie, donnez l'idée directrice, deux arguments, un auteur et un exemple.
\tcblower
\textbf{Solution détaillée.}

\textbf{I. Thèse : la responsabilité suppose nécessairement la liberté.}
\begin{itemize}
\item Idée directrice : on ne peut reprocher à quelqu'un que ce qu'il a choisi ; la liberté est donc la condition de la responsabilité.
\item Argument 1 : logiquement, imputer un acte à un agent, c'est affirmer qu'il en est l'auteur, donc qu'il pouvait ne pas le faire.
\item Argument 2 : juridiquement, le droit pénal déresponsabilise le fou, l'enfant en bas âge et celui qui agit sous la contrainte.
\item Référence : Kant — l'impératif moral (« tu dois, donc tu peux ») suppose que la volonté est libre.
\item Exemple : on ne juge pas le somnambule qui casse un objet, car il n'a pas choisi son geste.
\end{itemize}

\textbf{II. Antithèse : mais la liberté est-elle réelle ? Le déterminisme menace la responsabilité.}
\begin{itemize}
\item Idée directrice : si nos actes sont déterminés par des causes (hérédité, milieu, passions), la responsabilité perd son fondement.
\item Argument 1 : Spinoza — l'homme se croit libre parce qu'il a conscience de ses désirs mais ignore les causes qui le déterminent.
\item Argument 2 : les sciences humaines montrent le poids du milieu social et de l'éducation sur les conduites (le délinquant issu d'un quartier défavorisé).
\item Exemple : le juge tient compte du contexte social de l'accusé dans l'application de la peine.
\end{itemize}

\textbf{III. Synthèse : la responsabilité est proportionnée au degré de liberté.}
\begin{itemize}
\item Idée directrice : même si des déterminismes pèsent sur nous, une marge de choix subsiste ; la responsabilité s'exerce dans cette marge.
\item Argument 1 : Sartre — même conditionné, l'homme choisit toujours le sens qu'il donne à sa situation (« ce qui compte, c'est ce qu'on fait de ce qu'on a fait de nous »).
\item Argument 2 : c'est pourquoi la loi module la responsabilité (circonstances atténuantes, minorité, trouble psychique) au lieu de la supprimer.
\item Exemple : deux élèves issus du même quartier de Yaoundé ne réagissent pas de la même manière face à la tentation de la fraude.
\end{itemize}
\textbf{Conclusion du plan} : la liberté, même partielle et conquise, demeure le fondement de la responsabilité ; être responsable, c'est répondre de ce que l'on a fait de sa liberté.
\end{exoresolu}

\courssub{Méthode 4 : Rédiger une introduction complète de dissertation}

\begin{methode}[Les quatre moments de l'introduction]
\begin{enumerate}
\item \textbf{L'accroche (amorce)} : partir d'un constat général lié au sujet (fait de la vie courante, expérience commune), sans citer le sujet mot à mot.
\item \textbf{L'analyse des termes et la définition} : définir les concepts clés du sujet (ici : liberté, responsabilité) et dégager leur sens.
\item \textbf{La problématisation} : montrer la difficulté, le paradoxe ou la tension contenue dans le sujet.
\item \textbf{La question problématique et l'annonce du plan} : formuler la question centrale (souvent en transformant le sujet) et annoncer brièvement les étapes du développement.
\end{enumerate}
\textbf{Interdits} : pas de définition du dictionnaire recopiée, pas de « depuis la nuit des temps », pas d'annonce du plan sous forme de liste scolaire trop sèche.
\end{methode}

\begin{exoresolu}[Rédiger une introduction complète]
\textbf{Énoncé.} Sujet : « Peut-on être responsable sans être libre ? » Rédigez l'introduction complète de la dissertation (15 à 20 lignes).
\tcblower
\textbf{Solution détaillée (modèle rédigé).}

« Chaque jour, les tribunaux de notre pays demandent des comptes à des hommes et des femmes : le voleur est condamné, l'automobiliste imprudent est sanctionné, l'élève fraudeur est puni. Ces pratiques reposent sur une conviction profonde : celui qui a mal agi devait et pouvait agir autrement. Autrement dit, notre vie sociale toute entière suppose un lien entre la liberté et la responsabilité. La liberté est le pouvoir de se déterminer soi-même, de choisir ses actes sans contrainte fatale ; la responsabilité est l'obligation de répondre de ces actes et d'en assumer les conséquences devant sa conscience, autrui et la loi. Le sens commun affirme donc qu'on ne saurait être responsable de ce qu'on n'a pas librement accompli : on ne juge pas un fou comme on juge un homme sain. Pourtant, cette évidence se complique : nos choix ne sont-ils pas eux-mêmes déterminés par notre hérédité, notre éducation, notre milieu ? Et inversement, la société ne nous tient-elle pas parfois pour responsables d'actes que nous n'avons pas vraiment choisis ? Dès lors, une question s'impose : peut-on être responsable sans être libre ? Nous verrons d'abord que la responsabilité semble inséparable de la liberté ; nous montrerons ensuite que les déterminismes qui pèsent sur l'agent fragilisent ce lien ; nous soutiendrons enfin que la responsabilité, loin d'être abolie, se mesure au degré de liberté réel dont dispose l'agent. »

\textbf{Justification} : l'accroche part d'un fait social camerounais concret ; les deux concepts sont définis ; la tension (sens commun contre déterminisme) est dégagée ; la question problématique reprend le sujet ; le plan en trois temps est annoncé en une phrase fluide.
\end{exoresolu}

\courssub{Méthode 5 : Mobiliser correctement les références d'auteurs}

\begin{methode}[Citer un auteur sans se tromper]
\begin{enumerate}
\item \textbf{Choisir la référence pertinente} : pour cette leçon, retenir quatre références sûres : Sartre (l'homme est condamné à être libre, il est responsable de ce qu'il fait de lui-même) ; Kant (la moralité suppose la liberté : « tu dois, donc tu peux ») ; Spinoza (l'illusion de la liberté vient de l'ignorance des causes qui nous déterminent) ; Aristote (on n'est responsable que des actes volontaires, accomplis sans contrainte ni ignorance).
\item \textbf{Formuler la thèse de l'auteur en une phrase claire}, sans prétendre citer un texte mot à mot si l'on n'en est pas sûr.
\item \textbf{Expliquer la référence} : une citation non expliquée ne vaut rien. Après chaque référence, ajouter « autrement dit... » ou « cela signifie que... ».
\item \textbf{Relier la référence à l'argument} : l'auteur doit servir votre démonstration, non la remplacer.
\end{enumerate}
\textbf{Formule à retenir} : référence + explication + lien avec l'argument = un appui théorique valable.
\end{methode}

\begin{exoresolu}[Insérer une référence dans un paragraphe argumentatif]
\textbf{Énoncé.} Rédigez un paragraphe de développement (10 à 15 lignes) soutenant l'idée suivante : « On n'est responsable que des actes accomplis librement et en connaissance de cause », en mobilisant Aristote et un exemple de la vie courante.
\tcblower
\textbf{Solution détaillée (modèle rédigé).}

« La responsabilité ne peut s'attacher qu'à l'acte volontaire. Aristote l'a montré dans sa réflexion sur l'action morale : un acte ne peut être loué ou blâmé que s'il est accompli volontairement, c'est-à-dire sans contrainte extérieure et sans ignorance des circonstances. Autrement dit, celui qui agit sous la menace, ou qui ignore ce qu'il fait, n'est pas véritablement l'auteur de son geste : on ne saurait donc le lui imputer. Cette analyse se vérifie dans la vie quotidienne. Si un élève, bousculé dans la cour de récréation, renverse malgré lui un camarade et le blesse, nul ne songera à le punir : son geste n'était pas choisi, il a subi la contrainte physique. De même, si un commerçant de Mokolo vend sans le savoir un produit avarié que son fournisseur lui a livré, sa responsabilité sera atténuée par son ignorance, tandis que le fournisseur, qui connaissait l'état de la marchandise, répondra pleinement de son acte. La justice elle-même consacre ce principe en distinguant l'acte volontaire de l'accident et en prévoyant la force majeure comme cause d'irresponsabilité. Ainsi, la liberté et la connaissance sont les deux conditions sans lesquelles la responsabilité ne peut être engagée. »

\textbf{Justification} : la référence est exacte (Aristote et l'acte volontaire), elle est expliquée (« autrement dit... »), illustrée par deux exemples camerounais contrastés, et le paragraphe se referme sur l'idée directrice annoncée.
\end{exoresolu}

\courssub{Méthode 6 : Rédiger une conclusion et justifier sa position}

\begin{methode}[Construire une conclusion efficace]
\begin{enumerate}
\item \textbf{Répondre explicitement à la question posée} : reprendre la problématique et y apporter une réponse claire, nuancée mais tranchée.
\item \textbf{Bilan du parcours} : rappeler en une ou deux phrases ce que chaque étape du développement a établi.
\item \textbf{Justifier la position retenue} : montrer pourquoi la synthèse est supérieure aux deux positions opposées (elle conserve leur vérité et écarte leur excès).
\item \textbf{Ouverture} : terminer par une question nouvelle, liée au sujet (ex. : la responsabilité collective, la responsabilité à l'égard des générations futures), sans développer.
\end{enumerate}
\textbf{Interdits} : pas d'idée nouvelle argumentée, pas de simple répétition de l'introduction, pas de « en conclusion » suivi d'un vide.
\end{methode}

\begin{exoresolu}[Rédiger une conclusion complète]
\textbf{Énoncé.} Rédigez la conclusion de la dissertation sur le sujet : « La liberté est-elle le fondement de la responsabilité ? » (8 à 12 lignes), en justifiant la position synthétique et en proposant une ouverture.
\tcblower
\textbf{Solution détaillée (modèle rédigé).}

« Nous pouvons maintenant répondre à notre question : la liberté est bien le fondement de la responsabilité, mais d'une liberté comprise comme pouvoir réel, toujours partiel, de se déterminer soi-même. Notre parcours l'a montré : d'une part, imputer un acte à un agent n'a de sens que si cet agent pouvait faire autrement, ce que confirment la morale kantienne et le droit ; d'autre part, les déterminismes naturels, sociaux et psychologiques rappellent que cette liberté n'est jamais absolue. La synthèse s'impose alors : la responsabilité se mesure au degré de liberté et de connaissance de l'agent — pleine pour l'adulte lucide, atténuée pour celui que la contrainte ou l'ignorance a diminué. Cette position est préférable aux deux extrêmes, car elle évite à la fois l'injustice de rendre l'homme responsable de tout et le laxisme qui le déchargerait de tout. Être responsable, c'est donc répondre de l'usage que l'on a fait de sa liberté. Reste alors une question : si chacun est responsable de ses actes, que penser de la responsabilité collective d'une communauté ou d'une nation tout entière ? »

\textbf{Justification} : la réponse est explicite et nuancée, le bilan reprend les trois étapes, la supériorité de la synthèse est argumentée (elle évite les deux excès), et l'ouverture prolonge la réflexion sans l'alourdir.
\end{exoresolu}

\begin{attention}[Les 5 erreurs qui coûtent des points au Probatoire / Baccalauréat technique]
\begin{enumerate}
\item \textbf{Confondre liberté et libre arbitre avec le « faire n'importe quoi »} : définir la liberté comme « faire tout ce qu'on veut » est une erreur de définition. La liberté philosophique est le pouvoir de se déterminer soi-même, y compris selon la raison et le devoir (Kant). Une définition vague ou circulaire (« la liberté, c'est être libre ») est sanctionnée dès l'introduction.
\item \textbf{Affirmer que la contrainte sociale supprime toute responsabilité} : la pression du milieu, de la famille ou du métier atténue parfois la responsabilité, mais ne l'abolit pas, car une marge de choix subsiste. Ne réserver l'irresponsabilité totale qu'à la force majeure, à la contrainte physique et à l'ignorance invincible.
\item \textbf{Citer un auteur sans l'expliquer ni le relier à l'argument} : écrire « Sartre dit que l'homme est condamné à être libre » puis passer à autre chose ne rapporte aucun point. Toute référence doit être expliquée et intégrée au raisonnement.
\item \textbf{Attribuer une thèse au mauvais auteur} : c'est Spinoza (et non Sartre) qui soutient que l'homme se croit libre parce qu'il ignore les causes qui le déterminent ; c'est Aristote (et non Kant) qui lie la responsabilité à l'acte volontaire accompli sans contrainte ni ignorance. Une attribution fausse discrédite toute la copie.
\item \textbf{Négliger la problématique et la conclusion} : une introduction sans question problématique claire, ou une conclusion qui ne répond pas explicitement à la question posée, fait perdre des points de méthode même si le contenu est bon. Toujours transformer le sujet en question et y répondre en conclusion.
\end{enumerate}
\end{attention}

\newpage

\coursec{Exercices}

\courssub{Je vérifie mes connaissances}

\begin{exercice}[QCM : notions de liberté et de responsabilité]
Pour chaque question, entoure la (ou les) bonne(s) réponse(s).
\begin{enumerate}
\item La liberté physique est :
\begin{enumerate}
\item[a)] la capacité de faire tout ce que l'on veut sans limite ;
\item[b)] la capacité de se mouvoir et d'agir sans contrainte extérieure ;
\item[c)] la capacité de choisir le bien selon sa conscience.
\end{enumerate}
\item Le droit de vote, la liberté d'expression et la liberté de circulation relèvent de la liberté :
\begin{enumerate}
\item[a)] physique ; \quad b) politique ; \quad c) morale.
\end{enumerate}
\item La responsabilité pénale est engagée lorsque :
\begin{enumerate}
\item[a)] on a manqué à un devoir moral ;
\item[b)] on a causé un dommage à autrui sans intention ;
\item[c)] on a commis une infraction prévue et punie par la loi.
\end{enumerate}
\item Selon Sartre, l'homme est :
\begin{enumerate}
\item[a)] déterminé par sa nature ;
\item[b)] condamné à être libre ;
\item[c)] responsable seulement de ses actes volontaires.
\end{enumerate}
\end{enumerate}
\end{exercice}

\begin{exercice}[Vrai ou faux ? Justifie ta réponse]
Réponds par \textbf{vrai} ou \textbf{faux}, puis justifie en deux ou trois lignes.
\begin{enumerate}
\item La liberté consiste à faire tout ce qui nous plaît.
\item On peut être tenu pour responsable d'un acte accompli sous une forte contrainte physique.
\item La responsabilité morale s'adresse d'abord à la conscience de l'individu.
\item Il n'y a de responsabilité que là où il y a eu liberté de choix.
\end{enumerate}
\end{exercice}

\begin{exercice}[Définitions]
Définis chacune des notions suivantes en deux ou trois lignes, en donnant pour chacune un exemple tiré de la vie quotidienne camerounaise :
\begin{enumerate}
\item la liberté ;
\item la responsabilité ;
\item la liberté morale ;
\item la responsabilité civile.
\end{enumerate}
\end{exercice}

\begin{exercice}[Classement des situations]
Recopie et complète le tableau suivant en indiquant, pour chaque situation, le type de liberté ou de responsabilité en jeu (liberté physique, politique ou morale ; responsabilité morale, civile ou pénale).
\begin{center}
\begin{tabularx}{\linewidth}{|X|l|}
\hline
\rowcolor{popblueL} \textbf{Situation} & \textbf{Notion en jeu} \\
\hline
Un électeur choisit librement son candidat aux élections. & \\
\hline
Un commerçant rembourse la marchandise d'un voisin qu'il a détériorée. & \\
\hline
Un élève hésite entre tricher et rendre sa copie honnêtement, puis choisit l'honnêteté. & \\
\hline
Un chauffeur ivre est condamné par le tribunal après un accident. & \\
\hline
\end{tabularx}
\end{center}
\end{exercice}

\courssub{Je m'entraîne}

\begin{exercice}[Analyse d'une situation concrète]
Amina, élève en classe de Terminale technique à Douala, veut s'inscrire en section génie civil. Sa famille préférerait qu'elle devienne couturière, « métier de femme ». Amina maintient son choix et s'inscrit en génie civil.
\begin{enumerate}
\item Quel type de liberté Amina exerce-t-elle ? Justifie.
\item La pression familiale supprime-t-elle totalement sa liberté ? Argumente.
\item De quoi Amina devient-elle responsable en assumant son choix ?
\end{enumerate}
\end{exercice}

\begin{exercice}[Degrés de responsabilité]
Trois élèves cassent la vitre de la salle de classe : le premier l'a fait volontairement par colère ; le deuxième l'a brisée en jouant au ballon malgré l'interdiction ; le troisième, poussé par un camarade, l'a heurtée sans le vouloir.
\begin{enumerate}
\item Classe les trois élèves selon leur degré de liberté dans l'acte commis.
\item Classe-les ensuite selon leur degré de responsabilité.
\item Quelle conclusion tires-tu sur le rapport entre liberté et responsabilité ?
\end{enumerate}
\end{exercice}

\begin{exercice}[Tableau comparatif d'affaires]
Construis un tableau comparatif de trois affaires (réelles ou fictives) de ton choix. Pour chacune, tu indiqueras : les faits, le degré de liberté de l'agent (totale, diminuée, nulle), le type de responsabilité engagée (morale, civile, pénale) et la sanction ou la réparation qui te semble juste. Rédige ensuite un paragraphe de conclusion dégageant le critère qui permet d'attribuer une responsabilité.
\end{exercice}

\begin{exercice}[Réflexion argumentée]
Sujet : \emph{La pression sociale et familiale au Cameroun (mariage, orientation scolaire, traditions) supprime-t-elle la liberté individuelle ?}

Rédige une réponse personnelle et justifiée d'environ vingt lignes. Tu devras : définir la liberté individuelle, présenter au moins un argument en faveur de la thèse (la pression supprime la liberté) et un argument contre, puis prendre position de manière argumentée en t'appuyant sur un exemple concret.
\end{exercice}

\courssub{Je me prépare au Bac}

\begin{exercice}[Dissertation — type Baccalauréat technique]
Sujet : \emph{Peut-on être responsable sans être libre ?}

Travail demandé :
\begin{enumerate}
\item Analyse le sujet : définis les termes « responsable » et « libre », dégage l'enjeu philosophique et formule la problématique.
\item Rédige l'introduction complète de la dissertation (amorce, définitions, problématique, annonce du plan).
\item Construis un plan dialectique détaillé en trois parties :
\begin{enumerate}
\item[a)] thèse : on ne peut être responsable sans être libre (la liberté comme condition de la responsabilité) ;
\item[b)] antithèse : il existe des cas où l'on est tenu pour responsable sans avoir été pleinement libre (contrainte, déterminismes sociaux, responsabilité objective) ;
\item[c)] synthèse : ta position personnelle argumentée.
\end{enumerate}
\item Rédige la conclusion.
\end{enumerate}
Pour chaque partie, tu mobiliseras au moins un argument, un exemple concret et, si possible, une référence à un auteur étudié (Sartre, Kant).
\end{exercice}

\begin{exercice}[Explication de texte guidée — Sartre]
Texte : \emph{« L'homme est condamné à être libre. Condamné, parce qu'il ne s'est pas créé lui-même ; et par ailleurs libre, parce que, une fois jeté dans le monde, il est responsable de tout ce qu'il fait. »} (Sartre, \emph{L'Existentialisme est un humanisme}).

Questions :
\begin{enumerate}
\item Explique le sens de l'expression « condamné à être libre ». En quoi s'agit-il d'un paradoxe apparent ?
\item Pourquoi Sartre dit-il que l'homme « ne s'est pas créé lui-même » ? Cela supprime-t-il sa liberté ?
\item Selon Sartre, quel lien existe-t-il entre la liberté et la responsabilité ? De quoi l'homme est-il responsable ?
\item Discussion : cette conception de la responsabilité totale est-elle acceptable ? Un jeune élevé dans la pauvreté ou la violence est-il « responsable de tout ce qu'il fait » ? Rédige un paragraphe argumenté.
\end{enumerate}
\end{exercice}

\begin{exercice}[Étude de cas rédigée — l'accident du moto-taximan]
Faits : à Yaoundé, un moto-taximan roulant sans casque et en excès de vitesse renverse un piéton qui traversait sur un passage protégé. Le piéton est grièvement blessé. Le moto-taximan affirme qu'il « devait se dépêcher pour nourrir sa famille » et qu'un autre véhicule l'a serré sur la droite.

Travail demandé :
\begin{enumerate}
\item Rappelle les conditions qui rendent un agent responsable de ses actes (conscience, liberté, connaissance de la règle).
\item Examine si ces conditions sont réunies dans ce cas : le moto-taximan était-il libre ? Ses justifications (pauvreté, autre véhicule) diminuent-elles ou suppriment-elles sa liberté ?
\item Distingue les types de responsabilité qu'il encourt : morale, civile, pénale. Pour chacune, précise la sanction ou la réparation possible.
\item Rédige un jugement motivé (quinze à vingt lignes) : tu te placeras en juge qui prononce une décision juste et qui la justifie philosophiquement en montrant le lien entre liberté et responsabilité.
\end{enumerate}
\end{exercice}

\newpage

\coursec{Corrigés des exercices}

\begin{corrige}[Exercice 1 — QCM : notions de liberté et de responsabilité]
\begin{enumerate}
\item \textbf{Bonne réponse : b)}. La liberté physique est la capacité de se mouvoir et d'agir sans contrainte extérieure (enchaînement, emprisonnement, obstacle matériel). La réponse a) est fausse : la liberté n'est jamais absolue, elle rencontre toujours des limites (les lois, la liberté d'autrui). La réponse c) définit la liberté morale, non la liberté physique.
\item \textbf{Bonne réponse : b)}. Le droit de vote, la liberté d'expression et la liberté de circulation sont des libertés garanties par les lois et les institutions de l'État : ce sont des libertés politiques (ou civiques).
\item \textbf{Bonne réponse : c)}. La responsabilité pénale est engagée par la commission d'une infraction (crime, délit, contravention) prévue et punie par la loi pénale. Le manquement à un devoir moral relève de la responsabilité morale ; le dommage causé sans intention relève de la responsabilité civile (obligation de réparer).
\item \textbf{Bonne réponse : b)}. Pour Sartre, « l'existence précède l'essence » : l'homme n'a pas de nature fixée d'avance, il se fait lui-même par ses choix. Il est donc « condamné à être libre » et responsable de tout ce qu'il fait.
\end{enumerate}
\textbf{Points de vigilance} : ne pas confondre liberté physique (absence de contrainte extérieure sur le corps) et liberté morale (choix selon la conscience) ; bien distinguer les trois types de responsabilité par leur source (conscience, dommage, loi pénale).
\end{corrige}

\begin{corrige}[Exercice 2 — Vrai ou faux ?]
\begin{enumerate}
\item \textbf{Faux.} Faire tout ce qui nous plaît, c'est suivre nos désirs sans règle : c'est le caprice ou la licence, non la liberté. La vraie liberté suppose la capacité de choisir en connaissance de cause et rencontre des limites : la loi et la liberté d'autrui. « Ma liberté s'arrête où commence celle des autres. »
\item \textbf{Faux.} La responsabilité suppose la liberté : un acte accompli sous une contrainte physique irrésistible (on est poussé, menacé d'une arme) n'est pas un acte libre. La contrainte totale supprime la responsabilité ; une contrainte partielle peut seulement la diminuer (circonstances atténuantes).
\item \textbf{Vrai.} La responsabilité morale engage l'individu devant sa propre conscience : c'est elle qui juge si l'acte est conforme au devoir. Son « tribunal » est intérieur (remords, sentiment de culpabilité), même si la conscience peut être éclairée par la société.
\item \textbf{Vrai.} C'est le principe central de la leçon : on ne peut être tenu pour responsable que d'un acte que l'on a accompli librement et en connaissance de cause. La liberté de choix est la condition de la responsabilité : sans liberté, pas d'imputabilité.
\end{enumerate}
\end{corrige}

\begin{corrige}[Exercice 3 — Définitions]
\begin{enumerate}
\item \textbf{La liberté} : pouvoir de penser, de choisir et d'agir par soi-même, sans contrainte extérieure ni détermination absolue, dans le respect des lois et de la liberté d'autrui. \emph{Exemple} : un jeune de Bafoussam choisit librement sa filière de formation après le BEPC.
\item \textbf{La responsabilité} : obligation de répondre de ses actes devant soi-même, devant autrui ou devant la société, et d'en assumer les conséquences (réparation, sanction). \emph{Exemple} : un conducteur qui a endommagé un véhicule à un carrefour de Yaoundé doit réparer le dommage.
\item \textbf{La liberté morale} : capacité de choisir entre le bien et le mal selon sa conscience et sa raison, en s'affranchissant des passions et des pressions. \emph{Exemple} : un élève qui refuse de tricher à l'examen alors que ses camarades le font.
\item \textbf{La responsabilité civile} : obligation légale de réparer un dommage causé à autrui, même sans intention de nuire. \emph{Exemple} : un enfant casse la vitrine d'une boutique à Garoua en jouant au ballon ; ses parents doivent indemniser le commerçant.
\end{enumerate}
\end{corrige}

\begin{corrige}[Exercice 4 — Classement des situations]
\begin{center}
\begin{tabularx}{\linewidth}{|X|l|}
\hline
\rowcolor{popblueL} \textbf{Situation} & \textbf{Notion en jeu} \\
\hline
Un électeur choisit librement son candidat aux élections. & Liberté politique \\
\hline
Un commerçant rembourse la marchandise d'un voisin qu'il a détériorée. & Responsabilité civile \\
\hline
Un élève hésite entre tricher et rendre sa copie honnêtement, puis choisit l'honnêteté. & Liberté morale \\
\hline
Un chauffeur ivre est condamné par le tribunal après un accident. & Responsabilité pénale \\
\hline
\end{tabularx}
\end{center}
\textbf{Justifications} : le vote est un droit civique garanti par l'État (liberté politique) ; rembourser un bien détérioré, c'est réparer un dommage (responsabilité civile) ; le choix intérieur entre le bien et le mal engage la conscience (liberté morale) ; la conduite en état d'ivresse est une infraction punie par la loi (responsabilité pénale).
\end{corrige}

\begin{corrige}[Exercice 5 — Analyse d'une situation concrète (Amina)]
\begin{enumerate}
\item Amina exerce d'abord sa \cle{liberté morale} : elle choisit selon sa conscience et sa raison, contre les préjugés (« métier de femme »). Elle exerce aussi une forme de liberté politique ou civile, car le droit à l'éducation et à l'orientation est garanti à tous les citoyens, filles et garçons.
\item La pression familiale ne supprime pas totalement sa liberté : elle la \emph{rend difficile}, mais Amina conserve la possibilité de choisir, comme le prouve son inscription effective. La pression sociale est un obstacle, non une contrainte physique irrésistible. Tant que le choix reste possible, la liberté demeure, même diminuée.
\item En assumant son choix, Amina devient responsable de son avenir : responsable devant sa conscience (fidélité à elle-même), responsable de sa réussite scolaire en génie civil, et responsable des conséquences de sa décision (efforts à fournir, dialogue avec sa famille). Sa liberté de choisir fonde sa responsabilité d'assumer.
\end{enumerate}
\end{corrige}

\begin{corrige}[Exercice 6 — Degrés de responsabilité]
\begin{enumerate}
\item \textbf{Degré de liberté} (du plus libre au moins libre) : le premier élève (acte volontaire, pleine liberté) ; le deuxième (acte imprudent mais choisi : il jouait malgré l'interdiction, liberté réelle mais mal exercée) ; le troisième (poussé, acte involontaire, liberté nulle au moment du geste).
\item \textbf{Degré de responsabilité} : le premier est le plus responsable (intention de nuire : responsabilité morale et civile, sanction disciplinaire possible) ; le deuxième est responsable d'imprudence (responsabilité civile : réparation de la vitre) ; le troisième n'est pas responsable, ou très peu : c'est le camarade qui l'a poussé qui porte la responsabilité.
\item \textbf{Conclusion} : le degré de responsabilité est proportionnel au degré de liberté. Plus l'acte est libre et volontaire, plus l'agent est responsable ; la contrainte (être poussé) supprime ou diminue la responsabilité. Cela confirme que la liberté est le fondement de la responsabilité.
\end{enumerate}
\end{corrige}

\begin{corrige}[Exercice 7 — Tableau comparatif d'affaires]
Exemple de tableau attendu (les affaires peuvent varier) :
\begin{center}
\begin{tabularx}{\linewidth}{|X|l|l|X|}
\hline
\rowcolor{popblueL} \textbf{Faits} & \textbf{Degré de liberté} & \textbf{Responsabilité} & \textbf{Sanction / réparation juste} \\
\hline
Un élève vole délibérément le téléphone d'un camarade. & Totale & Morale et pénale & Restitution, sanction disciplinaire, poursuites possibles \\
\hline
Un cultivateur brûle son champ et le feu gagne la plantation voisine. & Diminuée (imprudence) & Civile & Indemnisation du voisin \\
\hline
Un homme bousculé dans la foule renverse une marchande. & Nulle & Aucune (ou responsabilité du bousculeur) & Pas de sanction ; excuses \\
\hline
\end{tabularx}
\end{center}
\textbf{Paragraphe de conclusion (modèle)} : ces trois affaires montrent que le critère d'attribution de la responsabilité est le degré de liberté de l'agent au moment de l'acte. Celui qui a agi volontairement et en connaissance de cause est pleinement responsable ; celui qui a causé un dommage par imprudence doit réparer, car il pouvait prévoir ; celui qui a été contraint n'est pas responsable. La justice consiste donc à mesurer la sanction à la liberté réelle de l'agent : c'est pourquoi les tribunaux recherchent l'intention et admettent les circonstances atténuantes.
\end{corrige}

\begin{corrige}[Exercice 8 — Réflexion argumentée (pression sociale et liberté)]
\textbf{Corrigé type (plan et arguments attendus)} :
\begin{itemize}
\item \textbf{Définition} : la liberté individuelle est le pouvoir de choisir sa vie (orientation, mariage, convictions) selon sa propre volonté éclairée.
\item \textbf{Argument pour la thèse} : au Cameroun, les traditions, le mariage arrangé dans certaines communautés, l'orientation imposée par les parents pèsent lourdement sur les jeunes ; la crainte du rejet familial peut rendre le choix impossible en pratique : la pression supprime alors la liberté de fait.
\item \textbf{Argument contre} : la pression sociale n'est pas une contrainte physique irrésistible ; elle influence sans déterminer totalement. De nombreux jeunes négocient, dialoguent, ou affirment leur choix (exemple d'Amina). La liberté intérieure (morale) demeure toujours : on peut toujours dire non, même si cela a un coût.
\item \textbf{Position personnelle argumentée} : la pression sociale \emph{diminue} la liberté sans la supprimer. Elle rend le choix coûteux, mais le choix reste possible. La solution n'est pas la rupture avec la famille ou la tradition, mais le dialogue et l'éducation, qui élargissent l'espace de liberté. \emph{Exemple} : une jeune fille de Maroua qui obtient, par la discussion avec ses parents, de poursuivre ses études avant le mariage concilie respect de la famille et liberté personnelle.
\end{itemize}
\textbf{Points de vigilance} : éviter les généralisations abusives sur « les traditions » ; distinguer influence (qui laisse le choix) et détermination (qui le supprime) ; conclure par une position nuancée et justifiée.
\end{corrige}

\begin{corrige}[Exercice 9 — Dissertation : « Peut-on être responsable sans être libre ? »]
\textbf{1. Analyse du sujet.} « Responsable » : qui doit répondre de ses actes et en assumer les conséquences. « Libre » : qui agit par choix propre, sans contrainte. Enjeu : la justice et la morale attribuent la responsabilité ; mais si l'homme n'est pas libre (déterminismes, contraintes), peut-on légitimement le punir ou le louer ? \textbf{Problématique} : la responsabilité suppose-t-elle nécessairement la liberté, ou peut-on être tenu pour responsable d'actes qui ne dépendaient pas entièrement de nous ?

\textbf{2. Introduction (modèle rédigé).} Chaque jour, les tribunaux jugent, les parents grondent, la conscience accuse ou absout : partout, nous tenons les hommes pour responsables de leurs actes. Être responsable, c'est devoir répondre de ce que l'on a fait ; être libre, c'est avoir pu choisir d'agir autrement. Mais tout acte est-il choisi ? La misère, la peur, l'éducation ne déterminent-elles pas nos conduites ? Dès lors, peut-on être responsable sans être libre ? Nous verrons d'abord que la liberté est la condition de la responsabilité, puis qu'il existe des cas de responsabilité sans pleine liberté, avant de montrer que la responsabilité, pour être juste, doit être proportionnée à la liberté réelle de l'agent.

\textbf{3. Plan dialectique détaillé.}
\begin{itemize}
\item \textbf{Thèse} : pas de responsabilité sans liberté. Argument : on ne peut reprocher à quelqu'un que ce qu'il a choisi ; un acte sous contrainte totale (être poussé) n'est pas imputable. Exemple : l'élève poussé contre la vitre. Référence : Kant — la morale suppose la liberté, car le devoir (« tu dois ») implique le pouvoir (« tu peux »).
\item \textbf{Antithèse} : on est parfois responsable sans avoir été pleinement libre. Argument : la loi punit l'imprudence (responsabilité objective) ; les déterminismes sociaux (pauvreté, milieu) diminuent la liberté sans effacer la responsabilité. Exemple : le chauffeur ivre, même fatigué par la misère, est condamné. Référence : Sartre — « condamné à être libre », l'homme est responsable de tout ce qu'il fait, même de ce qu'il n'a pas choisi d'être.
\item \textbf{Synthèse} : la responsabilité pleine exige la liberté, mais la liberté humaine est toujours située et partielle. La justice doit donc proportionner la responsabilité au degré de liberté : d'où les circonstances atténuantes, l'irresponsabilité pénale du fou ou du contraint. Être responsable, c'est répondre de ce qui dépendait de nous.
\end{itemize}

\textbf{4. Conclusion (modèle).} On ne peut être pleinement responsable sans être libre : la liberté est le fondement de la responsabilité. Mais comme la liberté réelle admet des degrés, la responsabilité juste admet aussi des degrés. Reconnaître ce lien, c'est à la fois respecter la dignité de l'homme, capable de choisir, et faire preuve d'équité envers celui dont la liberté était diminuée.

\textbf{Barème indicatif (sur 20)} : analyse et problématique (4 pts) ; introduction rédigée (4 pts) ; thèse et antithèse argumentées avec exemples et références (8 pts) ; synthèse personnelle et conclusion (4 pts). \textbf{Points de vigilance} : définir les termes dès l'introduction ; ne pas opposer liberté absolue et déterminisme absolu sans nuance ; citer correctement Kant (impératif moral) et Sartre (existentialisme).
\end{corrige}

\begin{corrige}[Exercice 10 — Explication de texte guidée (Sartre)]
\begin{enumerate}
\item « Condamné à être libre » signifie que l'homme n'a pas choisi d'exister ni d'être libre : la liberté lui est imposée, comme une condamnation. Le paradoxe apparent tient à ce que « condamné » évoque une peine subie, tandis que « libre » évoque un pouvoir : Sartre réunit les deux pour dire que la liberté est notre condition inévitable, non un cadeau choisi.
\item L'homme « ne s'est pas créé lui-même » : il naît dans un corps, une famille, un pays, une époque qu'il n'a pas choisis (sa « facticité »). Mais cela ne supprime pas sa liberté : dès qu'il est « jeté dans le monde », il choisit ce qu'il fait de ce qu'on a fait de lui. La situation est donnée, le sens qu'on lui donne est libre.
\item Pour Sartre, liberté et responsabilité sont inséparables : être libre, c'est être responsable de \emph{tout} ce qu'on fait, sans pouvoir s'excuser sur une nature, Dieu ou la société. L'homme est responsable de ses actes, de ses choix, et même de ce qu'il devient, car il se définit par ses actes.
\item \textbf{Discussion (modèle de paragraphe)} : la responsabilité totale de Sartre est exigeante mais contestable. Dire qu'un jeune élevé dans la pauvreté et la violence est « responsable de tout ce qu'il fait » ignore le poids réel des conditions sociales : la misère réduit les choix possibles, et la violence subie marque les conduites. Cependant, la thèse de Sartre a une force : elle refuse de traiter l'homme comme une simple chose déterminée et lui reconnaît toujours une marge de choix, donc une dignité. La position juste semble intermédiaire : la liberté n'est jamais nulle, mais elle est inégale selon les conditions ; la responsabilité doit donc être réelle mais proportionnée, comme le reconnaît le droit avec les circonstances atténuantes.
\end{enumerate}
\textbf{Points de vigilance} : ne pas faire de Sartre un partisan de la toute-puissance individuelle (il reconnaît la situation) ; distinguer facticité (ce qui est donné) et transcendance (ce que l'on en fait).
\end{corrige}

\begin{corrige}[Exercice 11 — Étude de cas rédigée : l'accident du moto-taximan]
\begin{enumerate}
\item Un agent est responsable de ses actes lorsque trois conditions sont réunies : la \cle{conscience} (il sait ce qu'il fait), la \cle{liberté} (il pouvait agir autrement, sans contrainte irrésistible) et la \cle{connaissance de la règle} (il connaît la loi ou le devoir qu'il viole).
\item Ces conditions sont réunies. Le moto-taximan était conscient et connaît le code de la route (casque obligatoire, limitation de vitesse, priorité au passage protégé). Il était libre : rouler vite était un choix, non une nécessité absolue. La pauvreté (« nourrir sa famille ») explique sa hâte mais ne le contraignait pas physiquement : c'est une circonstance qui \emph{diminue} légèrement la culpabilité sans supprimer la liberté. Le véhicule qui l'a « serré » peut partager la responsabilité s'il est identifié, mais ne rend pas l'excès de vitesse involontaire.
\item \textbf{Responsabilité morale} : il doit se reprocher son imprudence devant sa conscience (remords, faute contre le devoir de prudence). \textbf{Responsabilité civile} : il doit réparer le dommage — payer les soins du piéton, indemniser le préjudice (assurance ou paiement direct). \textbf{Responsabilité pénale} : blessures involontaires aggravées (vitesse, absence de casque) : amende, suspension du permis, voire peine de prison selon la gravité.
\item \textbf{Jugement motivé (modèle)} : « Le tribunal constate que le prévenu, conscient et libre de ses actes, a violé délibérément les règles de circulation : vitesse excessive, absence de casque, non-respect d'un passage protégé. Ses justifications ne tiennent pas : la pauvreté, si elle mérite la compassion sociale, n'autorise pas à mettre la vie d'autrui en danger, car nul n'est contraint par la misère d'enfreindre la loi ; quant au véhicule gênant, il ne supprime pas le choix de ralentir. La liberté du prévenu étant établie, sa responsabilité est engagée : il est condamné à indemniser intégralement la victime (responsabilité civile), à une amende et à une suspension de permis (responsabilité pénale), la peine étant atténuée en raison de sa situation sociale. Ce jugement illustre le principe philosophique selon lequel la responsabilité se fonde sur la liberté : parce que le prévenu pouvait choisir de respecter la règle, il doit répondre de ne pas l'avoir fait ; et parce que sa liberté, bien que réelle, s'exerçait dans des conditions difficiles, la sanction demeure mesurée. Ainsi la justice punit l'acte libre tout en tenant compte des conditions de l'agent. »
\end{enumerate}
\textbf{Barème indicatif} : conditions rappelées (3 pts) ; examen critique des justifications (5 pts) ; distinction des trois responsabilités avec sanctions (6 pts) ; jugement rédigé, cohérent et philosophiquement justifié (6 pts). \textbf{Points de vigilance} : ne pas confondre excuse (qui supprime la responsabilité) et circonstance atténuante (qui la diminue) ; citer la sanction adaptée à chaque type de responsabilité.
\end{corrige}

\newpage

\coursec{Fiche bilan}

\begin{popfigure}
\begin{tikzpicture}[mindmap, grow cyclic, every node/.style={concept, font=\small},
  root concept/.append style={concept color=popblue, font=\bfseries},
  level 1/.append style={level distance=3.4cm, sibling angle=51.4},
  level 2/.append style={level distance=2.1cm, sibling angle=45, font=\scriptsize}]
\node[root concept]{Liberté et Responsabilité}
  child[concept color=poporange]{ node {Définitions}
    child { node {Liberté : pouvoir d'agir par soi-même} }
    child { node {Responsabilité : répondre de ses actes} }
  }
  child[concept color=popgreen]{ node {Formes de la liberté}
    child { node {Liberté physique} }
    child { node {Liberté morale et politique} }
  }
  child[concept color=poppurple]{ node {Formes de la responsabilité}
    child { node {Morale} }
    child { node {Juridique (civile, pénale)} }
    child { node {Politique et sociale} }
  }
  child[concept color=popgold]{ node {Liberté, fondement de la responsabilité}
    child { node {Pas de responsabilité sans liberté} }
    child { node {Kant, Sartre} }
  }
  child[concept color=poppink]{ node {Auteurs et thèses}
    child { node {Spinoza : liberté consciente} }
    child { node {Sartre : « condamné à être libre »} }
    child { node {Kant : autonomie} }
  }
  child[concept color=popblue!60]{ node {Applications}
    child { node {Vie quotidienne camerounaise} }
    child { node {École, famille, cité} }
  }
  child[concept color=popgreen!70!popdark]{ node {Méthode}
    child { node {Dissertation : problématique, plan} }
    child { node {Argumenter et conclure} }
  };
\end{tikzpicture}
\legende{Carte mentale de la leçon : les sept branches essentielles à mobiliser le jour de l'examen.}
\end{popfigure}

\begin{aretenir}[L'essentiel de la leçon]
\textbf{Définitions clés (à savoir par cœur)}
\begin{itemize}
  \item \cle{Liberté} : pouvoir de l'homme d'agir par lui-même, de se déterminer sans contrainte extérieure, selon sa volonté et sa raison.
  \item \cle{Liberté physique} : absence d'entrave matérielle (aller et venir librement).
  \item \cle{Liberté morale} : capacité de choisir entre le bien et le mal en conscience.
  \item \cle{Liberté politique} : droits reconnus au citoyen (expression, vote, association).
  \item \cle{Responsabilité} : obligation de répondre de ses actes et d'en assumer les conséquences devant soi-même, devant autrui et devant la loi.
  \item \cle{Responsabilité juridique} : obligation légale de réparer le dommage causé (civile) ou de subir une peine (pénale).
\end{itemize}

\textbf{Thèses des auteurs}
\begin{center}
\begin{tabularx}{\linewidth}{|l|X|}
\hline
\rowcolor{popblueL} \textbf{Auteur} & \textbf{Thèse à retenir} \\
\hline
Spinoza & Les hommes se croient libres parce qu'ils ignorent les causes qui les déterminent : la vraie liberté est connaissance de la nécessité. \\
\hline
Kant & La liberté est \emph{autonomie} : obéir à la loi morale que la raison se donne à elle-même. \\
\hline
Sartre & L'homme est « condamné à être libre » : il est responsable de tout ce qu'il fait, sans excuse. \\
\hline
\end{tabularx}
\end{center}

\textbf{Idée-force de la leçon}
\begin{itemize}
  \item La \cle{liberté} est le \cle{fondement} de la responsabilité : on ne peut être tenu responsable que d'un acte accompli librement.
  \item La contrainte, la folie, la force majeure ou l'ignorance invincible atténuent ou suppriment la responsabilité.
  \item Réciproquement, être libre, c'est accepter d'assumer : liberté et responsabilité sont inséparables.
\end{itemize}

\textbf{Méthodes express}
\begin{itemize}
  \item \emph{Dissertation} : définir les deux termes du sujet dès l'introduction ; dégager la problématique (« la liberté suffit-elle à fonder la responsabilité ? ») ; plan dialectique (thèse, antithèse, synthèse).
  \item \emph{Application à une situation concrète} : 1) identifier l'acte ; 2) vérifier la liberté de l'agent (contrainte ? discernement ?) ; 3) conclure sur le degré de responsabilité.
  \item \emph{Conclusion} : répondre clairement à la problématique et ouvrir sur la dimension sociale ou citoyenne.
\end{itemize}
\end{aretenir}

\begin{aretenir}[Checklist avant le Bac]
\begin{itemize}
  \item[$\square$] Je sais définir la liberté et la responsabilité en une phrase chacune.
  \item[$\square$] Je distingue liberté physique, morale et politique avec un exemple pour chacune.
  \item[$\square$] Je distingue responsabilité morale, juridique (civile et pénale) et politique.
  \item[$\square$] Je sais expliquer pourquoi la liberté fonde la responsabilité.
  \item[$\square$] Je connais les causes d'atténuation ou d'irresponsabilité (contrainte, folie, minorité, force majeure).
  \item[$\square$] Je peux citer et expliquer la thèse de Spinoza sur l'illusion de la liberté.
  \item[$\square$] Je peux citer et expliquer l'autonomie chez Kant.
  \item[$\square$] Je peux citer et expliquer la formule de Sartre : « condamné à être libre ».
  \item[$\square$] Je sais analyser une situation concrète camerounaise (école, famille, circulation, réseaux sociaux) avec les deux notions.
  \item[$\square$] Je sais construire une problématique et un plan dialectique en trois parties.
  \item[$\square$] Je rédige une conclusion qui répond à la problématique posée.
  \item[$\square$] Je soigne l'orthographe des noms d'auteurs et des concepts clés.
\end{itemize}
\end{aretenir}

\findecours

\end{document}
